\chapter{医学影像}
\label{ch:19}

\paragraph*{本章阅读锚点}
本章的锚点是医学成像的物理约束。医学 3DGS 不能只追求视觉相似，还要处理体数据、投影成像、剂量、解剖结构和诊断可信度。


医学影像领域的三维重建面临独特挑战：稀疏且受控的视角、严格的实时性要求、以及对重建精度的高度敏感。从内窥镜导航到术前规划、从 CT/MRI 体积重建到显微成像，不同医学子领域对 3D 重建的需求差异巨大。3DGS 的实时渲染能力、紧凑的显式表示和可微渲染管线，使其在多个医学应用场景中展现出独特优势。

本章系统梳理了~91~篇相关论文，按~4~个研究方向展开讨论：内窥镜与手术场景重建（第~\ref{sec:19_1}~节）、CT/MRI~断层成像重建（第~\ref{sec:19_2}~节）、显微成像与生物结构重建（第~\ref{sec:19_3}~节）以及临床应用与专用场景（第~\ref{sec:19_4}~节）。

\begin{figure}[htbp]
\centering
\resizebox{\textwidth}{!}{\begin{tikzpicture}[
  route/.style={rectangle, rounded corners=5pt, draw=#1, fill=#1!10,
    text width=3.2cm, minimum height=0.9cm, align=center, font=\small},
  method/.style={rectangle, rounded corners=3pt, draw=#1, fill=#1!6,
    font=\scriptsize, text width=3.0cm, align=center, minimum height=0.6cm},
  arr/.style={-{Stealth[length=4pt]}, techgray!60},
  >=Stealth
]

  \node[route=techgray, text width=4cm, font=\small\bfseries] (root) at (19.7, 5.6) {医学影像方法分类};
  \node[route=techblue] (b0) at (7.9, 3.1) {内窥镜与手术};
  \node[method=techblue] (b0_l0) at (2.0, 0.6) {动态稀疏重建\\{\tiny ST-Endo4DGS, HFGS, EndoSparse, ColonSplat}};
  \node[method=techblue] (b0_l1) at (6.1, 0.6) {专用场景\\{\tiny SurGSplat, EndoLRMGS, Deform3DGS, SurgicalGaussian}};
  \node[method=techblue] (b0_l2) at (9.7, 0.6) {器械建模\\{\tiny Instrument-Splatting, G-SHARP}};
  \node[route=techgreen] (b1) at (17.4, 3.1) {CT/MRI断层成像};
  \node[method=techgreen] (b1_l0) at (13.5, 0.6) {X射线层析\\{\tiny R²-Gaussian, X²-Gaussian, FaCT-GS, DDGS-CT}};
  \node[method=techgreen] (b1_l1) at (17.5, 0.6) {MRI多模态\\{\tiny 3DGSMR, M-Gaussian, PINGS-X, MoRe-3DGSMR}};
  \node[route=techorange] (b2) at (21.9, 3.1) {显微成像};
  \node[method=techorange] (b2_l0) at (21.9, 0.6) {显微纳米结构\\{\tiny CryoSplat, GEM, TG-Field}};
  \node[method=techorange] (b2_l1) at (26.0, 0.6) {生物结构\\{\tiny Denza-Gaussian, 3DGAT}};
  \node[route=techpurple] (b3) at (31.7, 3.1) {临床应用};
  \node[method=techpurple] (b3_l0) at (29.8, 0.6) {可视化专科\\{\tiny DentalGS, AnatVis-GS}};
  \node[method=techpurple] (b3_l1) at (33.7, 0.6) {成像配准\\{\tiny Registering Neural 4D Gaussians}};
  \node[method=techpurple] (b3_l2) at (37.5, 0.6) {多模态融合\\{\tiny InsideOut, FoldGS, VolumeGS}};
  \draw[arr] (root) -- (b0);
  \draw[arr] (root) -- (b1);
  \draw[arr] (root) -- (b2);
  \draw[arr] (root) -- (b3);
  \draw[arr] (b0) -- (b0_l0);
  \draw[arr] (b0) -- (b0_l1);
  \draw[arr] (b0) -- (b0_l2);
  \draw[arr] (b1) -- (b1_l0);
  \draw[arr] (b1) -- (b1_l1);
  \draw[arr] (b2) -- (b2_l0);
  \draw[arr] (b2) -- (b2_l1);
  \draw[arr] (b3) -- (b3_l0);
  \draw[arr] (b3) -- (b3_l1);
  \draw[arr] (b3) -- (b3_l2);
\end{tikzpicture}}
\caption{医学影像方法分类}
\label{fig:ch19_medical_taxonomy}
\end{figure}

%=====================================================================
\section{内窥镜与手术场景重建}
\label{sec:19_1}

内窥镜与手术场景重建是~3DGS~在医学影像中最活跃的应用方向之一。内窥镜图像具有视野狭窄、光照不均匀（光源与相机共轴）、组织表面高度反光（镜面反射）、以及组织持续变形等特殊性。手术场景同时包含刚性器械和高度变形的软组织，需要实时、高保真、可交互的~3D~表示来支持手术导航和训练。这些因素导致标准~3DGS~管线在腔内场景中面临初始化困难、光照建模不准和动态重建质量下降等问题。

\subsection{动态与稀疏视角重建}
\label{ss:19_1}

本节讨论两条技术路线：（1）动态内窥镜重建与4D高斯表示；（2）稀疏视角与SLAM系统。

\subsubsection{动态内窥镜重建与4D高斯表示}

本路线介绍利用 4D 高斯表示建模内窥镜动态场景中组织变形与光照变化的方法。

4D Gaussian Splatting（4DGS）已成为内窥镜动态场景重建的主流范式。其核心思想是将~3DGS~的静态空间表示扩展到~4D~时空体积，使每个高斯基元同时编码空间位置和时间演化——这对于持续变形的组织表面至关重要。具体而言，4DGS~通过时间条件化将~4D~高斯投影为时间调制的~3D~高斯（$P_{3D}(t) = P_{4D} \cdot M(t)$），再经标准~3D~到~2D~投影完成渲染。

\textbf{ST-Endo4DGS}（2024.11）~\cite{2411_01218} 修改了~4DGS~的时空表示与外观编码环节——用非偏置~4D~高斯基元和柱面谐函数替代标准~4DGS~的分离式时空建模，并附加法向对齐约束——光栅化渲染管线保持不变。

\textbf{Baseline~做法与局限。}在~ST-Endo4DGS~之前，内窥镜~4DGS~方法（如~Endo-4DGS）主要采用变形场网络（deformation field）建模组织动态：通过多分辨率~HexPlane~编码时空特征，再将~3DGS~静态高斯沿变形场 warp~到当前时刻。这种分离式设计将外观建模（球谐函数）和运动建模（变形场）解耦为独立模块，在内窥镜场景中面临两个核心问题：一是光源与相机共轴导致强视角依赖的高光变化，标准~SH~无法建模时间维度上的光照演化；二是变形场以离散方式处理时间维度，难以捕获组织变形的连续时间特性。

\textbf{创新一：非偏置~4D~高斯基元。}ST-Endo4DGS~放弃变形场路线，直接在~4D~时空体积中定义高斯基元。每个基元参数化为各向异性~4D~椭圆，支持完整的~4D~旋转，不预先假设空间结构与时间动态可分离（因此称为``非偏置''）。渲染时，通过时间条件化将~4D~高斯转化为时间调制的~3D~高斯，再经~3D→2D~投影完成光栅化：

\begin{equation}
P_{3D}(t) = P_{4D} \times M(t), \quad P_{2D}(t) = W \cdot P_{3D}(t),
\end{equation}

其中~$M(t)$~编码时间依赖关系，$W$~为相机投影矩阵。这一变换的直觉是：在每一帧时刻，从~4D~时空``切片''出对应的~3D~高斯场景，然后按标准~3DGS~流程渲染。此外，通过时间边缘化~$P_{\text{time}} = R \cdot P_{4D}$~将~4D~表示压缩为~3D~分布，$R$~为时间积分变换矩阵，避免了在~4D~空间中直接操作带来的计算和显存膨胀。

\textbf{创新二：柱面谐函数（Spherindrical Harmonics）。}标准~3DGS~使用球谐函数~$Y_l^m(\theta,\phi)$~编码视角相关颜色，但该表示没有时间维度，无法捕获内窥镜场景中随时间变化的光照和外观。ST-Endo4DGS~将~SH~扩展为柱面谐函数——一种混合时空基函数，将空间角频率和时间频率统一在单一正交基中：

\begin{equation}
Z_{nl}^m(t, \theta, \phi) = Y_l^m(\theta, \phi) \cdot \sin(\omega_n t + \phi_n), \quad \omega_n = \frac{2\pi n}{T},
\end{equation}

其中~$Y_l^m(\theta,\phi)$~捕获视角相关外观（与标准~SH~一致），$\sin(\omega_n t + \phi_n)$~以角频率~$\omega_n$~和相位偏移~$\phi_n$~建模外观的时间周期演化，$T$~为时间周期。直觉上，柱面谐函数将``从哪个方向看''和``在哪个时刻看''编码在同一个紧凑系数矩阵中，使单个表示同时处理视角变化和时间变化，避免了分别维护外观网络和变形网络带来的参数冗余。该构造在时空域上构成正交基，保证了时间一致性。

\textbf{创新三：内窥镜法向对齐约束（ENAC）。}单目内窥镜场景缺少真值深度，依赖单目深度估计提供几何监督，但估计精度有限。ST-Endo4DGS~利用~3DGS~高斯基元天然携带几何法向（由旋转矩阵~$R$~和尺度矩阵~$S$~定义）的特性，引入~ENAC~对齐两种法向来源：

\begin{equation}
L_{\text{ENAC}} = \left| N(\text{depth}) - N(\text{Gaussian}) \right|_1,
\end{equation}

其中~$N(\text{depth})$~由渲染深度图求导得到，$N(\text{Gaussian})$~由高斯基元的几何参数直接计算。该约束的直觉是：如果渲染深度和高斯基元几何都正确，二者导出的表面法向应该一致；如果不一致，说明至少一方有误差，约束将其拉回对齐。ENAC~几乎不增加额外计算成本，因为它复用了~3DGS~渲染过程中已有的深度信息。总优化目标为：

\begin{equation}
L_{\text{total}} = L_{\text{render}} + \lambda_{\text{ENAC}} \cdot L_{\text{ENAC}},
\end{equation}

其中~$L_{\text{render}}$~为主渲染损失（$L_1$~光度损失加~D-SSIM~结构损失），$\lambda_{\text{ENAC}}$~控制法向对齐的权重。

\textbf{效果。}在~EndoNeRF~数据集（Cutting/Pulling~两个前列腺切除场景）上，ST-Endo4DGS~以~$1.6 \times 10^{-3}$~学习率训练，达到~PSNR~39.29/38.28、SSIM~0.973/0.966、LPIPS~0.016/0.024，相比~Endo-4DGS~的~PSNR~提升~6.5\%，LPIPS~改善~41.5\%，同时以~123 FPS~运行（比~Endo-4DGS~的~100 FPS~快~23\%），在动态内窥镜场景重建中建立了新的实时渲染基准。

HFGS（2024.05）~\cite{2405_17872} 从频率角度出发，认为内窥镜动态场景的信息集中分布在特定频段。为此，HFGS 分别引入空间高频强调重建（SHF）和时间高频强调重建（THF），利用光流先验聚焦于运动密集区域，在频域层面选择性增强关键信息。在~Baseline~做法中，现有~4DGS~内窥镜方法（如~EndoGS、EndoGaussian）采用变形场网络建模组织动态，通过多分辨率~HexPlane~编码时空特征。然而在高斯致密化过程中，标准逐像素~$L_1$~损失的平均梯度在重建不充分区域依然很小。这导致高斯基元无法触发分裂或克隆，产生持续性的欠重建，表现为渲染图像中的模糊和伪影。HFGS~的核心改进是从空间和时间两个频率维度同时解决欠重建问题：空间维度通过~SHF~模块在频域增强静态细节重建；时间维度通过~THF~模块利用光流先验增强动态区域的重建保真度。

在实现上，SHF~模块对~ground truth~图像经~FFT（快速傅里叶变换）~获取幅度谱，引入二值掩码提取高频幅度，再经逆~FFT~回到空间域得到高频图像，最终以高频图像为权重对逐像素误差加权，使优化重点聚焦于空间高频（边缘、纹理）区域。THF~模块利用预训练光流预测器（基于~UniMatch）分别对渲染图像和~ground truth~图像预测相邻帧间光流，以~Charbonnier~损失和~Census~损失之和约束两者一致，迫使运动剧烈区域的渲染质量提升。设计动机在于：标准~$L_1$~损失的均匀惩罚使高频结构和低频平坦区域获得相似权重，但高频区域正是欠重建最明显的位置；变形场网络对快速运动的组织重建不足，而光流天然编码了帧间运动幅度，可以自然定位运动密集区域。两个模块仅在训练阶段施加正则化，不引入额外推理开销。

总损失函数超参数设置为~$\lambda_{shf}=1$、$\lambda_{thf}=10$（THF~权重最大，表明光流约束对动态场景最为关键），在~EndoNeRF~和~SCARED~基准上分别达到~PSNR 38.14/27.47、SSIM 0.971/0.796，同时保持约~70 FPS~的实时渲染速度。EndoWave（2025.10）~\cite{2510_23087} 同样基于~ST-Endo4DGS~的~4D~高斯基元框架，但从频域监督角度提出了与~ST-Endo4DGS~互补的技术路线。ST-Endo4DGS~通过柱面谐函数在参数空间内隐式建模时空外观演化，而~EndoWave~认为内窥镜图像的关键信息——如镜面反射尖峰和平滑反照率——分布于不同频段，标准的~$L_1$~或~$L_2$~光度损失对它们施加均匀权重，无法在多尺度上精细约束重建质量。

为此，EndoWave~引入两项核心差异：(1) \textbf{时间演化外观柱面谐函数（TEASH）}——在~ST-Endo4DGS~的柱面谐函数基础上，为每个时间频率分量附加可学习相位~$\phi_n$，将固定的~$\sin(\omega_n t)$~推广为~$\sin(\omega_n t + \phi_n)$，使时间基函数能够与观测信号的相位对齐，在不增加系数数量的前提下提升时间拟合精度；(2) \textbf{有理小波分解损失}——采用有理尺度因子~$a = p/(p+1)$~替代传统的二进（dyadic）小波，生成比~$2^j$~更密集的尺度序列，从而更好地分离内窥镜图像中窄带镜面反射与低频反照率；对渲染图像和~ground truth~图像分别执行两级可分离~2D~有理小波分解，在~LL、LH、HL、HH~四个子带上分别计算~$L_2$~损失，以各子带独立权重~$\lambda_x$~控制约束力度。此外，EndoWave~引入光流几何约束——通过~GMFlow~计算~ground-truth~光流作为伪监督，对~4D~高斯基元投影后的~2D~位移施加~$L_2$~流损失，并使用双向一致性检查生成有效性掩码过滤不可靠对应点，确保流约束仅在可靠区域生效。

设计动机在于：ST-Endo4DGS~的~ENAC~法向约束仅作用于几何层面，而小波分解从外观的频率结构出发提供互补约束——高频子带聚焦边缘和纹理重建，低频子带保证全局色调一致。总体损失为~$\mathcal{L} = \lambda_{rgb}\mathcal{L}_{rgb} + \lambda_{depth}\mathcal{L}_{depth} + \lambda_{flow}\mathcal{L}_{flow} + \lambda_{wavelet}\mathcal{L}_{wavelet}$，在~EndoNeRF~和~SCARED~基准上实现了动态内窥镜场景的多尺度一致重建。

针对内窥镜特有的光照变化问题，一系列方法提出专门的解决方案。内窥镜的光照挑战源于其独特的光学配置：光源与相机共轴安装，随着探头移动，光源到组织的距离持续变化，导致照射强度呈平方反比衰减；同时湿润组织表面产生的镜面反射使表面外观高度视角依赖。ColIAGS（2025.10）~\cite{2510_18739} 针对结肠镜中移动光源导致的光照衰减，在锚基神经高斯（Scaffold-GS）框架上引入专门的照明模型。与~ST-Endo4DGS~的柱面谐函数隐式编码时间外观不同，ColIAGS~从物理分析入手：结肠镜中图像亮度同时受光源-组织距离~$d_l$（径向衰减，~$\propto d_l^{-1}$）和光线入射角~$\theta$~（光源扩散函数与相机渐晕的联合效应，~$\propto \cos^k\theta$）控制，两者可近似为统一角度依赖关系~$I \propto \cos^k\theta$。

由于直接估计指数~$k$~受相机自动增益干扰且高斯基元几何属性不精确，ColIAGS~不显式估计~$k$，而是通过两项改进隐式建模光照衰减：(1) \textbf{余弦嵌入几何建模（IMVE）}——在锚基高斯的几何预测~MLP~中，将相机到锚点的方向~$v$~经余弦基函数~$\Psi(v)$~投射到高维空间，使几何属性（位置、缩放、旋转）的预测能够感知视角变化，间接提升几何精度以满足光照建模的表面假设；(2) \textbf{光照衰减外观建模（IMIA）}——在颜色预测~MLP~中，将~$\Psi(v)$~和距离~$d$~同时作为输入，由~MLP~隐式学习~$\cos^k\theta$~衰减行为。这一设计的关键差异在于：ST-Endo4DGS~用柱面谐函数统一编码视角和时间外观，而~ColIAGS~将光照衰减从外观中解耦为独立的隐式分量，并通过物理近似将两个衰减源（径向距离、角度依赖）统一为单一角度参数化。训练损失包含标准~$L_1$~和~D-SSIM~重建损失、深度约束~$L_d$（对齐渲染深度与估计深度）以及缩放正则~$L_{scale}$（防止高斯过度重叠）。Endo-4DGX（2025.06）~\cite{2506_23308} 处理极端光照（低光和过曝光）条件下的重建，在~4DGS~基线上引入层次化光照校正管线。

与~ColIAGS~从物理模型出发建模稳态光照衰减不同，Endo-4DGX~聚焦于内窥镜手术中动态变化的极端光照场景——术中探头的移动导致同一场景在不同帧中出现低光和过曝光交替出现。其核心差异包含三级处理：(1) \textbf{每帧光照嵌入}——为每个训练视角初始化可学习嵌入~$e \in \mathbb{R}^{N \times 32}$，联合高斯基元一起优化，使模型能够记忆每个视角的特定光照状态；(2) \textbf{区域感知增强}——将训练图像按亮度中值分为"亮"和"暗"两级，分别训练两个独立的~3~层~MLP~遮蔽网络，对高斯的视角相关颜色施加仿射变换~$c_{tone} = c \cdot \gamma + \beta$（其中~$\gamma, \beta$~由网络根据嵌入~$e$~预测），解决单一网络处理两种极端光照时的训练不平衡和收敛失败问题；(3) \textbf{空间感知调整}——在区域增强之上，使用~3~层~MLP~（tanh~激活）从全局嵌入生成二次曲线映射，对渲染输出施加像素级亮度校正~$\hat{C}_{tone} = C_{tone} + C_{tone}(1 - C_{tone}) \cdot \delta$，模拟~Zero-DCE~的全局曝光控制。此外，Endo-4DGX~引入全局曝光控制损失~$\mathcal{L}_{Exposure} = \|\text{pool}_{1d}(\frac{1}{3}\sum_{j \in \{R,G,B\}} C_j) - E\|^2$（~$E=0.6$~），确保即使输入过曝或欠曝，重建结果仍保持正常曝光水平。

设计动机在于：ColIAGS~假设光照衰减可由单一物理模型描述，但极端光照条件超出了物理模型的适用范围，需要数据驱动的自适应校正。PR-ENDO（2024.11）~\cite{2411_12510} 利用内窥镜中单光源与相机紧密对齐的物理特性，将光照效果与组织固有的反照率（albedo）属性分离，实现物理基础的可重照明~3DGS——即在同一重建模型上模拟不同光照条件下的组织外观。与前述方法的关键差异在于：ColIAGS~通过隐式~MLP~建模光照衰减、Endo-4DGX~通过数据驱动校正极端亮度，两者均不具备在不同光照条件下重新渲染场景的能力；PR-ENDO~则为每个高斯基元赋予完整的~BRDF（Bidirectional Reflectance Distribution Function，双向反射分布函数）~材质属性——反照率~$a$、粗糙度~$r$~和基础反射率~$F_0$（对结肠组织~$F_0$~被限制为~0.03，对应真实折射率），结合~Cook-Torrance~微面元模型显式计算漫反射和镜面反射分量。

由于内窥镜光源近似点光源且与相机共置，PR-ENDO~用~MLP~增强传统~BRDF~的漫反射辐照度预测（输入高斯法向、光源方向和距离，并注入随机噪声提升泛化性），而非依赖解析光照系数。光源位置（相对相机中心微小偏移）与材质属性联合优化，训练损失包含图像重建、SSIM、深度、法向一致性，以及两个新引入的损失：漫反射一致性损失（约束~MLP~预测的辐照度与经典光传输方程的计算结果一致）和组织一致性损失（正则化同一组织区域的~$a$、$r$、$F_0$~趋于均值）。设计动机在于：物理基础的~PBR（Physically Based Rendering，基于物理的渲染）~模型允许分离光照和材质，使重建后的场景可在任意新光照条件下渲染——这对术中不同光源配置下的视觉一致性至关重要，是~ColIAGS~和~Endo-4DGX~的隐式方法无法实现的。

EndoGSim（2026.05）~\cite{2605_16022} 将~4DGS~与多模态大语言模型（MLLM）引导结合，通过物体级材料场和可微物质点法（MPM）实现内窥镜场景的物理感知重建和仿真。现有的内窥镜场景重建方法（如~Endo-4DGS、ForPlane~等）主要聚焦于视觉外观重建——利用~4DGS~表示动态手术场景的形变和外观，但无法描述组织的物理属性（如弹性模量、泊松比），因此无法驱动物理仿真来模拟器械与组织的交互力学行为。通用场景的物理仿真方法（如~PhysGaussian、Physics3D）虽将~3DGS~与可微~MPM（Material Point Method，物质点法）~结合，但依赖手动设置材料参数，且仅处理单一材质，无法应对手术场景中组织与器械的多种材质差异。

EndoGSim~的核心改进是提出物体级材料场（object-wise material field）：将场景中每个分割出的物体（软组织、手术器械）赋予独立的物理参数~$\theta_p = \{E, \nu\}$（Young's~模量~$E$~和泊松比~$\nu$），并通过两阶段策略自动推断——先由~MLLM（GPT-4o）根据分割图像粗估初始值，再通过可微~MLS-MPM~联合渲染损失和光流损失进行梯度优化。在技术实现上，EndoGSim~首先利用~Surgical-SAM-2（面向手术场景的专用分割模型）~对视频帧进行语义分割（区分组织~$M_{tissue}$~和器械~$M_{tools}$），然后用~$\pi^3$~模型估计深度；然后将分割掩码区域与原始图像一起输入~GPT-4o，获取各物体的初始材料参数；对于刚性器械，参数直接冻结，仅对可变形组织进行后续优化。优化过程采用可微~MLS-MPM~求解连续介质力学方程——粒子到网格（P2G）传递质量与动量，网格更新积分内力和外力，再通过网格到粒子（G2P）回传更新粒子状态。

联合损失函数为~$\mathcal{L} = \sum_t \|I_t - \hat{I}_t\|_2^2 + \lambda \sum_t \|U(I_t, I_{t+1}) - \hat{U}(\hat{I}_t, \hat{I}_{t+1})\|_2^2$，其中~$\lambda = 0.1$~平衡渲染外观与光流运动两项监督。物体级材料场使每个物体拥有独立的物理属性，MLLM~初始化提供了物理合理的起点（避免优化陷入局部最优），可微~MPM~则使材料参数能够通过视觉监督信号端到端优化，无需任何物理测量设备。消融实验表明：移除材料场后~RE~从~0.178~升至~0.433，EPE~从~0.132~升至~0.366，FID~从~4.01~升至~8.81。Feature-EndoGaussian（2025.03）~\cite{2503_06161} 首次实现实时语义分割与变形手术场景重建的统一管线，以~61 FPS~的渲染速度在~EndoNeRF~上达到~39.1 PSNR。与~ST-Endo4DGS~聚焦于动态外观的时空编码不同，Feature-EndoGaussian~的核心创新是将语义特征蒸馏嵌入~4DGS~的变形和渲染管线中，使每个高斯基元不仅携带颜色和几何信息，还携带来自~SAM（Segment Anything Model，通用图像分割模型）的~256~维语义特征。

其技术路线包含三个关键差异：(1) \textbf{特征-时空变形模块}——在~HexPlane~编码后，四路轻量几何变形分支分别预测位置~$\Delta\mu$、旋转~$\Delta R$、缩放~$\Delta S$~和不透明度~$\Delta o$~的更新，同时新增语义更新网络~$F_{feat}$~将四路变形特征拼接后预测每高斯~128~维语义特征增量~$\Delta z$，实现语义特征随组织变形同步更新；(2) \textbf{并行语义渲染}——修改可微光栅化器，在渲染颜色和深度的同时通过~$\alpha$-混合~渲染语义特征图~$S(x) = \sum_i z'_i \alpha_i \prod_{j=1}^{i-1}(1-\alpha_j)$，随后经~CNN（Convolutional Neural Network，卷积神经网络）~解码器上采样并与~SAM~输出的高维语义特征对齐，以逐像素~$L_1$~损失监督；(3) \textbf{粗到精训练策略}——粗阶段仅优化基础高斯参数，精阶段激活语义分支联合优化，避免语义监督信号在几何未稳定时干扰场景表示。设计动机在于：术中场景理解不仅需要重建外观，还需要实时识别和分割组织与器械；将语义特征嵌入高斯基元而非后处理分割，使语义推理与渲染同步完成，实现了~61 FPS~的端到端实时性能，推理时直接通过~SAM~解码器生成实例分割掩码。

\subsubsection{稀疏视角与SLAM系统}

本路线介绍稀疏视角条件下内窥镜重建和基于 3DGS 的 SLAM（Simultaneous Localization and Mapping，同步定位与建图）系统。

临床内窥镜常受限于稀疏视角输入，标准~3DGS~在此条件下重建质量显著下降。EndoSparse（2024.07）~\cite{2407_01029} 系统性地解决了稀疏视角（仅~3~个视角）内窥镜重建的核心难题。在基线做法中，标准~3DGS~依赖密集视角提供充分的多视角约束——高斯基元的位置和形状通过多视角光度一致性联合优化。当视角极度稀疏时，大量高斯缺乏充分观测，几何退化严重：不可见区域的基元自由漂移、表面法向不一致、渲染图像出现明显伪影。

EndoSparse 的核心改进是引入两类互补的先验知识弥补视角不足：\textbf{扩散先验}（Diffusion Prior）利用预训练的扩散模型对稀疏视角合成的图像施加正则化——将渲染图像送入扩散模型，通过~DDIM（Denoising Diffusion Implicit Models，去噪扩散隐式模型）~采样过程将其"投影"回自然图像流形，确保合成结果的视觉合理性，避免稀疏视角导致的退化解；\textbf{深度先验}利用基础模型（如~Depth-Anything）的单目深度估计能力——对每个输入视角估计稠密深度图，通过尺度对齐后作为高斯基元初始化和几何正则化的监督信号。两路先验互补：扩散先验约束外观合理性，深度先验约束几何正确性。仅用~3~个视角即可获得稳健的几何和外观质量，实现实时渲染。Gaussian Pancakes（2024.04）~\cite{2404_06128} 同样从稀疏视角挑战出发，但与~EndoSparse~的扩散+深度双先验策略不同，它采用~SLAM~驱动的增量重建路线。

Gaussian Pancakes~的核心差异在于将~RNN-SLAM~系统与~3DGS~渲染管线耦合：RNN-SLAM~提供实时相机位姿估计和深度预测（通过循环神经网络），3DGS~负责可微光栅化渲染和场景表示。其技术贡献包含两方面：(1) \textbf{深度正则化}——利用~RNN-SLAM~预测的深度图作为高斯基元深度渲染的监督信号，对齐渲染深度与~SLAM~深度，约束几何在稀疏视角下的合理性；(2) \textbf{几何约束}——引入法向一致性损失，对齐渲染深度导出的法向与~SLAM~深度导出的法向，防止稀疏视角下高斯基元的自由漂移。与~EndoSparse~需要预训练扩散模型不同，Gaussian Pancakes~的深度和位姿先验由轻量~RNN~在线生成，无需额外的大型基础模型，在结肠镜数据上实现~PSNR 18\%~的提升和~$100\times$~渲染加速。

针对内窥镜~SLAM~的特定应用需求，EndoFlow-SLAM（2025.06）~\cite{2506_21420} 系统性地提出了基于光流约束的~3DGS~框架。在~Baseline~做法中，现有~3DGS-SLAM~方法（如~EndoGSLAM、MonoGS）仅依赖光度损失和深度损失联合优化高斯基元和相机位姿。然而在内窥镜场景中存在两个核心挑战：非朗伯表面（湿润组织镜面反射）导致光度不一致，纯外观约束在视角变化剧烈处产生歧义；呼吸运动引起的组织动态变形影响跟踪稳定性。EndoFlow-SLAM~的核心改进是引入三个互补约束机制。首先，光流损失作为几何约束——计算每个高斯基元在相邻时刻之间的~"Gaussian flow"（高斯流），通过~alpha~加权混合得到复合流，以~$\mathcal{L}_{flow}$~约束其与预计算光流的一致性，在关键帧优化时同时约束~3D~场景结构和相机运动。

其次，深度正则化策略应对非朗伯面的光度不一致——采用深度梯度正则化对齐估计深度图与~ground truth~深度图的水平和垂直梯度，并结合尺度不变损失缓解单目深度估计的尺度歧义。第三，改进的全局精调策略——采用两阶段策略：第一阶段优先使用关键帧对应视角结合深度损失优化全局场景，第二阶段随机选择视角确保全面精调。设计动机在于：光流天然编码帧间运动幅度，既能约束~3D~结构一致性也能约束相机运动，弥补纯外观约束在非朗伯面处的不足；深度梯度正则化确保~3DGS~深度渲染在内窥镜场景中的有效性。NRGS-SLAM（2026.02）~\cite{2602_17182} 提出变形感知的~3D~高斯地图，为每个高斯基元附加可学习的变形概率，通过贝叶斯自监督策略优化，在单目非刚性~SLAM~中将位姿误差降低~50\%。

与~EndoSparse~从稀疏视角的先验注入不同，NRGS-SLAM~解决的是单目非刚性~SLAM~中的核心矛盾：内窥镜场景中刚性结构（如骨骼、器械）和可变形组织（如软组织蠕动）并存，而标准~3DGS-SLAM~对所有区域施加统一的光度约束，导致变形区域的帧间光度变化被错误归因于相机运动而非组织变形，严重影响位姿估计精度。

NRGS-SLAM~的核心创新包含三个层面：(1) \textbf{变形概率}——为每个规范空间高斯赋予可学习标量~$w_{d,i} \in [0,1]$（通过~sigmoid~参数化），$w_{d,i} \to 0$~表示刚性基元，$w_{d,i} \to 1$~表示完全可变形基元；通过可微光栅化合成~2D~变形置信度图~$M_{def}(u) = \sum_i w_{d,i} \alpha_i \prod_{j=1}^{i-1}(1-\alpha_j)$，在相机位姿估计时对高变形概率像素降权，使跟踪锚定在稳定解剖结构上；(2) \textbf{1D~高斯基函数变形场}——不同于~ST-Endo4DGS~的柱面谐函数或隐式~MLP~变形，NRGS-SLAM~用显式的~1D~高斯基函数加权求和参数化时间变形~$\Delta A_i(t) = \sum_{k=1}^{K^A} w_{i,k}^A \psi(t; \mu_{i,k}^A, \sigma_{i,k}^A)$，支持动态地图管理（新增/删除基函数），且变形概率作为软门控调制~$A_i(t) = A_i^c + w_{d,i} \cdot \Delta A_i(t)$，确保刚性区域的变形被抑制；(3) \textbf{贝叶斯自监督}——由于无~ground-truth~变形标签，通过单目残差推导后验变形似然作为伪标签监督~$w_{d,i}$~的学习。Endo-2DTAM（2025.01）~\cite{2501_19319} 基于~2DGS~构建表面法向感知管线，结合点对点和点对面距离度量实现实时跟踪，深度重建~RMSE~达~1.87 mm。

与~EndoSparse~从稀疏视角重建和~NRGS-SLAM~的非刚性变形建模不同，Endo-2DTAM~聚焦于~RGB-D SLAM~中的跟踪精度问题，其核心差异在于采用~2DGS~而非~3DGS~作为场景表示。2DGS~将高斯基元约束为平面椭圆（缩放矩阵~$S_i$~的第三维度为零），天然提供表面法向信息，避免了~3DGS~中法向需要从深度梯度间接推导的问题。在跟踪模块中，Endo-2DTAM~同时使用投影点对点距离~$d_{point\text{-}to\text{-}point} = \|D_{GT} - D\|_1$~和法向感知的点对面距离~$d_{point\text{-}to\text{-}plane} = \|(x_{GT} - x)^T N_{GT}\|_1$~（其中~$N_{GT}$~由~GT~深度的有限差分计算），后者约束渲染表面与~GT~表面切平面对齐，在法向方向上提供比纯深度约束更精细的几何约束。映射阶段引入深度扭曲损失（最小化沿射线高斯交点的深度间隔）和法向一致性损失（约束~2D~高斯法向与~GT~法向对齐）。此外，Endo-2DTAM~设计了位姿一致性关键帧采样策略——以平移距离、旋转误差和时间间隔的加权和计算各候选帧的选择概率，确保全局优化时采样到空间和时间上分布均衡的关键帧集合。

联合~Bundle Adjustment~每~100~次跟踪迭代执行一次，同时优化相机位姿和~2D~高斯场景表示。

\subsection{专用场景与变形建模}
\label{ss:19_2}

本节讨论两条技术路线：（1）专用内窥镜场景；（2）变形组织重建。

\subsubsection{专用内窥镜场景}

本路线介绍针对结肠镜、手术导航等特定内窥镜场景的专用重建方法。

不同内窥镜场景有各自独特的技术需求。

\textbf{ ColonSplat}（2026.03）~\cite{2603_06860} 修改了~3DGS~的动态变形建模环节——用时间相关的空间坐标偏移~$\Delta\mathbf{x}$~替代仅调整旋转/缩放/透明度的隐式变形策略——可微光栅化渲染管线保持不变。

\textit{问题与动机。}
结肠镜在狭窄管状结构中行进，相机持续平移且组织发生非刚性蠕动。现有动态内窥镜方法（如~EndoGaussian、Endo-4DGS）在处理此类场景时，\emph{不更新}高斯基元的空间坐标~$\mathbf{x}$，而是通过旋转~$\mathbf{r}$、缩放~$\mathbf{s}$~和透明度~$\alpha$~的调整来隐式编码变形。这种做法在宽腔体内窥镜中尚可接受，但在结肠场景中会产生物理上不合理的运动重建——高斯通过``旋转''和``缩放''而非``移动''来近似蠕动，全局视角下出现明显的几何不一致。ColonSplat~首次系统性地揭示了这一失败模式，并指出：结肠蠕动是真三维空间位移，必须显式更新~$\mathbf{x}$~坐标。为支持严格的全局几何评估，该工作还提出了~DynamicColon~合成数据集，提供每个时间步的~ground-truth~点云。

\textit{变形建模。}
ColonSplat~的核心是用一个多层感知机~(MLP)~$\theta$~预测时间相关的几何更新。直觉上，每个高斯基元在~``规范空间''~中有一个初始位置，MLP~根据时间~$t$~告诉它``移动到哪里''。具体地，给定归一化时间步~$t \in [0,1]$~和规范空间的高斯参数，变形网络预测三维偏移量：

\begin{equation}
    (\Delta\mathbf{x},\, \Delta\mathbf{s},\, \Delta\mathbf{r}) = \theta\!\bigl(\mathrm{G}(t),\, \mathrm{G}(\mathbf{x})\bigr),
\end{equation}

其中~$\mathrm{G}$~表示~HexPlane~网格编码，将时间和空间坐标映射为高维特征。变形后的几何参数通过\emph{加法}方式获得：

\begin{equation}
    \mathbf{x}' = \mathbf{x} + \Delta\mathbf{x},\quad
    \mathbf{s}' = \mathbf{s} + \mathrm{clip}(\Delta\mathbf{s}, -\tau_s, \tau_s),\quad
    \mathbf{r}' = \mathbf{r} + \mathrm{clip}(\Delta\mathbf{r}, -\tau_r, \tau_r).
\end{equation}

即高斯的平移是自由的，而缩放和旋转的偏移被裁剪阈值~$\tau_s$、$\tau_r$~限制在合理范围内。这个设计的核心动机是：\emph{强制网络通过空间位移~$\Delta\mathbf{x}$~来表达变形}，防止网络走捷径——通过大幅改变旋转或缩放来``补偿''空间变形的不足。此外，透明度~$\alpha$~在所有时间步保持固定，避免网络通过``隐藏/显示''高斯来规避真实的空间运动建模。

\textit{颜色建模。}
内窥镜成像中，相机与光源共置，导致强烈的视角相关外观变化，加上体液和组织（如息肉）的局部反射特性差异，颜色不能简单视为静态属性。ColonSplat~为此设计了独立的颜色分支~$\theta_{\mathrm{color}}$，利用每个高斯的可学习嵌入向量~$\mathbf{e}_i \in \mathbb{R}^d$~捕获局部时间外观变化。颜色更新采用\emph{乘法}形式~$\mathbf{c}' = \mathbf{c} \odot (1 + \Delta\mathbf{c})$，直觉上模拟了光-表面交互的乘性调制效应，而非简单地叠加偏移。

\textit{正则化约束。}
为确保变形在局部空间上的一致性，ColonSplat~引入~$K$-最近邻变形一致性损失。对每个高斯~$i$，在规范空间中找到~$K$~个最近邻~$\mathcal{N}_i$，惩罚其变形后位置与邻域均值的偏差：

\begin{equation}
    \mathcal{L}_{\mathrm{knn}} = \frac{1}{N}\sum_i
    \Bigl\|\mathbf{x}'_i(t) - \frac{1}{|\mathcal{N}_i|}\sum_{j\in\mathcal{N}_i}\mathbf{x}'_j(t)\Bigr\|_2^2.
\end{equation}

直观理解，这个损失确保相邻组织一起运动，避免个别高斯``跳走''产生不连续变形。此外还包含颜色偏移正则~$\mathcal{L}_{\mathrm{co}}$（防止通过颜色变化伪解释运动）、颜色方差正则~$\mathcal{L}_{\mathrm{cv}}$（避免未观测区域的颜色伪影）、深度正则~$\mathcal{L}_{\mathrm{depth}}$（归一化深度图的~$\ell_1$~损失）和全变分平滑~$\mathcal{L}_{\mathrm{TV}}$。总训练目标为：

\begin{equation}
    \mathcal{L} = \mathcal{L}_{\mathrm{rgb}} + \lambda_{\mathrm{TV}}\mathcal{L}_{\mathrm{TV}}
    + \lambda_{\mathrm{knn}}\mathcal{L}_{\mathrm{knn}}
    + \lambda_{\mathrm{depth}}\mathcal{L}_{\mathrm{depth}}
    + \lambda_{\mathrm{co}}\mathcal{L}_{\mathrm{co}}
    + \lambda_{\mathrm{cv}}\mathcal{L}_{\mathrm{cv}}.
\end{equation}

ColonSplat~在~C3VDv2~和~DynamicColon~数据集上均取得了优于现有动态内窥镜方法的几何保真度。SurGSplat（2025.06）~\cite{2506_05935} 针对术中无位姿标注的手术导航场景，提出了从局部到全局的渐进~3DGS~重建范式。与~ColonSplat~使用已知相机位姿进行~4D~变形建模不同，SurGSplat~面临更基本的挑战：术中场景无法预运行~SfM~获取位姿和初始点云，必须从单帧出发同时估计相机位姿和重建场景。

其核心技术路线包含四步：(1) \textbf{单帧初始化}——使用~DepthNet~估计首帧单目深度，假设单位位姿，将深度图反投影为初始~3D~高斯集，仅通过光度损失优化高斯属性；(2) \textbf{相对位姿估计}——将相机位姿估计转化为~3D~高斯核的仿射变换优化（数学上两者等价），固定高斯属性，仅优化仿射变换~$A_t$~使渲染结果与下一帧匹配；(3) \textbf{全局渐进增长}——将每帧的局部高斯场景通过相对位姿逐步融合到全局场景中，利用帧间视差梯度信息识别欠重建区域进行致密化，实现从稀疏到完整的场景覆盖；(4) \textbf{几何约束优化}——引入两个互补约束：投影几何一致性损失（PGC）利用密集特征匹配网络建立帧间~2D~对应关系，通过~3D~高斯投影计算重投影误差，确保~2D~特征点与~3D~空间投影一致；深度几何相关性损失（DGC）在每个随机采样的局部图块内计算~DepthNet~深度与渲染深度的~Pearson~相关系数，约束相对深度趋势一致而不强求绝对尺度对齐。

设计动机在于：内窥镜场景的纹理稀疏和结构复杂使传统~SfM~频繁失败，SurGSplat~通过交替优化位姿和场景（而非~SfM~先行、重建后行）实现了端到端的术中重建，特别适合血管结构等关键特征的高保真重建。EndoLRMGS（2025.03）~\cite{2503_22437} 将大重建模型（LRM）与~GS~结合用于完整手术场景重建，解决~ColonSplat~等纯~GS~方法无法处理器械等封闭物体的~3D~完整性问题。ColonSplat~仅用高斯基元重建组织表面，对于手术器械等需要完整~3D~几何（包括不可见面）的物体无能为力。EndoLRMGS~的核心差异是引入混合表示策略：\textbf{EndoGaussian~重建组织（表面物体）}，\textbf{LRM~重建器械（封闭物体）}。

具体而言，首先使用~DEVA~分割网络对每帧中的手术器械生成唯一掩码，区分组织与器械；对于组织区域，沿用~EndoGaussian~的反投影初始化和高斯渲染方案，以工具掩码排除器械区域后训练组织的高斯表示；对于每个器械，利用~DINOv2（Distillation with NO labels v2，自监督视觉Transformer）~将分割掩码转换为~patch~特征~token~，投影为三平面表示，再经~MLP~预测~384$^3$~分辨率的~3D~密度场和颜色场，最终通过~Marching Cubes~提取完整水密网格。LRM~生成的器械~3D~模型在固定网格内缺乏物理尺度和空间定位，为此~EndoLRMGS~提出\textbf{正交透视联合投影优化（OPjPO）}：利用~GS~重建的组织点云构建正交投影矩阵，将器械和掩码区域的组织投影到~2D~平面（消除透视畸变），通过面积比估计器械的缩放因子~$\gamma = A_{mask}/A_{tool}$，然后在透视投影下通过最大化器械投影与实际掩码的~IoU~沿~$x$、$y$、$z$~三轴精调器械位置。

设计动机在于：术中场景同时包含变形组织和刚性器械，两者需要截然不同的~3D~表示——高斯基元擅长开放表面渲染，但无法保证封闭物体的拓扑完整性；LRM~擅长从单视图生成完整~3D~模型，但无法直接与~GS~场景配准。EndoLRMGS~通过~OPjPO~弥合两种表示的尺度-位置鸿沟，输出组织-器械一体的完整手术场景~3D~模型。

此外，Endo-G$^{2}$T（2025.11）~\cite{2511_21367} 提出几何引导和时间感知的训练方案，在~ColonSplat~的~4D~变形建模基础上引入三项互补改进。ColonSplat~依赖已知相机位姿进行变形训练，而~Endo-G$^2$T~面向无位姿或位姿不准确的通用内窥镜场景。

其核心差异为：(1) \textbf{几何引导先验蒸馏（GPD）}——利用单目深度估计（附置信度图~$C^*$）作为外部先验，通过帧自适应置信度阈值和物理深度范围约束定义有效像素集，仅在有效像素比例~$\geq 10\%$~时激活先验损失，采用~warm-up-to-cap~调度使先验权重从零渐进增长至上限值，避免训练初期对噪声深度估计的过拟合；先验损失包含尺度不变的~log~深度差（SILog）和深度梯度对齐两项，分别在全局尺度和局部边缘层面约束几何；(2) \textbf{时间嵌入高斯场（TEGF）}——为每个高斯基元维护速度~$v_i(t)$~和旋转增量~$\omega_i(t)$，位置通过速度积分更新，旋转通过最小旋转子（rotor-like）操作更新，保持旋转的流形约束；引入不透明度熵损失鼓励~$\alpha_i$~趋向~0~或~1（消除半透明模糊基元），以及局部速度相干性损失平滑时空邻域的运动场；(3) \textbf{关键帧约束流式处理}——按步长~$w$~将视频划分为关键帧和候选帧，关键帧执行完整优化（含致密化和剪枝），候选帧仅执行轻量图像空间更新且保持高斯数量固定；全局高斯预算~$G_{max}$~在任何时刻约束活跃基元总数，防止长序列的内存膨胀。

SAGS（2025.10）~\cite{2510_27318} 引入注意力驱动的动态加权~4D~变形解码器，从变形场网络架构层面改进~ColonSplat~和~ST-Endo4DGS~使用的固定权重~MLP~解码器。SAGS~的核心差异在于提出\textbf{自适应加权变形解码器}：将多头自注意力机制与~MLP~分支通过可学习的动态权重~$\lambda_1$、$\lambda_2$~进行融合，使模型在内窥镜场景的大尺度组织运动和细微纹理变化之间自适应切换——自注意力捕获长程依赖保持全局几何一致性，MLP~处理局部非线性变形保留精细细节。解码过程采用双阶段仿射变换（Affine）稳定特征缩放，残差输出通过~$\lambda_1$、$\lambda_2$~加权后直接更新高斯的位置、旋转、缩放、球谐系数和不透明度。此外，SAGS~针对内窥镜场景特有的视觉伪影（镜面反射高光、窄立体基线和快速非刚性运动引起的锯齿和闪烁）引入抗锯齿处理：\textbf{3D~平滑滤波器}根据~Nyquist~采样定理对每个高斯基元施加低通正则化~$G_k(x)_{reg} = (G_k * G_{low})(x)$~抑制体积域高频伪影；\textbf{2D Mip~滤波器}在屏幕空间对每个像素区域进行盒式滤波近似，消除缩放视图或相机距离变化时的锯齿。

两个滤波器联合作用于~HexPlane~的六面体特征，确保~4D~变形场在空间和时间维度上均保持抗锯齿质量。Hayoz et al.（2024.09）~\cite{2409_06037} 提出基于控制点的在线~3D~重建和密集表面跟踪框架，与~ColonSplat~通过~HexPlane~编码时空特征驱动全局变形场不同，Hayoz et al.~采用稀疏控制点驱动的局部变形建模。其核心差异在于将规范场景（canonical scene）通过~3DGS~建模，而组织变形通过一组控制点~$D = \{(p_k, t_k, q_k)\}_{k=1}^K$~的平移和旋转偏移来表达——每个控制点是一个锚高斯，平移场~$\tau^t(x)$~和旋转场~$\tau^q(x)$~通过高斯核加权平均插值~$\tau^t(x) = \frac{1}{\sum_k w(x, p_k)} \sum_k w(x, p_k) t_k$（$w$~为高斯距离衰减核），确保变形在空间上平滑传播。

在线拟合分三步进行：(1) \textbf{规范场景扩展}——对新帧中不透明度低于~0.95~的像素添加新高斯，位置由反投影确定，尺度设为最近高斯距离，渐进覆盖未见区域；(2) \textbf{控制点初始化}——从高斯集合中随机采样锚点，用~RAFT~计算光流初始化平移偏移（闭式最小二乘解）；(3) \textbf{能量最小化}——外部能量（图像和深度的~MSE~）与内部能量（刚性项惩罚相邻锚点间相对位置/旋转变化、等距项鼓励局部距离保持、可见性项惩罚不可见区域变形漂移）联合优化。特别地，通过梯度调制因子~$\eta_i = 2(1 - \sigma(c_1 v_i - c_2))$~随更新次数~$v_i$~增加逐步冻结规范场景参数，防止漂移。最终实现任意表面点的密集跟踪——通过找到最近高斯并随变形场传播其~3D~位置。Local-EndoGS（2026.02）~\cite{2602_17473} 处理任意相机运动的单目内窥镜序列，通过渐进窗口表示和粗到细策略实现高质量~4D~重建。与~ColonSplat~使用单一规范空间和全局变形网络不同，Local-EndoGS~的核心差异在于\textbf{渐进窗口全局场景表示}：将输入序列自适应划分为~$M$~个连续局部窗口，每个窗口由独立的规范空间和变形网络建模，逐窗口渐进优化参数。

窗口划分策略同时考虑相机运动（平移超过阈值~$\epsilon_t$~或旋转角超过~$\epsilon_r$~时创建新窗口）和帧内容变化（当前帧与窗口首帧的~RGB~差异超过~0.05~时切分），确保窗口内视角一致且内容相似。每个局部窗口使用~EH-SurGS~的变形模型，包含~Gaussian~基函数的时间变形（支持不可逆变形如组织剪切）和生命周期机制（高斯仅在有效时间范围内激活）。针对单目序列的尺度歧义和稀疏~SfM~问题，Local-EndoGS~提出\textbf{粗到精初始化}：粗阶段利用~Track-Any-Point（TAP）模型建立跨帧像素级时序轨迹（取代不可靠的~SIFT~匹配），通过多视角三角化生成尺度一致的密集点云，同时从前一窗口传递已优化的局部模型作为跨窗口先验；精阶段计算光度误差图，对高误差区域用尺度对齐后的单目深度估计替换初始点云中的对应点。

训练损失设计三项：(1) 渲染损失（~$L_1$ + D-SSIM），不含深度监督（实验表明深度先验对该框架有负面影响）；(2) 2D~跟踪损失，利用~TAP~模型的像素轨迹约束规范空间与观测空间的~2D~一致性；(3) 物理正则化，包含局部刚性损失、局部旋转相似损失和长期等距损失，通过~kNN~高斯权重在邻域内施加物理约束。Beltran et al.（2026.04）~\cite{2604_28179} 将呼吸建模嵌入~CT~知情的~3DGS~框架，用于动态支气管镜导航。该方法与~ColonSplat~的~MLP~驱动自由变形和~Hayoz et al.~的控制点驱动局部变形有根本性差异——Beltran et al.~将变形空间约束为一维呼吸相位子空间，消除了对几何正则化（ARAP、等距、刚性损失）的需求。

其核心技术路线包含三个层面：(1) \textbf{网格锚定高斯表示}——每个高斯基元通过可学习重心坐标~$b_k$~锚定到~CT~导出的气道网格~$M = (V, F)$~的父三角面上，位置~$\mu_k = \sum_{j=0}^{2} \beta_j v_j^{(f_k)}$~（$\beta$~经~softmax~约束在三角面内），法向对齐三角面法向且法向缩放固定为零形成薄盘基元；当网格随呼吸变形时，每个高斯自动继承父三角面的局部拉伸、旋转和平移，无需逐高斯变形参数；(2) \textbf{CT~知情呼吸变形}——通过可变形~CT~配准获得吸气和呼气网格间的固定位移场~$\Delta = V_{exp} - V_{insp}$，任意呼吸相位~$\phi_t \in [0,1]$~的网格为~$V(\phi) = V_{insp} + \phi_t \Delta$，所有预测配置均位于患者特定解剖轨迹上；(3) \textbf{余弦-线性混合激活}——用~$(1-\epsilon)\frac{1}{2}(1-\cos\theta_t) + \epsilon\theta_t/\pi$~替代~sigmoid~参数化呼吸相位，余弦项提供与呼吸周期一致的平滑遍历，线性泄漏项（~$\epsilon=0.05$~）确保在吸气/呼气极值处梯度不为零，避免优化停滞。损失函数仅含光度损失和时间正则化（相邻帧~$\phi_t$~平滑），无需额外几何正则化。

此外，该工作还提出~RESPIRE~仿真框架，利用~Blender Cycles~路径追踪引擎（含次表面散射模拟黏膜半透明性和镜面反射模拟湿润表面眩光）生成具有真实呼吸变形和逐帧~GT~几何的支气管镜视频，填补了该领域无~GT~变形数据集的空白。Yang（2025.09）~\cite{2509_17027} 提出基于虚拟相机的正则化方法，从稀疏内窥镜视角高效重建可交互的手术仿真场景。与~ColonSplat~处理动态变形场景不同，Yang~聚焦于稀疏视角下的静态场景重建和后续物理仿真。其技术贡献分为重建和仿真两部分。在重建方面，Yang~针对标准~3DGS~在稀疏视角下过拟合的问题引入三级正则化：(1) 利用~Depth Anything-V2~的单目深度作为深度监督损失~$L_{depth}$；(2) 引入深度扭曲损失~$L_{dist} = \sum_{i,j} w_i w_j |z_i - z_j|$~最小化沿射线高斯交点的深度间隔，改善几何质量；(3) 引入深度图全变分（TV）正则化~$L_{TV}$~消除不自然的锐利边缘。

更进一步，Yang~提出\textbf{虚拟相机深度正则化}——在训练时从最近两个训练视角通过球面线性插值（Slerp）和线性插值（lerp）生成虚拟相机，在虚拟视角下仅施加~$L_{dist}$~和~$L_{TV}$~（无需~GT），强制高斯几何在训练视角之外的视角下也保持合理，防止过拟合到稀疏训练视图。在仿真方面，Yang~针对~PhysGaussian~直接对所有高斯施用~MPM~导致大规模高斯集仿真过慢的问题，提出\textbf{稀疏控制节点方法}：通过最远点采样（FPS）从重建的高斯集中提取少量控制节点送入~MPM~仿真，每个高斯的变形通过~KNN~找到~$K=4$~个最近控制节点，以其加权极分解变形梯度更新位置和协方差矩阵。这使仿真速度与高斯总数解耦，实现高效的实时软组织交互仿真。Chen and Wang（2026.02）~\cite{2602_20718} 引入基于符号距离场的表面感知重建和半刚性变形约束，实现单目内窥镜组织的~3D~重建。

与~ColonSplat~用自由变形场建模动态组织不同，Chen and Wang~的核心差异在于\textbf{将网格表面先验显式嵌入高斯表示}——对首帧使用~NeuS2~（带~SDF~和~Eikonal~正则化的哈希编码隐式表面模型）重建三角网格，然后将高斯基元一一绑定到网格三角面质心，初始半径设为三角面内切圆大小，建立高斯-三角面对应关系。这种绑定带来两个优势：(1) 变形时可直接通过三角面拓扑确定邻域，无需额外~kNN~计算；(2) 表面法向信息自然可用。对高斯施加尺度正则~$L_{scale}$~（约束高斯尺寸不超过绑定三角面外接圆半径）和位移正则~$L_{shift}$~（惩罚高斯偏离绑定三角面），维持局部连续性。变形建模采用\textbf{半刚性策略}：局部刚性约束使用~ARAP（As-Rigid-As-Possible）损失，仅在~SIFT~关键点周围的有限邻域内施加，利用~SVD~分解计算最优旋转矩阵后惩罚偏离刚性的变形；全局非刚性约束包含旋转相似性损失（鼓励~kNN~邻域高斯的旋转变化一致）和长期等距损失（保持规范空间和当前空间中邻域高斯间距离不变）。

此外，预处理阶段引入基于~Transformer~的视频修复网络去除器械遮挡（在~StereoMIS~数据上微调），为后续重建提供无遮挡的组织图像。设计动机在于：标准~3DGS~的离散高斯基元无法准确表示连续表面，通过网格-高斯绑定将~SDF~的表面精度与~GS~的渲染效率结合，同时~ARAP~和等距约束确保变形的物理合理性。

\subsubsection{变形组织重建}

本路线介绍手术场景中软组织动态变形的 3DGS 建模与重建方法。

Deform3DGS（2024.05）~\cite{2405_17835} 首次将~3DGS~引入手术场景的快速重建。它修改了~3DGS~的变形建模环节，用灵活变形建模（FDM）替代~HexPlane~隐式变形场和傅里叶/多项式显式变形；3D~高斯基元的光照/透明度属性和可微渲染保持不变。

\textbf{Baseline 做法及其局限。}动态场景重建的主流变形建模有两条路线：(1)~基于~HexPlane~的隐式变形场（如~4DGS、LerPlane），将时空信息编码为共享特征平面并通过~MLP~解码时间变形。虽然表示能力强，但特征平面插值和~MLP~推理的计算开销大，难以满足手术场景的实时需求。(2)~基于傅里叶/多项式基函数的显式变形，为每个高斯基元赋予可学习权重~$\omega$~对固定基函数进行线性组合来回归运动曲线~$\psi(t;\omega)=\sum_{j=1}^{B}\omega_j b_j(t)$。虽然计算高效，但存在两个根本局限：(i)~基函数在给定时间点取值固定、不随训练调整，时间变形被约束为规范运动轨迹的一致表示，模型被迫放弃某些特定运动以保持整体轨迹的连贯性，无法精确捕捉器械-组织交互等精细局部变形；(ii)~某一时刻学到的变形会全局影响整条轨迹，不同查询时间的局部变形表示相互干扰。

\textbf{改了什么：FDM 新机制。}FDM~的核心改进是将固定基函数替换为可学习参数化的~Gaussian~基函数——每个基函数的中心~$\mu$~和宽度~$\sigma$~均可在训练中优化，使变形曲线不再被固定函数形状束缚，而是自适应地调整每个基函数在时间轴上的"作用区域"和"影响范围"。直觉上，每个~Gaussian~基函数就像时间轴上的一个"可移动、可伸缩的灯泡"，灯泡中心~$\mu$~决定它照亮哪个时间段，灯泡宽度~$\sigma$~决定它照亮多宽的时间窗口，而权重~$w$~决定亮度——多个灯泡叠加即可拟合任意复杂的变形曲线。

\textbf{实现原理。}具体地，FDM~为每个高斯基元赋予一组~Gaussian~基函数（论文实验中取~$B=17$），以可学习中心~$\mu_j$~和方差~$\sigma_j$~参数化每个基函数：

\begin{equation}
    \tilde{b}(t;\mu_j,\sigma_j) = \exp\left(-\frac{1}{2\sigma_j^2}(t-\mu_j)^2\right),
\end{equation}

即每个基函数在时间轴上以~$\mu_j$~为中心、$\sigma_j$~为宽度呈钟形分布。时间变形曲线由~$B$~个基函数的加权线性组合给出：

\begin{equation}
    \Delta(t;\theta) = \sum_{j=1}^{B} w_j \cdot \exp\left(-\frac{1}{2\sigma_j^2}(t-\mu_j)^2\right),
\end{equation}

即总变形是各基函数在当前时刻~$t$~的响应值的加权和——只有中心靠近~$t$~的基函数才产生显著贡献，远离~$t$~的基函数几乎为零。FDM~分别为位置、旋转和缩放学习独立的变形参数集~$\theta^\mu$、$\theta^r$、$\theta^s$，最终通过~$\mu = \mu_c + \Delta^\mu(t;\theta^\mu)$、$r = r_c + \Delta^r(t;\theta^r)$、$s = s_c + \Delta^s(t;\theta^s)$~将规范高斯变换到时间相关状态，其中~$\mu_c$、$r_c$、$s_c$~为规范空间的初始属性。由于~Gaussian~函数的局部支撑特性，相邻时刻的变形保持连续性，而大时间间隔的变形几乎解耦——这使模型能独立捕捉不同时间段的精细变形动力学。

\textbf{设计动机。}手术场景的变形具有两个独特特征：(i)~器械-组织交互引起的变形高度局部化（仅发生在接触点附近），而~HexPlane~等全局时空分解将所有时空位置编码为共享特征平面，无法为局部变形分配额外的表示容量；(ii)~不同时间段的变形动力学可能截然不同（如器械夹持前的缓慢变形和夹持瞬间的急剧变形），需要时间轴上的自适应分辨率。FDM~选择高斯基函数正是利用其局部支撑特性——每个基函数仅影响时间轴上以~$\mu_j$~为中心、$\sigma_j$~为宽度的局部区域，使模型能独立捕捉不同时间段的精细变形动力学而不引入全局时间耦合。此外，该方法还提出运动感知点融合（MAPF）初始化方案：以第一帧的深度点云作为规范空间初始化，然后计算运动感知遮挡掩码识别具有显著运动的区域（像素颜色差异超过阈值~$\tau$~的区域），选择性融合其他帧的~3D~点填补器械遮挡造成的初始化空洞。

\textbf{性能。}在~DaVinci~机器人手术视频（EndoNeRF~和~StereoMIS~数据集）上，Deform3DGS~达到~PSNR 37.90~和~338.8 FPS~的渲染速度，训练时间仅需~1~分钟（对比~EndoGaussian~的~138~秒训练时间和~128 FPS~渲染速度），消融实验表明~FDM~相比~HexPlane~变形场（PSNR 37.13, 108~秒）和傅里叶/多项式基函数（PSNR 37.08, 60~秒）在保持相近训练效率的同时获得了更高的重建质量。

SurgicalGaussian（2024.07）~\cite{2407_05023} 针对手术场景的软组织重建提出前映射变形架构。其核心设计包括：(1) \textbf{深度引导初始化}（GIDM）——多帧点云累积初始化策略，在所有帧上累积深度和颜色信息（而非单帧~SfM），同时利用工具掩码排除器械区域，确保被器械遮挡的组织区域仍有~3--5$\times$~更多的初始点覆盖；(2) \textbf{前映射变形~MLP}——将~3D~高斯基元的规范状态通过时间条件的~MLP~映射到变形状态，建模软组织的时空特征；(3) \textbf{工具感知训练}——通过工具掩码引导训练过程，在渲染损失中排除器械区域，避免刚性器械对变形场优化的干扰。

SurgicalGS（2024.10）~\cite{2410_09292} 发现已有~3DGS~手术重建方法使用逆深度损失函数压缩深度变化，导致精细几何细节丢失。其引入归一化深度正则化损失（将深度归一化到~$[0,1]$~后再计算平滑约束，消除远近区域的梯度不均）和无监督深度平滑约束，在两个真实手术数据集上实现最先进的几何精度。

\textit{与~Deform3DGS~的差异。}Deform3DGS~的~FDM~虽能灵活建模动态变形，但其深度优化仅依赖标准逆深度损失~$L_{\text{depth}} = \|1/d - 1/\hat{d}\|$。逆深度变换~$1/d$~对远处深度值施加较大压缩（例如从~100 mm~到~101 mm~和从~1 mm~到~2 mm~的逆深度变化分别为~$10^{-4}$~和~0.5，相差三个量级），导致近处组织的精细几何细节在优化中被忽略。SurgicalGS~的核心改进是将逆深度损失替换为归一化深度正则化~$\mathcal{L}_{\text{NDR}} = \|\hat{D}/\max(\hat{D}) - D/\max(D)\|_1$，通过逐帧归一化消除绝对距离的梯度不均；同时引入无监督深度平滑约束~$\mathcal{L}_{\text{smooth}} = \|\nabla_x \hat{D}\|_1 + \|\nabla_y \hat{D}\|_1$~直接在深度图梯度域正则化几何。设计动机在于：手术导航对几何精度（尤其是近处组织）的要求远高于外观质量，而逆深度损失的压缩效应恰好破坏了最关键的近场精度。此外，SurgicalGS~采用二值运动掩码识别显著深度变化的像素区域，跨帧融合深度图点云进行初始化。这一初始化策略比~Deform3DGS~的~MAPF~更充分利用了深度先验。在~EndoNeRF~和~SCARED~数据集上，归一化深度正则化的消融实验表明其对几何精度的贡献最大。

EH-SurGS（2025.01）~\cite{2501_01101} 提出~3D~高斯生命周期感知的变形建模，同时捕捉常规和不可逆变形（如组织剪切），结合自适应运动层级策略区分静态和变形区域以提高渲染速度。\textit{与~Deform3DGS~的差异。}Deform3DGS~的~FDM~中每个高斯基元一旦创建后在整个时间序列中保持固定，无法表示组织剪切等不可逆拓扑变化。EH-SurGS~引入~3D~高斯生命周期感知机制：为每个高斯赋予出生时间~$t_{\text{birth}}$~和消亡时间~$t_{\text{death}}$~参数，当组织发生不可逆变形时，旧高斯通过消亡参数退出渲染，新高斯通过动态分裂机制产生。此外，EH-SurGS~提出自适应运动层级策略：根据运动幅度将高斯分为静态和动态两类，静态区域跳过变形网络直接渲染，将通过变形网络的高斯数量减少约~40\%。

Diff2DGS（2026.02）~\cite{2602_18314} 采用两阶段框架解决器械遮挡下的手术场景重建。\textit{与~Deform3DGS~的差异。}Deform3DGS~及后续方法在器械区域施加掩码跳过遮挡像素，但被遮挡组织始终缺乏监督信号。Diff2DGS~用扩散视频修复模块填补这一空白：第一阶段利用预训练扩散模型对被器械遮挡的组织区域进行时序一致的视频修复，生成伪真值组织外观；第二阶段以修复后的完整组织视频作为监督训练~2DGS~结合可学习变形模型捕获动态变形。选择~2DGS~的动机在于其天然编码表面法向信息。首次在~SCARED~数据集上执行定量深度精度分析，发现仅优化图像质量并不自动保证~3D~重建精度的提升。在~EndoNeRF~上达到~38.02 dB PSNR，在~StereoMIS~上达到~34.40 dB PSNR。

ColorGS（2025.08）~\cite{2508_18696} 针对手术场景中细微颜色变化和全局变形的挑战，引入着色高斯基元和增强变形模型（EDM），在~EndoNeRF~上达到~39.85 PSNR。\textit{与~Deform3DGS~的差异。}Deform3DGS~的颜色通过~SH~系数在规范空间中固定表示，时间变化仅通过几何变形间接影响渲染，无法捕获组织颜色的细微变化。ColorGS~引入两个核心改进：(1)~\textbf{着色高斯基元}——用携带可学习颜色参数的动态锚点替代固定~SH~系数，使每个高斯的外观能根据局部空间位置自适应调整；(2)~\textbf{增强变形模型}（EDM）——将~FDM~的单组~Gaussian~基函数扩展为时间感知~Gaussian~基函数与可学习时间无关变形的组合，前者捕获局部组织形变，后者建模器械-组织交互引起的一致性全局运动。在~EndoNeRF~上达到~39.85 PSNR（比~Deform3DGS~高~1.5 dB），SSIM~达~97.25\%。

LGS（2024.06）~\cite{2406_16073} 专注于轻量化~4DGS~框架，通过变形感知剪枝和~4D~特征场压缩实现~9$\times$~以上的压缩率。\textit{与~Deform3DGS~的差异。}Deform3DGS~虽实现~338.8 FPS~渲染速度，但存储开销对资源受限的手术设备仍构成瓶颈。LGS~从三个维度压缩~4DGS~存储冗余：(1)~\textbf{变形感知剪枝}——量化每个高斯基元对变形结果的贡献度，剪除贡献低于阈值的高斯；(2)~\textbf{属性维度剪枝}——对非关键区域的~SH~系数降维，在纹理简单的区域减少参数量；(3)~\textbf{4D~特征场凝缩}——对~HexPlane~多分辨率时空编码器的特征网格执行基于~SVD（Singular Value Decomposition，奇异值分解）~的低秩分解。在~EndoNeRF~上实现~9$\times$~以上压缩率，为便携式手术机器人的部署提供了可行性。

\subsection{器械建模与场景配准}
\label{ss:19_3}

本节介绍基于几何先验和关节运动学的手术器械建模与术中场景配准方法。

\subsubsection{手术器械建模与场景配准}

本路线介绍基于 CAD 几何先验和正向运动学的手术器械建模，以及术中场景配准方法。

手术器械的精确~3D~建模是手术仿真的关键环节——与软组织不同，器械是具有已知~CAD~模型的刚性物体，需要将高斯表示绑定到器械的已知几何结构上。

\textbf{Baseline 做法。}在标准~3DGS~中，场景由一组自由高斯基元表示，每个高斯拥有独立的位置~$\bm{\mu}$、旋转~$\bm{r}$、缩放~$\bm{s}$、不透明度~$\alpha$~和球谐系数~$sh$~等可学习属性。这种表示对软组织等非结构化场景非常有效，但无法直接建模具有已知几何结构和关节连接的刚性器械——因为标准~3DGS~的高斯基元是"自由漂浮"的，没有任何物理约束将它们锚定到器械的~CAD~几何上。此外，现有手术场景~3DGS~方法（如~Deform3DGS、EndoGaussian）用时间戳条件化变形场来驱动高斯基元的运动，这种时间驱动方式虽然能捕捉组织的非刚性变形，却无法让用户以关节角度为输入主动控制器械姿态，即无法实现~Real2Sim~的可控重建需求。

Instrument-Splatting（2025.03）~\cite{2503_04082} 正是针对上述两个瓶颈——高斯基元缺乏几何锚定和运动缺乏关节可控性——提出~Real2Sim~方法，修改了~3DGS~的高斯基元初始化和变换机制——用~CAD~几何预训练替代随机/结构光初始化，用正向运动学驱动替代自由优化——高斯的~$\alpha$-混合~渲染管线保持不变。

\textbf{几何预训练：从~CAD~到语义嵌入高斯。}该方法的第一步是将高斯点云精确绑定到器械各部件的~CAD~网格上。直觉上，这相当于给每个高斯基元"刷上"一层几何先验的外衣，使其分布贴合器械的真实形状而非自由飘散。具体实现为：对器械的每个刚性部件（shaft、wrist、gripper），在~CAD~网格表面通过光线追踪（ray-tracing）密集采样点作为高斯基元的初始位置——该采样策略有效避免了网格三角形过大导致的点分布稀疏问题。然后用~Blender~从多视角渲染~CAD~模型的~ground-truth~轮廓图，以~$\mathcal{L}_1$~损失优化高斯的不透明度渲染与~GT~轮廓之间的差异，使高斯点云精确学习部件的几何形状。预训练完成后，为每个高斯基元附加一个语义标签~$f \in \{1,2,3\}$（分别对应~shaft、wrist、gripper），形成"语义嵌入高斯"（semantics-embedded Gaussians）——即每个高斯的属性扩展为~$\Theta = \{\mu, r, s, \alpha, sh, f\}$。这种语义标签使后续的位姿估计和渲染能够区分不同部件，是该方法的核心设计之一。

\textbf{正向运动学：关节驱动高斯变换。}该方法的核心创新在于：用器械的正向运动学链替代标准~3DGS~中高斯基元的自由优化。以~da Vinci~的~EndoWrist~大针驱动器（LND）为例，其运动学结构包含三个关节变量~$q = \{\theta_1, \theta_2, \theta_3\}$，分别控制~shaft~绕~$z_1$~轴的旋转、wrist~绕~$y_2$~轴的旋转以及~gripper~的开合角度。正向运动学链定义为一系列刚体变换：

\begin{equation}
{}^{s}T_w(\theta_1)\;{}^{w}T_{g}(\theta_2)\;{}^{g}T_{l}(\theta_3),
\end{equation}

其中~${}^{a}T_b \in SE(3)$~表示从坐标系~$a$~到坐标系~$b$~的刚体变换，$s$、$w$、$g$、$l$~分别对应~shaft、wrist、gripper~和左工具尖端的坐标系。给定关节状态~$q$~和~shaft~到相机的变换~${}^{cam}T_s(\xi)$（其中~$\xi \in se(3)$~为位姿参数），任意部件~$j$~上的高斯基元从局部坐标系变换到相机坐标系的变换函数为：

\begin{equation}
\bm{T}_j(x_j;\;q,\xi) = {}^{cam}T_s(\xi)\;{}^{s}T_j(q)\;x_j,
\end{equation}

即先根据关节角度将局部点变换到~shaft~坐标系，再由~shaft~到相机的位姿变换到相机坐标系。这意味着：给定关节角度和~shaft~位姿，每个高斯基元的新位置和旋转可由刚体变换直接计算，而非通过梯度下降自由优化。在纹理学习阶段，位置~$\mu$~被冻结，仅优化~$\{r, s, \alpha, sh\}$~以学习外观——这保证了几何的精确性不会在学习纹理时退化。

\textbf{位姿估计与跟踪：render-and-compare~策略。}对于无位姿标注的真实手术视频，该方法设计了一种基于~render-and-compare~框架的位姿估计模块。核心思路是：用当前关节状态~$(q, \xi)$~驱动高斯渲染语义轮廓图，与~GT~分割掩码比较，通过最小化损失来优化位姿参数：

\begin{equation}
\hat{\xi},\;\hat{q} = \arg\min\;\mathcal{L}(\xi,\;q), \quad \mathcal{L} = \mathcal{L}_{mask} + \mathcal{L}_{tip},
\end{equation}

其中~$\mathcal{L}_{mask}$~为渲染语义轮廓与分割掩码的~$\mathcal{L}_1$~损失；$\mathcal{L}_{tip}$~是利用器械几何先验设计的工具尖端正则项，包含距离损失~$\mathcal{L}_{dist} = \text{ReLU}(d - r)$~（$d$~为渲染尖端与检测尖端的欧氏距离，$r$~为对齐阈值）和结构损失~$\mathcal{L}_{struct} = \text{ReLU}(-(\theta_l + \theta_r))$~（确保左尖端始终位于右尖端左侧），防止因错误的尖端检测陷入局部最优。对于帧间大运动，通过特征匹配~+~PnP（Perspective-n-Point）~提供粗初始位姿，再用~render-and-compare~精修——这是一种"粗到精"的位姿跟踪策略。

\textbf{设计动机。}标准~3DGS~和现有手术~3DGS~方法将高斯基元视为自由优化的~3D~点，既无法将高斯锚定到已知几何上，也无法以关节角度控制高斯的变换。这在软组织重建中不是问题（组织本身就是自由变形的），但对具有已知~CAD~模型和明确关节连接的器械来说是根本性的不匹配。Instrument-Splatting~通过"几何预训练~+~语义嵌入~+~正向运动学驱动"的三阶段设计解决了这一不匹配：预训练提供几何先验，语义标签支持部件级位姿估计，正向运动学确保关节可控性。在~EndoVis 2017/2018~数据集上，该方法在~wrist~\&~gripper~区域达到~25.81~dB~PSNR~和~0.064~LPIPS，显著优于~Deform3DGS~（19.65~dB）和~EndoGaussian（18.90~dB）；几何预训练贡献~+9.24~dB~PSNR（从~14.6~提升至~23.84~dB），$\mathcal{L}_{tip}$~使~gripper~的语义~Dice~从~63.17\%~提升至~73.32\%。

其升级版~Instrument-Splatting++（2026.03）~\cite{2603_22792} 进一步引入语义感知位姿估计（SAPET）和鲁棒纹理学习（RTL），从无位姿视频中学习真实纹理。\textit{与~Instrument-Splatting~的差异。}Instrument-Splatting~的~render-and-compare~位姿估计依赖人工标注的分割掩码提供监督信号——在真实手术室中获取逐帧掩码不现实。Instrument-Splatting++~的核心改进~SAPET~利用纯合成语义训练的~gripper-tip~网络替代人工掩码：该网络仅在~CAD~模型渲染的合成语义图上训练，通过检测真实视频中的工具尖端位置为位姿优化提供鲁棒监督——消除了对真实数据标注的依赖。此外，SAPET~引入松约束正则化（loose regularization），通过约束关节角度在物理合理范围内抑制``奇异关节''（如~gripper~过度张开）等非物理位姿解。RTL~模块解决~Instrument-Splatting~未涉及的核心难题：从无位姿视频中学习真实纹理。

由于位姿估计存在噪声，直接用噪声位姿优化纹理会导致模糊伪影。RTL~采用交替优化策略——每轮先精修~N~步位姿参数，再固定位姿优化~M~步纹理参数——逐步消除位姿噪声对纹理质量的影响。在~EndoVis17/18、SAR-RARP~和自建数据集上验证了方法的有效性，并首次展示了下游关键点检测任务中~unseen-pose~数据增强的性能提升。G-SHARP（2025.12）~\cite{2512_02482} 构建首个商业兼容的实时手术重建框架，标志着~3DGS~手术重建从学术概念验证向手术室实际部署的关键一步。在~Baseline~做法中，现有~3DGS~手术重建方法（如~EndoGS、Deform3DGS、LGS~等）虽在~EndoNeRF~基准上取得优异重建质量（PSNR~达~37.9~dB），但均依赖原始~3DGS~代码库的~fork——该代码库使用非商业许可证或包含第三方受限组件，每个方法各自创建独立分支，形成碎片化的~fork~生态，渲染路径互不兼容，无法在手术室等临床环境中合法部署。

G-SHARP~是首个原生构建于~GSplat（Apache-2.0~许可证）可微高斯光栅化器的手术重建管线——GSplat~是~Nerfstudio~项目维护的~PyTorch~兼容后端，提供融合前向/反向传播的高效~CUDA~内核和数值稳定的梯度计算。基于此后端，G-SHARP~实现三重创新：（1）面向~EndoNeRF-Pulling~的变形与掩码方案；（2）两阶段训练管线——粗阶段（~200~次迭代，约~30~秒）仅初始化几何，精阶段（~1500~次迭代，约~1.5~分钟）引入时间变形；（3）基于~Holoscan SDK~的实时推理部署。

变形网络由~HexPlane~时空特征网格和~MLP~解码器组成——HexPlane~将~4D~时空体分解为~6~个~2D~特征平面（3~个空间平面~XY/XZ/YZ + 3~个时空平面~XT/YT/ZT），空间分辨率~$64^3$、时间维度~100，采用多分辨率网格（尺度~$[1,2,4,8]$）捕获不同粒度特征；MLP~解码器（8~层，256~维隐藏单元）从双线性插值特征中预测每个高斯基元在时刻~$t$~的变形偏移~$\Delta\bm{\mu}$（位置）、$\Delta\bm{s}$（缩放）、$\Delta\bm{q}$（旋转四元数）和~$\Delta\alpha$（不透明度）。训练采用掩码监督策略——工具掩码应用于~GT~和渲染图像后再计算损失，确保梯度仅通过组织区域流动；对所有帧工具掩码取并集生成"不可见掩码"，对始终被遮挡区域施加~TV~正则化鼓励平滑插值。设计动机在于：手术场景的器械-组织遮挡和软组织非线性变形与自动驾驶中刚性环境根本不同，原始~3DGS~的紧耦合~C++/CUDA~内核虽针对固定~splat~公式优化，但难以扩展以支持工具感知遮挡推理和逐高斯变形场等手术特需功能。

在~EndoNeRF-Pulling~基准上达到~37.98~dB PSNR~和~$>$60 FPS~的实时渲染速度，训练仅需~2~分钟，并成功部署到~NVIDIA IGX Orin~和~Thor~边缘硬件。

场景配准方面，Registering Neural 4D Gaussians（2024.07）~\cite{2407_20213} 提出空间感知特征聚合方法（SWC），实现~4D~手术神经场景的动态配准。\textit{与~Instrument-Splatting~的差异。}Instrument-Splatting~聚焦于单个器械的~Real2Sim~可控重建，而场景配准关注的是将两个独立重建的~4D~手术场景在空间中对齐——这在术前规划与术中仿真中至关重要（例如将患者特异的~4D~解剖模型与标准手术流程的~4D~模板对齐）。SWC~方法首先用~4DGS~分别表示两个手术场景，然后在~4D~特征空间中执行空间感知的特征聚合——通过空间加权聚类（Spatially Weighted Clustering）对齐两个场景的特征分布，再利用可变形场景配准策略联合利用空间和时间信息进行对应匹配。此外，该方法提出变形场景配准的新范式：不同于传统刚性配准仅求解全局~SE(3)~变换，它估计两个动态场景之间的稠密变形场，使一个场景中的解剖结构能非线性变形以匹配另一个场景。

BridgeSplat（2025.09）~\cite{2509_18501} 实现术中~3D~重建与术前~CT~数据的双向耦合，将~3D~高斯绑定到~CT~网格上实现关节优化。\textit{与~Instrument-Splatting~的差异。}Instrument-Splatting~将高斯绑定到器械的~CAD~几何上，而~BridgeSplat~解决的是更通用的术前-术中数据融合问题：如何将术中~RGB~视频重建的~3D~高斯场景与术前~CT~体积关联起来，实现双向信息传播。其核心设计是将每个~3D~高斯基元参数化为相对于其父~CT~网格三角形的局部坐标~$(b_1, b_2, b_3)$，而非~Instrument-Splatting~的全局~SE(3)~关节变换。这一参数化确保高斯始终贴合~CT~网格表面，变形可通过雅可比链传播回~CT~网格更新术前数据。通过光度监督（术中~RGB~图像的渲染损失）同时优化高斯参数和网格变形，实现了从单目~RGB~数据驱动~CT~体积变形的管线。在猪内脏手术和合成人肝数据上验证了方法的有效性。

SurgCalib（2026.03）~\cite{2603_08983} 利用~3DGS~可微渲染管线实现~da Vinci~手术机器人的手眼标定。\textit{与~Instrument-Splatting~的差异。}Instrument-Splatting~估计的是器械的关节角度和~shaft~位姿，而~SurgCalib~解决的是更基础的机器人学问题：估计机器人基座到内窥镜相机之间的刚性变换~$AX = XB$~。传统手眼标定需要~fiducial markers（标记板），违反手术室无菌协议。SurgCalib~提出无标记方案：首先用原始运动学测量初始化器械位姿，然后在~3DGS~可微渲染管线中通过两阶段优化精修位姿——第一阶段在~RCM（远程中心运动）约束下联合优化手眼变换和器械位姿，第二阶段用~render-and-compare~策略精修。在~SurgPose~基准上，左右器械的工具尖端~2D~重投影误差分别达到~12.24 px~（2.06 mm）和~11.33 px~（1.9 mm），3D~距离误差分别为~5.98 mm~和~4.75 mm。

数据生成方面，Zeng et al.（2024.07）~\cite{2407_14846} 首次探索基于~3DGS~的合成手术数据集生成。\textit{与~Instrument-Splatting~的差异。}Instrument-Splatting~追求高保真重建，Zeng et al.~则从数据稀缺角度出发：手术~CV~模型（如工具检测、分割）需要大量标注数据，但真实手术数据获取昂贵且标注困难。该方法提出从~3DGS~重建的手术场景中提取器械和背景的高斯表示，然后通过变换和重组生成新的合成手术图像。具体流程为：分别重建器械和背景环境的~3DGS~模型，然后将器械高斯变换到新位姿（平移、旋转、关节角度变化）并嵌入背景场景渲染合成图像。开发了数据采集系统同时获取图像和器械/相机位姿，使合成图像与真实图像可进行直接对比（PSNR 29.592）。在~YOLOv5~工具检测任务中，合成数据训练的模型在真实测试集上比真实数据训练的模型性能提升~12\%。

NeeCo（2025.08）~\cite{2508_07897} 在~Zeng et al.~基础上提出动态高斯模型，支持渲染未见过视角和变形状态的器械。\textit{与~Zeng et al.~的差异。}Zeng et al.~仅处理静态器械的位姿变换，无法建模器械在操作过程中的关节运动（如~gripper~的开合）。NeeCo~引入动态训练调整策略：在训练过程中动态调整高斯基元的位置和属性以适应相机位姿不准带来的挑战——真实场景中的位姿估计不可避免地存在误差，静态高斯模型在位姿偏差处会产生伪影。NeeCo~通过联合优化高斯属性和位姿修正量来缓解这一问题。此外，NeeCo~提出基于动态高斯的自动标注方法：利用高斯基元的~3D~位置和语义属性自动生成分割掩码和关键点标注，消除了人工标注需求。在包含~7~个场景、14000~帧的数据集上验证，合成图像~PSNR~达~29.87，训练于合成数据的~YOLOv5~模型比标准数据增强提升~10\%，总体性能提升近~15\%。

\begin{takeaway}
内窥镜与手术场景重建将 3DGS 引入微创手术：器械感知的联合优化实现手术工具的高保真重建，术前 CT 与术中 RGB 的双向耦合为手术导航提供精确配准，合成数据生成则有效缓解了手术 AI 的数据稀缺瓶颈。
\end{takeaway}

%=====================================================================

\subsection{近期方法脉络：Endoscopy novel view synthesis benchmarks}

这一组近期工作延续 3DGS，主要把改动集中在评测等环节。下面只保留相对既有路线的变化，便于和前文方法直接对照。

\paragraph{Gastroendoscopy View Synthesis}
Gastroendoscopy View Synthesis（2026.06）~\cite{2606_25427} 延续 3DGS 在“Endoscopy novel view synthesis benchmarks”上的基本设定，主要改变是：构建 GastroNVS 真实数据集并比较 3DGS、NeRF 等 NVS 方法。3DGS 的显式高斯表示和可微渲染能力保持为基础。
其设计动机是：提供真实胃镜图像、相机位姿和点云，用于评估 NVS 在胃镜数字孪生和训练中的潜力。
具体而言，论文不是提出新的 3DGS 表示，而是把 3DGS-based methods 作为对比对象，在胃镜稀疏、非刚性和低纹理环境下测试其效果。
适用边界是：侧重数据集和评测，方法贡献较少；医学内窥镜中的组织形变和相机标定误差仍可能限制 3DGS 表现。
\section{CT/MRI断层成像重建}
\label{sec:19_2}

将~3DGS~从光学渲染扩展到断层成像重建，是~3DGS~在医学影像中最具方法论创新的方向之一。核心挑战在于：3DGS~的~$\alpha$-混合~光栅化假设不适用于~X~射线投影的积分物理和~MRI~的~k~空间重建。X~射线投影的像素值是密度沿射线路径的线积分；MRI~需要处理复数域的磁共振信号和~k~空间的欠采样模式。两者都需要重新推导适合各自成像物理的渲染方程。

\subsection{X射线层析重建}
\label{ss:19_4}

本节讨论三条技术路线：（1）辐射高斯与层析重建；（2）4D CT与DSA重建；（3）专用层析成像。

\subsubsection{辐射高斯与层析重建}

本路线介绍基于辐射高斯表示的静态 X 射线层析重建方法。

R$^2$-Gaussian（2024.05）~\cite{2405_20693} 是首个将~3DGS~用于稀疏视角层析重建的框架。其核心创新在于提出了"辐射高斯"（Radiative Gaussian）表示——与标准~3DGS~的高斯包含视角相关颜色（SH~系数）不同，辐射高斯移除了视角依赖的颜色分量，因为~X~射线衰减仅取决于各向同性密度。每个辐射高斯基元定义一个局部~Gaussian~形密度场：$G_i^3(x | \rho_i, p_i, \Sigma_i) = \rho_i \exp(-\frac{1}{2}(x - p_i)^T \Sigma_i^{-1}(x - p_i))$，其中~$\rho_i$~为中心密度、$p_i$~为位置、$\Sigma_i = R_i S_i S_i^T R_i^T$~为协方差。整体密度通过求和计算：$\mu(x) = \sum_{i=1}^{M} G_i^3(x)$。

在~X~射线光栅化推导中，R$^2$-Gaussian~发现了一个关键性质：归一化~3D~高斯分布沿一个坐标轴的积分直接得到归一化~2D~高斯分布。因此，X~射线投影可以通过简单求和~2D~高斯来渲染（而非~alpha-compositing）：$I_r(r) = \sum_{i=1}^{M} G_i^2(\tilde{x} | \tilde{\rho}_i, \tilde{p}_i, \tilde{\Sigma}_i)$，其中~$\tilde{\rho}_i = \rho_i \cdot (2\pi)^{3/2}|\Sigma_i|^{1/2}$~是投影后的~2D~高斯中心密度。更关键的是，R$^2$-Gaussian~发现了标准~3DGS~中未知的"积分偏差"（integration bias）——3DGS~实际上学习的是~2D~图像平面上的积分密度而非~3D~空间中的实际密度，这导致不同视角恢复出不同的~3D~密度值，违反了~X~射线衰减的各向同性本质。R$^2$-Gaussian~通过将实际~3D~密度分配给高斯基元并正向计算~2D~投影，从根本上解决了这一问题。此外，该方法开发了可微体素化器（voxelizer）用于~3D~密度体积查询和全变分（TV）正则化，初始化采用~FDK~解析重建的快速低质量体积，在~4~分钟内完成高质量重建，比~NeRF~方法快~12$\times$。

在此基础上，Radioactive 3D Gaussian Ray Tracing（2026.02）~\cite{2602_01057} 用光线追踪替代~splatting~，通过解析计算~3D~高斯基元的线积分避免局部仿射近似，提供更物理一致的正向投影模型。\textit{与~R$^2$-Gaussian~的差异。}R$^2$-Gaussian~采用~EWA splatting~策略：将~3D~高斯通过局部仿射近似坍缩为~2D~高斯再叠加渲染。这一近似在投影角度变化平缓处足够精确，但当高斯基元距离探测器较远或投影角度极端时，局部仿射近似引入系统性误差。Radioactive 3D Gaussian Ray Tracing~的核心改进是沿射线路径解析计算~3D~高斯的线积分——对射线上每个采样点直接评估~3D~高斯密度函数并沿射线方向积分，避免了~3D~到~2D~的仿射坍缩。这一改进带来两个优势：(i)~投影模型更加物理一致，定量精度更高；(ii)~光线追踪公式化提供对射线起点和方向的显式控制，便于集成非线性几何校正——例如~PET~成像中的弧形校正（arc-correction），将~splatting~方法难以处理的探测器几何直接纳入正向投影模型。

FaCT-GS（2026.04）~\cite{2604_01844} 通过优化体素化和光栅化管线，在~512$\times$512~投影上比已有~GS-CT~方法快~4$\times$~以上。\textit{与~R$^2$-Gaussian~的差异。}R$^2$-Gaussian~的可微体素化器将高斯场转换为离散体素网格时，采用逐体素评估所有高斯基元贡献的策略——计算复杂度为~$O(N \cdot V)$（$N$~为高斯数量，$V$~为体素数量），在大体积或高分辨率投影场景中成为瓶颈。FaCT-GS~的核心贡献是对体素化和光栅化管线的深度优化：(1)~将体素化从逐体素评估改为基于~3D~包围盒的快速排序策略，仅计算对当前体素有贡献的高斯；(2)~优化光栅化中的~2D~投影排序和~$\alpha$-混合~内核。这两项优化使~FaCT-GS~在~512$\times$512~投影上比~SOTA~快~4$\times$，在~2K~投影上快~13$\times$。此外，改进的体素化使~FaCT-GS~能快速将已有~CT~体积拟合为高斯表示——可作为重建的热启动先验，或作为压缩表示替代原始体积。

针对稀疏视角下的伪影问题，GR-Gaussian（2025.08）~\cite{2508_02408} 引入基于图的梯度策略抑制针状伪影，PSNR~提升~0.67--0.92 dB。\textit{与~R$^2$-Gaussian~的差异。}R$^2$-Gaussian~在致密化过程中使用视角内所有点的平均梯度幅值决定高斯的分裂——在稀疏视角条件下，这种平均化策略无法区分真正需要增加表示能力的区域和已被充分约束的区域，导致高斯沿射线方向过度延伸产生针状伪影。GR-Gaussian~提出两个核心改进：(1)~\textbf{去噪点云初始化}——对~FDK~初始重建执行结构感知的去噪处理，减少初始点云中的伪结构点，加速后续收敛；(2)~\textbf{像素图感知梯度策略}——构建像素邻域图，利用图中相邻像素间的密度差异替代简单的平均梯度幅值——图结构捕获了密度场的局部不连续性，使致密化决策能精确聚焦于重建不充分的区域而非已充分约束的区域。在~X-3D~和真实数据集上，PSNR~分别提升~0.67 dB~和~0.92 dB，SSIM~分别提升~0.011~和~0.021。

Residual Gaussian Splatting（2026.04）~\cite{2604_27552} 提出频谱解耦高斯表示，将体积场分层为几何基分量和残差细节分量，解决~X~射线衰减的非负性与小波高频系数的双极性之间的数学矛盾。\textit{与~R$^2$-Gaussian~的差异。}R$^2$-Gaussian~用单一密度值~$\rho_i$~表示每个高斯基元的衰减，通过直接优化~$\rho_i$~拟合投影数据。然而在超稀疏视角条件下，标准光度优化存在频谱偏差（spectral bias）——网络倾向于优先拟合低频结构而忽略高频细节，导致过度平滑。Residual Gaussian Splatting~的核心创新是频谱解耦：将体积场分解为几何基分量（捕获大尺度低频结构）和残差细节分量（捕获小尺度高频纹理），两者由独立的高斯组分别表示。设计动机在于解决一个数学矛盾：X~射线衰减系数必须非负（物理约束），但小波变换的高频系数天然具有双极性（正负值表示边缘的明暗跳变）。直接用非负密度拟合高频系数在数学上不可能。Residual GS~通过频谱-空间协同优化策略协调几何锚定与纹理精修，有效防止频谱串扰，在骨小梁和血管结构等复杂解剖细节上实现卓越的视觉保真度。

TG-Field（2026.02）~\cite{2602_11705} 引入几何感知的高斯变形框架，用多分辨率哈希编码捕获局部空间先验，支持静态和动态~CT~重建。\textit{与~R$^2$-Gaussian~的差异。}R$^2$-Gaussian~仅适用于静态~CT~重建，且在超稀疏视角下缺乏空间正则化导致高斯基元自由漂移。TG-Field~的核心改进是用多分辨率哈希编码为高斯基元提供局部空间先验——在超稀疏视角下，大量高斯缺乏充分的多角度约束，哈希编码通过空间邻域信息正则化基元参数，抑制欠约束区域的自由漂移。此外，TG-Field~将框架扩展到动态~CT~重建：引入时间条件化表示和时空注意力块自适应聚合时空特征，解决时空模糊并强制时间一致性；运动流网络（motion-flow network）建模精细呼吸运动以跟踪局部解剖变形。在~AAAI 2026~发表，在合成和真实数据集上均超越现有方法。

\subsubsection{4D CT与DSA重建}

本路线介绍处理呼吸运动的时间维度变化和 DSA 血管动态重建的方法。

4D CT~重建需要处理呼吸运动带来的时间维变化。

\textbf{Baseline 做法。}传统~4D CT~重建遵循相位分箱（phase-binning）工作流：首先借助外部呼吸门控设备（需直接接触患者）将~X~射线投影按呼吸周期划分为~10~个离散相位，然后对每个相位独立执行~3D~重建（如~FDK~算法或~TV~最小化），得到一组~3D~体积序列。这种离散化策略存在两个根本缺陷：(1)~每个相位仅分到约~$1/10$~的投影数据，重建中产生严重的条纹伪影；(2)~离散相位无法建模解剖结构在呼吸周期中的连续时空演化，相位间插值引入时间不一致性。近期基于~3DGS~的静态~CT~重建方法（如~R$^{2}$-Gaussian）虽已实现高质量的~3D~体积恢复，但无法处理呼吸运动的时间维度变化；而并发工作仅将时间离散化为~10~个时间戳对应~10~个相位，本质上仍是离散表示。

X$^{2}$-Gaussian（2025.03）~\cite{2503_21779} 首次实现真正的连续时间~4D CT~重建——在任意时刻查询动态~CT~体积，彻底绕过相位分箱。其核心是将静态辐射高斯~Splatting~扩展到时间域，并引入自监督呼吸周期学习消除对外部门控设备的依赖。

\textbf{动态高斯运动建模。}直觉上，该方法为每个~3D~高斯基元赋予一个时间变形场：给定任意时间戳~$t$，变形场预测该高斯的位置偏移、旋转偏移和缩放偏移，使其从规范状态变换到~$t$~时刻的变形状态。形式化地，变形后的高斯定义为~$G_i' = G_i + \Delta G_i = (\bm{\mu}_i + \Delta\bm{\mu}_i,\; \bm{R}_i + \Delta\bm{R}_i,\; \bm{S}_i + \Delta\bm{S}_i,\; \rho_i)$，其中~$\Delta\bm{\mu}_i$、$\Delta\bm{R}_i$、$\Delta\bm{S}_i$~分别为位置、旋转和缩放的变形偏移量。值得注意的是，中心密度~$\rho_i$~保持时间不变——因为~CT~密度由组织本身决定，不随呼吸改变。变形场~$\mathcal{D}$~由两部分组成：$\mathcal{D} = \mathcal{F} \circ \mathcal{E}$，其中~$\mathcal{E}$~为时空编码器，$\mathcal{F}$~为变形感知解码器。

在编码端，为避免直接用神经网络参数化时空特征导致的过拟合和低渲染速度，X$^{2}$-Gaussian~借鉴动态场景重建中的~K-Planes~分解策略，将~4D~特征空间因式分解为六组多分辨率特征平面：三个空间平面~$\mathcal{P}_{xy}$、$\mathcal{P}_{xz}$、$\mathcal{P}_{yz}$~捕获局部解剖关系，三个时空平面~$\mathcal{P}_{xt}$、$\mathcal{P}_{yt}$、$\mathcal{P}_{zt}$~捕获全局运动模式。给定高斯中心~$\bm{\mu}=(x,y,z)$~和时间戳~$t$，将~4D~坐标~$\bm{v}=(x,y,z,t)$~投影到各平面上，通过双线性插值提取特征，再经~Hadamard~积拼接后通过单层~MLP~融合，得到时空特征~$\bm{f}_h$。

在解码端，轻量级多头解码器从~$\bm{f}_h$~分别预测三类变形参数：$\Delta\bm{\mu} = \mathcal{F}_\mu(\bm{f}_h)$，$\Delta\bm{R} = \mathcal{F}_R(\bm{f}_h)$，$\Delta\bm{S} = \mathcal{F}_S(\bm{f}_h)$。这种解耦设计使平移运动、方向变化和体积膨胀/收缩可独立学习，同时支持时间轴上的连续插值，实现平滑的时间过渡。

\textbf{自监督呼吸周期学习。}呼吸运动的准周期性为重建提供了强大的生理先验：解剖结构在每个呼吸周期后回到近似相同的位置，即~$t$~时刻的解剖状态应与~$t + nT$~时刻一致（$T$~为呼吸周期，$n$~为整数）。基于此，X$^{2}$-Gaussian~将呼吸周期~$T$~建模为可学习参数~$\hat{T} = \exp(\hat{\tau})$（对数空间参数化保证正值和数值稳定），通过生理驱动周期一致性损失从投影数据中自监督学习：

\begin{equation}
\mathcal{L}_{pc} = \mathcal{L}_1(I(t),\; I(t + n\hat{T})) + \lambda_1 \mathcal{L}_{ssim}(I(t),\; I(t + n\hat{T})),
\end{equation}

其中~$n \in \{-1, 1\}$（限制为相邻周期以避免学到谐波的周期估计）。最优周期通过~$\tau^* = \arg\min_{\hat{\tau}} \mathcal{L}_{pc}$~获得，$T^* = \exp(\tau^*)$。论文展示了两种关键数值技巧的必要性：无界循环偏移导致~$\hat{T}$~趋向真值的子谐波（如~$T/2$）；线性空间参数化导致优化轨迹剧烈振荡——两者结合后~$\hat{T}$~稳定收敛到正确呼吸周期。

训练采用渐进式策略：先用~5000~次迭代训练静态~3D~辐射高斯~Splatting~捕获基础解剖结构，再扩展为完整~4D~框架联合优化高斯参数、时空编解码器和呼吸周期参数。总损失函数为~$\mathcal{L}_{total} = \mathcal{L}_{render} + \alpha \mathcal{L}_{pc} + \beta \mathcal{L}_{TV}^{3D} + \gamma \mathcal{L}_{TV}^{4D}$，其中~$\mathcal{L}_{render} = \mathcal{L}_1 + \lambda_2 \mathcal{L}_{ssim}$~为渲染损失，$\mathcal{L}_{TV}^{3D}$~为~3D~全变分正则化，$\mathcal{L}_{TV}^{4D}$~为多分辨率~K-Plane~网格正则化。在三个公开~4D CT~数据集（DIR、4DLung、SPARE）的~13~例患者上达到~SOTA，比传统方法提升~9.93 dB PSNR，比先前高斯~Splatting~方法提升~2.25 dB。

DIGS（2025.06）~\cite{2506_22280} 将自由变形（FFD）模型与~4DGS~结合用于动态~CBCT（Cone Beam CT，锥束CT）~重建，比~HexPlane~快~6$\times$。\textit{与~X$^{2}$-Gaussian~的差异。}X$^{2}$-Gaussian~采用~K-Planes~特征分解作为时空编码器——将~4D~时空体因式分解为~6~个~2D~特征平面，通过双线性插值提取时空特征。这种隐式表示虽然灵活，但计算开销大且缺乏空间正则化：高斯基元的运动在相邻空间位置可能不一致。DIGS~的核心差异是引入基于~FFD（Free-Form Deformation）的空间基函数作为变形模型——FFD~将变形场定义在均匀~3D~控制点网格上，每个控制点的位移通过~B-spline~插值传播到整个空间，天然提供平滑的空间正则化。此外，DIGS~提出变形知情框架（deformation-informed framework）：将高斯基元的位置、缩放和旋转的时变演化统一在同一个~FFD~变形场下约束，而非~X$^{2}$-Gaussian~中三组独立解码器分别预测~$\Delta\bm{\mu}$、$\Delta\bm{R}$、$\Delta\bm{S}$~。这种耦合设计确保了几何参数变化的空间一致性，同时比~HexPlane~快~6$\times$。在~6~个~CBCT~数据集上验证了方法的优越性（MICCAI 2025）。

DGR（2024.11）~\cite{2411_04844} 提出离散高斯表示用于端到端的体积~CT~重建，引入高度并行化的快速体积重建技术。\textit{与~R$^{2}$-Gaussian~的差异。}R$^{2}$-Gaussian~采用~splatting~渲染：将~3D~高斯投影到~2D~探测器平面后叠加渲染。DGR~采用离散高斯表示——将~3D~高斯直接参数化为~CT~体积中的密度分布，通过光线追踪计算沿射线路径的积分。DGR~的核心技术贡献是快速体积重建（Fast Volume Reconstruction）——一种高度并行化的技术，将高斯基元的贡献高效聚合到体素网格上，最小化计算开销。相比~R$^{2}$-Gaussian~的~splatting~路线，DGR~的离散化方案天然避免了~EWA splatting~中的局部仿射近似误差。在真实和合成数据集上均展现了优越的重建质量和运行时性能（ICCV 2025）。

DSA~血管重建方面，TOGS（2024.03）~\cite{2403_19586} 引入时间不透明度偏移表建模造影剂的时间辐射变化，实现~4D DSA~的实时渲染。\textit{与~R$^{2}$-Gaussian~的差异。}R$^{2}$-Gaussian~的辐射高斯基元拥有固定密度~$\rho_i$，适用于静态~CT~重建。DSA~中造影剂在血管中随时间流动，导致同一位置的~X~射线衰减在不同时间步发生显著变化。TOGS~的核心创新是为每个高斯基元引入时间不透明度偏移表（temporal opacity offset table）：为每个高斯维护一个离散的偏移值表~$\{\Delta\alpha_t\}$，在任意时间步~$t$~通过线性插值查询该时刻的不透明度偏移，将静态不透明度~$\alpha$~修正为~$\alpha + \Delta\alpha_t$。这一设计将时间动态建模从变形场（~X$^{2}$-Gaussian~的路线）简化为逐高斯的不透明度调制——血管结构本身是静态的，仅衰减值随时间变化，因此无需建模空间变形。此外，TOGS~引入平滑损失缓解稀疏视角下的过拟合，训练阶段通过随机剪枝高斯减少存储开销。4DRGS（2024.12）~\cite{2412_12919} 针对稀疏视角动态~DSA~血管重建提出~4D~辐射高斯~Splatting~框架。

在~Baseline~做法中，现有~DSA~重建方法（如~VPAL）基于~NeRF~建模整个~3D~空间的密度场，需要逐点体渲染，单病例处理耗时数小时；而~R$^2$-Gaussian~等~3DGS~层析方法仅适用于静态~CT~重建，无法处理~DSA~中造影剂随时间流动导致的衰减变化。4DRGS~的核心改进是将~3D~辐射高斯扩展为~4D~辐射高斯核：每个核的几何参数（位置、旋转、缩放）保持时间不变，用于建模静态血管结构；而中心衰减~$\rho(\bm{p}_i, t)$~为时间相关，由紧凑神经网络预测以捕获造影剂流动引起的时间变化衰减响应。4D~核的响应函数定义为~$G_i(\bm{x}, t) = \rho(\bm{p}_i, t) \cdot \exp(-\frac{1}{2}(\bm{x} - \bm{p}_i)^\top \Sigma_i^{-1}(\bm{x} - \bm{p}_i))$。

在技术实现上，4DRGS~提出三个关键组件：（1）动态神经衰减场（DNAF）——采用~3D~哈希编码捕获时间不变的空间特征，4D~哈希编码捕获时空变化特征，两者拼接后经~MLP~解码输出衰减值；（2）有界缩放激活——将标准~3DGS~的指数激活替换为~$\bm{s}_i = (s_{max} - s_{min})\sigma(\tilde{\bm{s}}_i) + s_{min}$，将核尺度约束在范围内，避免生成过长的高斯核产生针状伪影；（3）累积衰减剪枝——统计每个高斯核在所有视角和时间步上的累积衰减贡献，剪除贡献低于阈值的非血管核。设计动机在于：血管在扫描过程中保持静态结构，仅衰减随造影剂流动而变化——将几何与衰减解耦避免了~4D~变形场的高计算开销；基于核的表示天然跳过空背景区域，相比~NeRF~的全空间建模效率更高。训练仅需~5~分钟达到~SOTA，比~VPAL~快~32$\times$。DSA-SRGS（2026.03）~\cite{2603_04770} 首次提出超分辨率~3DGS~框架用于动态稀疏视角~DSA~重建。\textit{与~4DRGS~的差异。}4DRGS~通过~DNAF~建模造影剂的时空变化衰减，但其重建分辨率受限于输入投影的分辨率——对低分辨率投影执行朴素上采样会导致严重模糊和锯齿伪影。DSA-SRGS~的核心创新是引入多保真度纹理学习模块（Multi-Fidelity Texture Learning Module）：从微调的~DSA~专用超分辨率模型提取高质量先验，将其融入~4D~重建优化过程。为缓解伪标签可能引入的幻觉伪影，该模块采用置信度感知策略（Confidence-Aware Strategy）自适应加权原始低分辨率投影和生成的高分辨率伪标签之间的监督信号。此外，DSA-SRGS~提出辐射亚像素致密化（Radiative Sub-Pixel Densification）：利用高分辨率亚像素采样累积的梯度信息细化~4D~辐射高斯核，使高斯分布能更精确地匹配亚像素级血管细节。在两个临床~DSA~数据集上显著超越~SOTA。

\subsubsection{专用层析成像}

本路线介绍各向异性 X 射线交互建模、CT 新视角合成、非破坏性检测等专用层析成像方法。

除标准~CT~重建外，3DGS~还被适配到多种专用层析成像场景，包括各向异性~X~射线交互建模、CT~新视角合成、非破坏性检测和超分辨率重建等。本节介绍~8~种代表性方法。

\textbf{Baseline 做法。}数字重建射线图（DRR，Digitally Reconstructed Radiograph）是从~3D CT~体积模拟生成的~2D X~射线图像，广泛用于术前规划。现有~3DGS~用于~DRR~渲染的方法存在两类局限：一方面，物理精确的~Monte Carlo~模拟虽然能正确建模~Compton~散射等各向异性~X~射线交互，但计算量极大（需~$>10^8$~次迭代渲染一张图像），无法用于实时应用；另一方面，解析式~DRR~渲染器和早期~3DGS~方法（如~X-Gaussian、GaSpCT）采用光线追踪沿射线累积~HU~值，虽然高效但假设~X~射线衰减是各向同性的，忽略了散射等方向相关的物理效应。此外，这些方法的高斯点云初始化采用均匀间距或随机采样，忽略了~CT~数据的几何和材料属性，导致收敛慢、精度低。

DDGS-CT（2024.06）~\cite{2406_02518} 提出方向解耦~3DGS（DDGS），在高效可微~DRR~生成中首次近似各向异性~X~射线交互（如~Compton~散射），同时适配~CT~数据特性的初始化策略。其核心思想是将~X~射线与物质的交互分解为各向同性和各向异性两个独立分量，用独立的~3D~高斯组分别建模。

\textbf{方向解耦辐射度函数。}直觉上，~X~射线成像中的光子-物质交互同时包含各向同性现象（如光电吸收、Rayleigh~散射）和各向异性现象（如~Compton~散射）。标准~3DGS~的球谐（SH）分解虽然也包含各向同性和各向异性分量，但假设它们共处同一个~3D~高斯基元。DDGS~的关键创新在于打破这种共位约束：用两组独立的高斯~$g_i^{\text{iso}}$~和~$g_j^{\text{dir}}$~分别建模两个分量。它们的辐射度函数分别定义为：

\begin{equation}
c_i^{\text{iso}} = \text{sigmoid}(\bm{b}^{\text{iso}} \cdot \bm{f}_i^{\text{iso}}); \quad c_j^{\text{dir}} = \text{sigmoid}(Y_{1..L}(\theta,\phi) \cdot \bm{B}^{\text{dir}} \bm{f}_j^{\text{dir}}),
\end{equation}

其中~$\bm{f}_i^{\text{iso}}, \bm{f}_j^{\text{dir}} \in \mathbb{R}^k$~为特征向量，$\bm{b}^{\text{iso}} \in \mathbb{R}^k$~为全局各向同性基向量（所有~splat~共享），$\bm{B}^{\text{dir}} \in \mathbb{R}^{k_L \times k}$~为各向异性基矩阵，$Y_{1..L}(\theta,\phi)$~为去掉零阶常数项的~$L$~阶球谐分解（$k_L = L(L+2)$），$\theta,\phi$~为射线方向角。

设计动机在于：各向同性分量捕获~X~射线衰减的主体（光电效应），仅依赖材料本身；各向异性分量近似散射效应，依赖射线入射方向。两者解耦后，当训练~DRR~无散射（scatter-free）时，方向相关分量可优雅退化为轻量形式（类似~X-Gaussian~的单基向量），避免不必要的计算开销。实验表明这种解耦在提升图像精度的同时降低了表示复杂度。

\textbf{放射密度感知双采样（RADS）。}标准~3DGS~使用~SfM~点云初始化（仅采样场景表面），或均匀/随机采样~3D~空间——这些策略对~CT~体积数据均不适用，因为~CT~的体素包含~HU~放射密度信息而非表面几何。DDGS~提出放射密度感知双采样（Radiodensity-Aware Dual Sampling, RADS）：基于~CT~体积的~HU~值，以两种策略分别采样~3D~点集~$\mathbf{P}^{\text{mc}}$~和~$\mathbf{P}^{\text{dw}}$~——前者针对不同解剖元素的接触区域（此处光子传输可能产生不连续性），后者针对均匀区域（主导衰减可稀疏建模但受益于方向相关噪声建模）。这种初始化策略同时服务于双高斯模型：各向同性高斯从接触区域的密度梯度采样，各向异性高斯从均匀区域采样以建模散射效应。相比均匀采样，RADS~显著加速~3DGS~收敛并提升运行时性能。

在多个医学基准数据集上，DDGS~的图像精度超越~SOTA~方法，同时在位姿配准等逆问题中展现出实时术中应用的潜力——比解析式~DRR~方法提供更高的配准精度和运行时性能。

GaSpCT（2024.04）~\cite{2404_03126} 首次将~3DGS~用于~CT~新视角合成，用~$\beta$~损失和~TV~正则化替代~SfM~初始化。\textit{与~R$^{2}$-Gaussian~的差异。}R$^{2}$-Gaussian~提出了辐射高斯的完整数学框架和积分偏差修正，但~GaSpCT~作为~3DGS~在~CT~领域的早期探索，贡献在于适配策略而非数学推导。GaSpCT~的两个核心差异是：(1)~用~$\beta$~损失替代标准~$L_2$~损失——~$\beta$~损失~$\sum |x|^\beta$（$\beta < 1$）对稀疏信号施加更强惩罚，鼓励高斯分布集中于高密度区域（骨骼、组织边界），实现背景-前景的强区分；(2)~用均匀先验替代~SfM~初始化——将高斯中心初始化在脑部~CT~预期的~3D~空间范围内，而非~SfM~点云（~CT~投影无需~SfM~）。在~PPMI~脑部~CT~数据集上，GaSpCT~验证了~3DGS~在~CT~新视角合成中的可行性，内存需求比等效体素网格降低~17\%。

X-GRM（2025.05）~\cite{2505_15235} 提出大规模前馈模型从稀疏视角~X~射线投影重建~CT~体积，引入体素化高斯~Splatting（VoxGS）表示。\textit{与~R$^{2}$-Gaussian~的差异。}R$^{2}$-Gaussian~是逐案例优化方法——对每例患者独立训练高斯参数，无法泛化到新患者。X-GRM~是前馈模型：用可扩展~Transformer~编码稀疏视角~X~射线输入，跨视角~token~高效融合后解码为~VoxGS~表示。VoxGS~将每个体素位置关联一组高斯参数（位置偏移、缩放、密度），实现高效~CT~体积提取和可微~X~射线渲染。高容量~Transformer~架构和灵活的~VoxGS~表示使~X-GRM~能从域内和域外~X~射线投影产生高质量重建，展现了~3DGS~在~CT~重建中的泛化潜力。

IXGS（2025.04）~\cite{2504_14699} 将~3DGS~扩展到术中稀疏任意姿态~X~射线的~3D~重建。\textit{与~R$^{2}$-Gaussian~的差异。}R$^{2}$-Gaussian~假设~X~射线投影的几何参数（视角、姿态）是精确已知的。术中场景（如脊柱手术中的~C~臂~X~射线）面临两个额外挑战：投影姿态可能是任意且不规则的，不同视角间的视觉一致性差（不同曝光、不同组织重叠）。IXGS~在~R$^{2}$-Gaussian~框架上引入解剖引导的射线标准 化（anatomy-guided radiographic standardization）：使用风格迁移统一不同视角间的视觉外观，提升跨视角一致性从而增强重建质量。无需预训练，天生适应新患者和解剖结构。在~ex-vivo~数据集上验证，外科专家评估确认~20--30~视角的~3D~重建具有临床导航价值。

LamiGauss（2025.09）~\cite{2509_13863} 针对非破坏性检测中的~X~射线层析成像提出专用方法。\textit{与~R$^{2}$-Gaussian~的差异。}R$^{2}$-Gaussian~面向标准~CT~几何（旋转台式扫描，射线源绕物体旋转）。层析成像（Laminography）面向平板状结构（微芯片、复合材料），射线源和探测器以倾斜角照射——这一几何差异使标准~CT~的~splatting~公式不再适用。LamiGauss~引入专用检测器-世界坐标变换模型，将层析倾斜角纳入投影计算。此外，LamiGauss~提出层析伪影感知的初始化策略：对~FDK~初始重建显式滤除层析特有的伪影结构，防止高斯被错误分配到伪结构上，将模型容量集中于真实物体表示。仅使用~3\%~的全视角数据即可超越全数据优化迭代方法的性能。

Noh et al.（2025.08）~\cite{2508_15151} 结合扩散先验和有符号~3D~高斯实现零样本~CT~超分辨率。\textit{与~R$^{2}$-Gaussian~的差异。}R$^{2}$-Gaussian~从低密度投影重建~CT~体积，但不处理重建后体积的分辨率提升问题。Noh et al.~提出零样本~CT~超分辨率框架——无需配对的高低分辨率训练数据。核心创新是~NAB-GS（Negative Alpha Blending）：标准~3DGS~的~$\alpha$-混合~假设密度非负，但扩散模型上采样后的~2D~投影与低分辨率上采样投影之间的残差可正可负——需要高斯基元同时建模正密度（添加结构）和负密度（移除伪影）。NAB-GS~允许高斯密度取负值，以学习扩散生成高分辨率投影与上采样低分辨率投影之间的有符号残差。框架包含两阶段：(1)~LR CT~投影超分辨率——在丰富~X~射线数据上训练扩散模型上采样低分辨率投影；(2)~3D CT~体积重建——用~NAB-GS~从增强投影中重建~3D~体积。在两个公开数据集上展现优越性能（MICCAI 2026~早期接收）。

InsideOut（2025.10）~\cite{2510_17864} 融合~RGB~和~X~射线的辐射高斯~Splatting，桥接表面外观与内部结构。\textit{与~R$^{2}$-Gaussian~的差异。}R$^{2}$-Gaussian~仅处理~X~射线模态的重建，InsideOut~解决跨模态融合问题：如何用统一的高斯表示同时编码~RGB~表面外观和~X~射线内部结构。核心挑战在于~RGB~和~X~射线成像的物理差异——~RGB~捕获视角相关的表面反射，X~射线捕获各向同性的密度积分。InsideOut~提出层次化拟合策略：先独立训练~RGB~和~X~射线辐射高斯，再通过对齐将两种表示融合到统一坐标系中。引入~X~射线参考损失确保内部结构的一致性。该方法在医学诊断、文化遗产修复和制造检测等领域展现了统一~RGB-X~表示的潜力（ICCV 2025）。

\subsection{MRI与多模态重建}
\label{ss:19_5}

本节介绍将 3DGS 适配到欠采样 k 空间 MRI 重建的方法。

\subsubsection{MRI重建}

本路线介绍将 3DGS 适配到欠采样 k 空间 MRI 重建和动态 MRI 重建的方法。

3DGSMR（2025.02）~\cite{2502_06510} 首次将~3DGS~用于欠采样~k~空间数据的~3D MRI~各向同性分辨率重建。此前，MRI~重建主要依赖两类表示：体素网格（每个体素~1~个参数，但需要数百万体素）和隐式神经表示（MLP~参数化，但推理慢且难以控制）。3DGSMR~的核心思想是用~3D~高斯点云替代传统体素网格来表示~MR~体积——每个体素位置的值由~$M$~个可学习~3D~高斯的叠加近似：$x_j = \sum_{i=1}^{M} G_i^3(j \mid \sigma_i, p_i, \Sigma_i)$，其中~$\sigma_i$~为复值中心密度（匹配~MR~图像的复数性质），$p_i$~为中心位置，$\Sigma_i = R_i S_i S_i^T R_i^T$~为协方差矩阵（可分解为旋转~$R_i$~和缩放~$S_i$）。该方法在自监督框架下运行——不需要完整的~k~空间数据作为监督，而是利用欠采样数据本身的数据一致性约束：优化目标为~$\mathcal{L} = \|A(\hat{X}) - b\|^2 + \lambda \cdot \text{TV}(|\hat{X}|)$，其中~$A$~是由线圈灵敏度图、欠采样掩码和~FFT~组成的前向算子，$b$~为实际采集的欠采样~k~空间数据，TV~正则化作用于重建体积的幅度以抑制噪声。

在实现层面，3DGSMR~有三个关键设计。(1)~\textbf{初始化}：利用欠采样~k~空间数据的~iFFT~结果作为粗重建，从中随机采样~$M$~个网格点作为高斯中心，缩放初始化为到最近三个邻居的平均距离，密度值按因子~$k$~缩放——这一选择利用了~iFFT~结果在~MRI~场景中免费可获取的特性。(2)~\textbf{长轴分裂策略}：标准~3DGS~的克隆+分裂机制在~MRI~中产生过度重叠和高频细节模糊，3DGSMR~改为沿高斯的最长轴一分为二（其余两轴缩放~$0.85$~倍），确保分裂后两个子高斯在最长轴上不重叠，维持整体形状一致性——这一改进显著提升了高加速因子下的高频细节重建。(3)~\textbf{瓦片化体素化管线}：借鉴~3DGS~的光栅化策略，将~3D~体积划分为非重叠~tile，每个~tile~并行处理其关联的高斯点，利用~GPU~共享内存实现高效片上数据复用。这一管线是连接高斯表示与体素化输出的关键桥梁——训练完成后，需要将连续高斯场转换为离散体素以供临床使用。对于~$400 \times 303 \times 165$~的~MR~体积，3DGSMR~将高斯点数限制在~40~万以内（共~$400\text{K} \times 12 = 4.8\text{M}$~参数），远少于密集体素网格所需的~$39.9\text{M}$~参数。

M-Gaussian（2026.03）~\cite{2603_00145} 适配多堆栈厚切片~MRI~采集场景，引入磁高斯基元（Magnetic Gaussian Primitives）和物理一致体积渲染，专门处理~MRI~采集几何中厚切片导致的层面间分辨率差异，在~FeTA~数据集上达到~40.31 dB PSNR~且快~14$\times$。\textit{与~3DGSMR~的差异。}3DGSMR~假设~MRI~数据在~3D~空间中均匀采样——即~k~空间在三个维度上同等加速。然而临床实践中，多堆栈厚切片采集（如胎儿脑~MRI）在层面内（in-plane）和层面间（through-plane）的分辨率严重不对齐——层面内分辨率约~1 mm~而层面间可达~3--5 mm~。3DGSMR~的各向同性高斯基元无法适应这种各向异性采样。M-Gaussian~引入三个核心改进：(1)~\textbf{磁高斯基元}——针对~MRI~采集几何设计的高斯基元参数化，适配厚切片层面间分辨率差异；(2)~\textbf{神经残差场}——在主高斯场上叠加轻量神经网络预测的高频残差，弥补高斯基元在表示精细纹理方面的不足；(3)~\textbf{多分辨率渐进训练}——从低分辨率开始逐步提升训练分辨率，先捕获全局结构再细化局部细节。这是~3DGS~首次成功适配多堆栈~MRI~重建，在~FeTA~数据集上达到~40.31 dB PSNR~且比~INR~方法快~14$\times$。

Terpstra et al.（2026.03）~\cite{2603_21980} 将~3DGS~扩展到动态~MRI~重建，训练~60~秒即可获得~0.5 ms/动态的推理速度，利用~3DGS~的显式表示优势避免了~NeRF~逐点采样的推理瓶颈。\textit{与~3DGSMR~的差异。}3DGSMR~处理静态~MRI~重建（单时间点体积），Terpstra et al.~解决动态~MRI~问题（如心脏~MRI~的时间序列体积重建）。核心差异在于将时间维度纳入高斯表示：为每个高斯基元附加时间变形参数，通过时间条件化预测位置和密度的时变演化。在~R=16~高加速因子下获得高质量心脏~MRI。更重要的是，该工作发现高斯基元的物理属性（位置、密度、运动幅度）直接编码了心脏的生理特性——无需额外后处理即可从训练好的高斯模型中提取心室壁运动、射血分数等临床指标。训练仅需~60~秒，推理速度~0.5 ms/动态，远超~NeRF~方法的逐点采样推理速度（ISMRM 2026）。

PINGS-X（2025.11）~\cite{2511_11048} 针对~4D Flow MRI~超分辨率提出轴对齐时空高斯表示，引入归一化~3DGS~和收敛保证。\textit{与~3DGSMR~的差异。}3DGSMR~重建~MR~图像的标量强度值（解剖结构），PINGS-X~重建~4D Flow MRI~的速度场向量（血液流速的~3D~分量~$v_x, v_y, v_z$~随时间~$t$~变化）——这是一个~4D~向量场重建问题。PINGS-X~引入三个关键创新：(1)~\textbf{归一化高斯~Splatting}——对高斯基元的贡献进行归一化，提供形式化的收敛保证，确保优化过程的稳定性；(2)~\textbf{轴对齐高斯}——将~4D~时空高斯简化为沿各轴独立的~1D~高斯乘积，大幅简化高维数据的训练复杂度同时保持精度和收敛保证；(3)~\textbf{高斯合并过程}——防止退化解并提升计算效率。在~CFD~和真实~4D Flow MRI~数据集上，PINGS-X~相比~PINN~方法大幅减少训练时间，同时实现更优的超分辨率精度（AAAI 2026）。

MoRe-3DGSMR（2025.05）~\cite{2505_04959} 利用~3DGS~实现自由呼吸肺部~MRI~的运动分辨重建，通过患者特异性~CNN~估计变形矢量场。\textit{与~3DGSMR~的差异。}3DGSMR~处理静态~MRI~重建，假设患者在扫描过程中保持静止。肺部~MRI~面临的核心挑战是呼吸运动——患者无法屏气足够长时间获取高质量~3D~数据。MoRe-3DGSMR~提出运动分辨重建管线：首先用~golden-angle~径向采样轨迹采集自由呼吸数据，从~k~空间中心提取呼吸运动信号；然后将~k~空间数据按运动状态排序为多个呼吸相位；用~3DGS~重建参考运动状态的图像体积；最后训练患者特异性~CNN~估计变形矢量场（DVFs），将参考体积变形到其他运动状态。设计动机在于：3DGS~提供连续空间表示实现体素间的平滑插值，而~CNN~建模呼吸运动的非线性变形——两者结合在无需屏气的情况下实现运动分辨的高分辨率肺部~MRI~。在~6~例受试者数据上验证，信噪比和对比度噪声比均优于现有~SOTA~方法。

Krainovic et al.（2024.12）~\cite{2412_18584} 探究了多种~3D MRI~分辨率鲁棒重建方法，发现~3DGS~和神经表示等模型方法在分辨率泛化上不及数据驱动的多分辨率训练策略。\textit{研究性质。}不同于上述方法提出新框架，该工作是系统性基准研究。核心发现是：现有基于~2D~扩散先验的~3D MRI~重建方法（包括~3DGS~和~INR~等模型方法）均对固定体素尺寸敏感——在训练分辨率处性能最优，但当体素尺寸变化时（临床常见情况）性能显著退化。研究比较了~3DGS、神经表示和无穷维扩散模型等模型方法与简单的数据驱动策略（在多种分辨率上训练扩散模型），发现数据驱动的多分辨率训练策略在不损失精度的前提下实现了分辨率鲁棒性，而所有模型方法均无法弥合性能差距。这一发现对~3DGS~在~MRI~重建中的实际部署具有重要指导意义：模型内在的泛化能力可能不及数据层面的多样化训练。

\begin{takeaway}
CT/MRI 断层成像重建是 3DGS 方法论创新的试验田：X 射线体渲染和 k 空间 splatting 需重新推导成像物理的渲染方程，3DGS 的显式表示使高斯属性直接编码心脏生理指标，实现了从重建到诊断的一步跨越。
\end{takeaway}

%=====================================================================

\subsection{近期方法脉络：高斯断层重建、Volumetric MRI resampling with Gaussian Splatting、Dental volume generation with Gaussian Splatting}

这一组近期工作延续 3D Gaussian tomography、3DGS，主要把改动集中在分析、适配等环节。下面只保留相对既有路线的变化，便于和前文方法直接对照。

\paragraph{Projection-Volume Fidelity Divergence}
Projection-Volume Fidelity Divergence（2026.06）~\cite{2606_22525} 延续 3D Gaussian tomography 在“高斯断层重建”上的基本设定，主要改变是：Projection-Volume Fidelity Divergence(PVFD)：刻画各向异性高斯变形和视角特化primitive协同导致的优化漂移。3DGS 的显式高斯表示和可微渲染能力保持为基础。
其设计动机是：揭示3DGS类显式高斯在Radon稀疏约束下可能过度适应投影域而损害真实体积，提供表示级漂移分析。
具体而言，论文分析各向异性高斯变形和view-specific co-adaptation如何让投影损失持续下降但体积质量退化，并设计PVFD mitigation策略。
适用边界是：聚焦tomography场景；缓解策略可能引入额外正则和计算开销。

\paragraph{Rendering Novel Views of MRI Using 3D Gaussian Splatting}
Rendering Novel Views of MRI Using 3D Gaussian Splatting（2026.06）~\cite{2606_26236} 延续 3DGS 在“Volumetric MRI resampling with Gaussian Splatting”上的基本设定，主要改变是：将 3D Gaussian Splatting 适配为体积重建并渲染解剖对齐切面。3DGS 的显式高斯表示和可微渲染能力保持为基础。
其设计动机是：将 3DGS 从自然图像 NVS 扩展到 MRI 重采样，服务于脊柱狭窄等局部病变严重程度评估。
具体而言，从稀疏 MRI 图像形成 Gaussian 体积重建，再根据目标解剖结构渲染更优切面，供后续 ordinal severity grade 预测。
适用边界是：医学体数据与 3DGS 原始相机成像假设差异较大，临床可信度、各向异性采样和诊断安全性需要更充分验证。

\paragraph{X-Splat}
X-Splat（2026.07）~\cite{2607_02099} 延续 3DGS 在“Dental volume generation with Gaussian Splatting”上的基本设定，主要改变是：X-Splat：面向牙科体生成的 Gaussian Splatting 框架。3DGS 的显式高斯表示和可微渲染能力保持为基础。
其设计动机是：把 PXR 到 CBCT 的高度欠定问题转化为几何驱动的 Gaussian 体表示生成，避免隐式/生成方法过平滑或解剖幻觉。
具体而言，针对曲线路径 X-ray 投影导致深度解剖不可观测的问题，引入 Gaussian Splatting 的显式局部表示来定位尖锐解剖边界并提供几何监督。
适用边界是：单视图医学体生成天然欠定，仍需临床验证和不确定性估计以避免有害幻觉。
\section{显微成像与生物结构重建}
\label{sec:19_3}

3DGS~在显微成像中的应用涵盖冷冻电镜（cryo-EM）、电子断层扫描和光学显微镜等多个子方向，核心挑战在于从极端噪声和稀疏投影中重建纳米级~3D~结构。

\subsection{显微与纳米结构重建}
\label{ss:19_6}

本节讨论三条技术路线：（1）冷冻电镜重建；（2）电子断层扫描与原子重建；（3）光学显微与体积电镜。

\subsubsection{冷冻电镜重建}

本路线介绍将 3DGS 应用于冷冻电镜实空间重建和层次构象建模的方法。

\paragraph{GEM：实空间 cryo-EM 重建}

Cryo-EM~重建方面，GEM（2025.09）~\cite{2509_25075} 将~3DGS~直接应用于冷冻电镜的实空间重建。理解~GEM~需要先了解~cryo-EM~的成像链路：蛋白质在溶液中被快速冷冻后，由电子束照射产生~2D~投影图像，该图像可建模为~$I_i = C_i * \text{Proj}(V_i; \phi_i, t_i) + \epsilon_i$，其中~$C_i$~为对比度传递函数（CTF），$\text{Proj}(\cdot)$~为沿电子束方向的投影算子，$\phi_i$~和~$t_i$~为旋转和平移参数，$\epsilon_i$~为加性噪声。传统方法分为两类：傅里叶空间方法利用~Fourier~切片定理，只需预测频域切片，避免实例化完整体积；实空间方法实例化整个~3D~体积再执行投影，内存需求为立方量级。GEM~选择实空间路线，但用~3D~高斯替代体素网格来解决内存瓶颈——每个高斯仅需~11~个参数（3~位置+3~缩放+4~旋转+1~密度系数），蛋白质密度表示为所有高斯的叠加：$V(x) = \sum_{j=1}^{M} \sigma_j \exp(-\frac{1}{2}(x-p_j)^T \Sigma_j^{-1}(x-p_j))$。与原始~3DGS~不同，GEM~采用各向同性的、基于密度的公式（不含视角相关效果），以匹配~cryo-EM~的纯密度积分成像模型。

GEM~最关键的技术贡献是\textbf{闭合形式的~3D$\to$2D~投影}：将~3D~高斯沿光轴~$z$~积分，可直接得到~2D~高斯混合——$I_i(r) = \sum_{j=1}^{M} \sigma_j \sqrt{2\pi|\Sigma_{j}^{2D}|^{1/2}} \exp(-\frac{1}{2}(r-\mu_j^{2D})^T (\Sigma_j^{2D})^{-1}(r-\mu_j^{2D}))$，其中~$\mu_j^{2D}$、$\Sigma_j^{2D}$~分别是~3D~对应量去掉第三行/列后的~2D~均值和协方差。这一闭合形式边缘化避免了在~3D~网格上显式查询密度再聚合，大幅降低计算量和内存。投影后，GEM~在傅里叶空间施加~CTF：$\hat{I}_i^{\text{pred}} = \mathcal{F}^{-1}(\mathcal{F}(C_i) \cdot \mathcal{F}(\tilde{I}_i))$，然后用标准~$\ell_2$~损失与带噪观测比较——虽然~3DGS~通常在干净自然图像上训练，但~GEM~发现标准~MSE~损失对带噪~cryo-EM~图像同样有效。为降低内存开销，GEM~还引入阈值剪枝策略：仅保留对给定射线贡献非可忽略（$G_j^{2D} > \tau$）的高斯，按~$z$~轴排序后从最低~$z$~值（最靠近传感器）到最高值累积——优先处理影响最大的高斯，进一步提升计算效率。训练后，蛋白质密度可通过查询训练好的高斯提取。相比传统体素表示，GEM~的训练速度提升~48\%、内存降低~12\%、局部分辨率提升~38.8\%。

\paragraph{CryoSplat：正交投影适配}

CryoSplat（2025.08）~\cite{2508_04929} 针对标准~3DGS~与冷冻电镜物理的三个根本不匹配提出专用适配：(i) 将透视投影替换为正交投影——cryo-EM~中电子束近似平行投影，镜头模型为正交而非透视；(ii) 将~$\alpha$-混合~替换为物理正确的线积分——电子透射成像的像素值是~3D~电位沿光束方向的线积分，而非~alpha~合成的前到后累积；(iii) 将光栅化坐标系与~FFT~网格对齐——cryo-EM~重建广泛使用~FFT~加速，坐标系必须与频域计算一致才能保证梯度传播的正确性。这些适配使~CryoSplat~能够从随机初始化（无需外部体积或原子模型）进行端到端可微重建，充分利用~3DGS~的瓦片光栅化加速策略。

\paragraph{CryoSPIRE：层次构象建模}

CryoSPIRE（2025.06）~\cite{2506_09063} 提出\textbf{层次高斯混合模型}处理非刚性构象灵活性和组成变化。\textit{与~GEM~的差异。}GEM~假设蛋白质具有单一刚性结构，所有高斯基元共同表示一个静态~3D~密度——这在蛋白质存在多种构象状态（如~IgG~抗体~Fab~域的摆动、核糖体亚基的逐步组装）时不成立。CryoSPIRE~的核心改进是引入\textbf{两阶段层次框架}：第一阶段进行\textbf{部件发现}，第二阶段构建\textbf{部件感知的层次~GMM}。

\textbf{部件发现阶段。}给定初始刚性重建（由反投影或~CryoSPARC~获得），CryoSPIRE~首先优化一个粗粒度~GMM（$G=2{,}048$~个各向同性高斯，每个高斯额外携带一个~24~维可学习特征向量~$f_i \in \mathbb{R}^E$）。关键发现是：训练完成后，这些特征向量\textit{自发地}编码了结构异质性信息——一致运动或同步出现/消失的高斯会获得相似的特征表示。这一现象无需任何显式监督即可涌现，使得后续通过~k-means++~对特征空间聚类即可自动发现~3D~部件。空间上过大的聚类进一步按~3D~位置再分割，确保每个部件被锚点充分覆盖。聚类中心的特征和最近高斯的位置分别初始化锚点的特征~$\phi_a$~和位置~$\mu_a$，形成一个锚点集合~$\mathcal{A} = \{(\mu_a, \phi_a)\}_{a=1}^A$。

\textbf{部件感知层次~GMM。}基于发现的部件，CryoSPIRE~构建层次化的~GMM——每个高斯关联到一个锚点（即一个部件），其位置和特征以锚点为参考系的偏移量表示：$c_i = \mu_{a_i} + \Delta c_i$，$f_i = \phi_{a_i} + \Delta f_i$。异质性建模在\textbf{两个层级}展开：\textbf{锚点级}——对每个锚点施加刚性体变换（旋转~$R_a^{(n)} \in SO(3)$~和平移~$t_a^{(n)}$），由~MLP~以锚点特征和每图像潜变量~$z^{(n)} \in \mathbb{R}^4$~为输入预测，捕获部件级别的大尺度刚体运动；\textbf{高斯级}——在同一~MLP~的独立输出分支上为每个高斯生成局部偏移~$t_i^{(n)} \in \mathbb{R}^3$，捕获部件内部的细粒度柔性形变。对于\textbf{组成异质性}（部件的有/无），CryoSPIRE~通过~sigmoid~激活的~MLP~分支对每个高斯的振幅进行调制：$m_i^{(n)} = m_i \cdot \sigma_i^{(n)}$，其中~$\sigma_i^{(n)} \in (0,1)$——接近~0~表示该高斯在当前构象下缺失，接近~1~表示存在。这种层次设计避免了扁平~GMM~在状态数增多时的模式崩塌：锚点级变换保证大尺度运动的物理合理性（整体刚性），高斯级偏移则允许局部区域的独立柔性。

\textbf{设计动机。}选择层次而非扁平结构的原因是：蛋白质的结构异质性天然具有\textit{分段光滑性}——邻近区域可能以截然不同的方式运动或出现/消失，纯空间位置编码无法捕获这种不连续的分段行为。层次结构通过可学习特征（而非固定的位置编码）实现分段光滑性：属于同一部件的高斯共享相似的锚点特征，从而产生一致的形变模式；不同部件则通过不同的锚点特征实现独立的运动。消融实验证实，使用位置编码替代可学习特征会导致定量性能下降，且无法识别有意义的部件。

\textbf{实验结果。}在~CryoBench~基准的三个数据集上，以~Per-Conformation~FSC~曲线的~AUC~为度量（越高越好），CryoSPIRE~取得~IgG-1D~0.402（均值$\pm$标准差~0.002）、IgG-RL~0.386（0.014）、Ribosembly~0.427（0.014），全面超过广泛使用的~3DVA（0.351/0.341/0.375）、3DFlex（0.336/0.339/无法处理组成异质性）、CryoDRGN（0.366/0.349/0.415）和~DRGN-AI（0.366/0.355/0.372），与线性子空间方法~RECOVAR~在~IgG~数据集上以显著优势胜出（0.402~vs.~0.391，0.386~vs.~0.372），在~Ribosembly~上表现持平（0.427~vs.~0.430，无统计差异）。CryoSPIRE~是首个在~CryoBench~上成功评估的~GMM~方法。在真实实验数据上（EMPIAR-10076~大肠杆菌核糖体亚基组装体、EMPIAR-10180~前催化剪接体），CryoSPIRE~自动发现的部件与原始研究中手动标注的结构模块高度一致（如剪接体的足部、体部、解旋酶和~SF3b~四个区域），进一步验证了部件发现机制的有效性。

\subsubsection{电子断层扫描与原子重建}

本路线介绍基于 3DGS 的暗场电子断层扫描重建和原子级直接重建方法。

\paragraph{DenZa-Gaussian：暗场层析重建}

电子断层扫描方面，DenZa-Gaussian（2026.04）~\cite{2604_04693} 将~3DGS~适配到暗场扫描透射电子显微镜（ADF-STEM）层析重建。\textit{与~GEM~的差异。}GEM~面向明场~cryo-EM~成像，像素值为~3D~电位的直接线积分，散射物理与~X~射线~CT~类似。ADF-STEM~采用环形暗场检测器，仅收集高角度散射电子——像素值不是简单的线积分，而是~3D~电位的复杂非线性函数（涉及电子多重散射）。这一根本不匹配意味着不能直接复用~GEM~的闭合形式投影公式。DenZa-Gaussian~从~Rutherford~弹性散射物理出发，将~ADF-STEM~的散射强度建模为局部原子密度与~$Z^\alpha$~（$Z$~为原子序数，$\alpha \in [1.3, 2.0]$~为修正指数，论文取~$\alpha = 1.7$）的乘积。由于纳米材料的元素分布未知且当前目标是形貌重建而非定量元素分析，DenZa-Gaussian~将物理量~$\rho \cdot Z^\alpha$~简化为单一可学习标量场，称为\textbf{denza~系数}~$d_i$——每个~3D~高斯携带一个~$d_i$，作为该纳米区域综合散射能力的"宏观代理"，将不可观测的微观原子排列与可测量的~ADF-STEM~强度直接关联。这使得高斯基元从原始~3DGS~的颜色载体转变为~ADF-STEM~散射强度载体，能够自适应匹配异质区域的散射变化（如核壳纳米颗粒的壳层与核心、复合材料中的孔隙与基体）。

除~denza~系数外，DenZa-Gaussian~还引入\textbf{散射视角一致系数}~$\eta_i$——一个基于协方差行列式比的缩放因子。标准~3DGS~在~3D$\to$2D~投影时忽略了协方差相关的缩放变化，导致不同倾斜角下同一高斯的散射贡献不一致（违反~Rutherford~散射强度守恒）。$\eta_i$~校正这一偏差，确保~$d_i$~的散射强度在所有投影角度下保持物理一致。投影过程建模为锥形电子束沿路径的线性积分：$P(u) \approx \sum_i \eta_i d_i G_{2D,i}(u)$，将物理散射强度直接转化为图像对比度。

在架构上，DenZa-Gaussian~采用\textbf{两阶段流水线}：第一阶段用~FDK（Feldkamp-Davis-Kress）快速生成粗~3D~体积以捕获整体结构拓扑（而非随机初始化），第二阶段将粗体积转换为~3D~高斯点云，联合优化~denza~系数和~$\eta_i$~系数。损失函数结合四个互补约束：（1）逐像素~$L_1$~损失保证灰度一致性；（2）\textbf{2D~傅里叶幅度损失}——这是应对稀疏视图的关键设计，逐像素损失只能约束局部灰度但无法约束全局频域结构，而缺失楔形伪影本质上是频域中特定角度范围的频率缺失，傅里叶幅度损失通过加权强调高频区域（更易受角度欠采样和噪声影响），促进空间频率相干性，恢复被缺失楔形效应抑制的精细结构细节；（3）基于~SSIM~的感知损失保持高对比度边缘（如核壳界面）；（4）\textbf{3D~全变分正则化}（3D TV）在体素化体积上惩罚相邻体素的突变，推广至~3D 8~邻域配置（含轴向和对角方向），在鼓励局部平滑的同时保留尖锐过渡。

这些设计使得~3DGS~首次在~ADF-STEM~层析重建中实现了端到端的可微优化。实验在~PtNi~纳米催化剂的真实倾斜序列上进行评估：45~视图充足采样下训练视图~PSNR~达~33.02、测试视图~29.17，显著优于~GENFIRE（26.05）、SIRT（23.58）和~FDK（15.53）；15~视图极度稀疏条件下仍达~PSNR~27.53（测试），验证了对缺失楔形伪影的鲁棒性。值得注意的是，标准~3DGS~在~ADF-STEM~任务上几乎完全失败（PSNR~仅~8.93），表明物理适配的必要性——没有~denza~和~$\eta_i$~的散射建模，3DGS~的渲染管线无法匹配~ADF-STEM~的成像物理。消融实验进一步证实~$\eta_i$~的关键作用：移除后~PSNR~从~31.73~骤降至~10.32，说明~3D-to-2D~投影缩放偏差的校正是不可或缺的。

\paragraph{Singh et al.：原子级直接重建}

Singh et al.（2025.12）~\cite{2512_15034} 从原子电子断层扫描（AET）的重建需求出发，将~3DGS~从连续密度场表示重新定位为\textbf{离散原子结构直接优化}——每个高斯基元严格对应一个原子，位置~$\mu$~即为原子坐标，缩放~$\sigma$~编码原子电子密度的空间展布，振幅~$q$~编码散射强度。原子的静电势可建模为~$f_a(r) = q_a \exp(-\frac{1}{2}(r - r_a)^T \Sigma_a^{-1}(r - r_a))$，整个纳米颗粒的总势场为所有原子势的叠加~$f(r) = \sum_{a=1}^{K} f_a(r)$。\textit{与~GEM~的差异。}GEM~用大量小高斯叠加近似蛋白质的连续密度场，高斯数量在训练中自由增长，最终的高斯基元没有明确的物理对应关系；Singh et al.~反过来用少量高斯直接表示离散原子，要求每个高斯与真实原子一一对应，参数量大幅减少且物理意义明确。这一根本差异带来了两个关键技术设计：

(i)~\textbf{原子物理先验}。由于原子电子密度是球对称的，Singh et al.~强制各向同性约束~$\Sigma_a = \sigma_a I$（GEM~不做此约束，允许各向异性高斯以更好地拟合密度场细节）。同时引入最小原子间距约束：周期性计算所有高斯之间的成对距离，将距离小于阈值的相邻高斯合并为一个（取平均位置、振幅和缩放），从而保证最终结构中没有不物理的过近原子。这两条约束共同确保了高斯到原子的~1:1~对应——没有物理约束时，优化倾向于用多个高斯拟合单个原子（产生等效投影但无法用于结构鉴定），并且不同元素的强度分布无法分离。

(ii)~\textbf{消除~atom tracing~两步流程}。传统~AET~先重建体素化~3D~体积，再通过~atom tracing（局部峰值检测~+ 3D~高斯拟合~+ 模板匹配~+ 人工校验）提取原子位置，过程耗时且易受缺失楔形伪影干扰。Singh et al.~的优化目标直接是原子参数~$\min_{K, r_a, q_a, \Sigma_a} L[p, \hat{p}(\sum_{a=1}^{K} f_a)]$，其中~$K$~（原子数）也是待优化变量——训练结束后，高斯中心即原子坐标，振幅和缩放可用于元素识别（不同元素的~$q_a$~和~$\sigma_a$~分布可分离且遵循原子序数排序），完全绕过了~atom tracing~步骤。在正交层析投影的闭合形式梯度下进行优化（与~GEM~类似，不使用自动微分），损失函数仅用~L1~损失而移除了~3DGS~标准中的~SSIM~项——实验表明~SSIM~会鼓励高斯过拟合投影图像中的噪声细节。

在模拟数据（涵盖面心立方金纳米颗粒~Au、铁铂~FePt、非晶钯~aPd~和高熵合金~HEA~四种典型结构）上，完整投影角度下各方法性能相当；在缺失楔形和稀疏角度条件下，Singh et al.~在重建质量和假阳性率上显著优于~FBP、SART~和~INR~基线。在实验~Zr-Te~双壁碳纳米管数据上验证了概念可行性，训练仅需~$\sim$7~分钟（NVIDIA A6000），最大粒子占用~5~GB GPU~内存。

\paragraph{EMGauss：体积电镜切片插值}

EMGauss（2025.12）~\cite{2512_06684} 将体积电子显微镜（vEM）中的各向异性~3D~重建重新定义为~4DGS~的切片插值问题，利用可变形~3D~高斯的时间变形场建模切片间的~3D~连续结构。\textit{与~GEM~的差异。}GEM~通过闭合形式~3D→2D~投影从多视角投影重建单一静态~3D~密度。EMGauss~的输入不是多视角投影，而是一组沿~$z$~轴稀疏采样的~2D~切片（典型~10\%--20\%），目标是推断未观测切片位置的~3D~体积。

\textit{切片到~4DGS~的映射。}EMGauss~的核心洞察是将切片索引~$t \in [0,1]$~视为~4DGS~中的时间坐标，将切片间插值等价为动态场景的时间插值。具体而言，在规范空间中维护一组~3D~高斯基元~$G_c = \{(\mu_i, S_i, R_i, o_i, c_i)\}_{i=1}^M$，轻量~MLP~变形网络~$\phi$~根据切片坐标~$t$~预测每高斯的偏移：$(\Delta\mu_i, \Delta S_i, \Delta o_i) = \phi(\mu_i, t)$，其中~$\Delta\mu_i = (\Delta x_i, \Delta y_i, 0)$~仅编码横向位移，$\Delta S_i = (\Delta s_x, \Delta s_y, 0)$~仅调整平面内缩放。关键约束是~$z$~坐标和~$z$~缩放被固定为全局常数（$z_i = z_0$，$s_{z,i} = s_{z,0}$），旋转不依赖时间（$\Delta R_i = 0$），确保高斯沿轴向对齐且仅发生平面内变形。这一约束的动机是：体积~EM~的切片间变化主要来自组织结构的横向位移和局部形变，而非沿轴向的整体运动。

\textit{Teacher-Student~自举框架。}由于仅~10\%--20\%~的切片被观测，大量~$z$~轴位置缺乏监督。EMGauss~引入基于~EMA~的~Teacher-Student~自举：维护一个稳定的教师模型（学生模型的指数移动平均，衰减率~$\beta = 0.995$），教师对未观测切片生成伪标签监督学生训练。训练采用三阶段管线：（1）预热阶段（~2k~迭代）冻结变形网络，仅优化规范高斯集；（2）短联合训练（~1k~迭代）启用变形网络与高斯集联合优化；（3）伪监督阶段（~15k~迭代）激活~EMA~教师，在逐步扩大的未观测切片子集上提供伪监督，伪监督损失权重从~$0.1$~线性增长至~$1.0$。

\textit{实验验证。}在~EPFL~小鼠脑（555\,nm~分辨率）和~FIB-25~果蝇脑（888\,nm）数据集上，以~6$\times$~下采样模拟各向异性，EMGauss~在~PSNR/SSIM~上优于~CycleGAN-IR~和~EMDiffuse，且是唯一支持沿~$z$~轴任意位置连续生成切片的方法。在真实各向异性~FANC~数据集（各向异性比~10:1）上同样展现有效的重建质量。下游~SAM2~分割实验证实重建体积的结构连续性优于基线方法。

\subsubsection{光学显微与体积电镜}

本路线介绍基于 3DGS 的光场显微镜重建和体积电子显微镜切片插值方法。

\paragraph{3DGAT：光场显微重建}

光学显微镜方面，3DGAT（2025.05）~\cite{2505_12875} 为傅里叶光场显微镜（Fourier Light-Field Microscopy, FLFM）提出基于~3DGS~的自监督~3D~重建框架。\textit{与~GEM~的差异。}GEM~处理电子显微镜的正交投影（沿光轴直接线积分），3DGAT~处理光学光场显微镜的角度-空间混合数据——FLFM~通过微透镜阵列在频域同时捕获~4D~光场（2D~空间~$\times$~2D~角度），而非逐视角投影。标准~3DGS~的透射光栅化仅支持~3D~到~2D~的单视角投影渲染，无法处理光场的角度维信息。3DGAT~的核心改进是将~3DGS~的透射光栅化替换为\textbf{傅里叶域光场渲染}：将~3D~高斯基元投影到频域平面上，通过傅里叶切片定理将~4D~光场渲染分解为频域操作，使渲染过程与~FLFM~的光学成像物理一致。设计动机在于：标准~3DGS~在空域逐像素渲染，无法区分来自不同角度方向的光线——这正是光场显微镜捕获的核心信息；傅里叶域渲染天然编码角度分辨能力，使高斯基元能同时拟合所有视角的观测。训练完全自监督，以多视角聚焦堆栈为监督信号（无需~3D GT~数据），在空间分辨率和重建精度上优于传统去卷积和基于~NeRF~的光场重建方法。

\begin{takeaway}
显微成像与生物结构重建从极端噪声和稀疏投影中重建纳米级 3D 结构：冷冻电镜的 3D 高斯原子模型超越传统方法达到近原子级分辨率，傅里叶域光场渲染则将 3DGS 的适用范围扩展到光学显微成像。
\end{takeaway}

%=====================================================================


\section{临床应用与专用场景}
\label{sec:19_4}

3DGS~的实时渲染、紧凑显式表示和灵活适配能力使其在多个临床专用场景中展现出独特优势。从交互式解剖可视化和稀疏视角牙科重建，到超声成像和医学图像配准，每个子领域对~3D~重建的需求差异巨大。

\subsection{可视化与专科重建}
\label{ss:19_7}

本节讨论两条技术路线：（1）解剖可视化与医学教育；（2）牙科与口腔重建。

\subsubsection{解剖可视化与医学教育}

本路线介绍利用 3DGS 的实时渲染和显式表示实现交互式解剖可视化和医学教育的方法。

3DGS~的实时渲染和显式~3D~表示特性使其成为交互式解剖可视化的理想工具，为医学教育、手术规划和诊断提供沉浸式体验。

ClipGS（2025.07）~\cite{2507_06647} 引入可裁剪高斯~Splatting~框架，支持~CT~体积数据的交互式电影渲染。此前的基线做法是硬截断（Hard Truncation）：直接用高斯中心~$\mu$~的位置判断可见性——中心在裁剪平面以下则可见，否则不可见。这种方法虽然直观，但将问题简化为二元决策，忽略了高斯与裁剪平面的空间交叉关系，容易陷入局部最优并产生裁剪面上的毛刺伪影。ClipGS~的核心创新是\textbf{可学习截断（Learnable Truncation）方案}：为每个高斯基元引入一个可学习属性~$m \in \mathbb{R}$，将可见性判断从"中心是否在裁剪面下方"改为"高斯的有效重心是否在裁剪面下方"，即~$M = \mathbb{1}[(\mu + \delta) \cdot n < z]$，其中~$\delta$~是相对于均值~$\mu$~的偏移量，$n$~为裁剪面法线，$z$~为裁剪面距原点的距离。为使阶跃函数~$\mathbb{1}[\cdot]$~可微分，ClipGS~采用直通估计器（Straight-Through Estimator）重新定义截断函数：$M = \text{sg}(\text{sign}(m - z)) \cdot \sigma(-(m-z)) + (m-z)$，其中~$\sigma$~为~sigmoid~函数，$\text{sg}(\cdot)$~为停止梯度算子——前向传播保持阶跃行为，反向传播用~sigmoid~梯度近似阶跃梯度，从而端到端优化~$m$~值。

进一步，可学习截断本质上是离散的可见性控制，在裁剪面连续移动时会导致渲染不一致。为此，ClipGS~设计了\textbf{自适应变形调整模型（Adaptive Deformation Adjustment）}：用位置编码+轻量~MLP~将高斯位置和裁剪面参数编码为特征~$f$，再通过三个独立的~MLP~头分别预测位置偏移~$\Delta\mu$、旋转偏移~$\Delta q$~和缩放偏移~$\Delta s$，将最终渲染改写为仅对可见高斯的加权合成：$C = \sum_{i=1}^{N} M_i T_i \alpha_i G_i'(r; \mu_i + \Delta\mu_i, q_i + \Delta q_i, s_i + \Delta s_i) c_i$。这使得裁剪面移动时高斯能连续变形，消除离散跳变。训练采用两阶段策略：先单独训练可学习截断获得粗初始化，再加入自适应变形细化渲染质量，损失函数为~$\mathcal{L} = (1-\lambda)\mathcal{L}_1 + \lambda \mathcal{L}_{\text{D-SSIM}}$，平衡像素保真度与感知质量。该方法使医生能够像"解剖"一样逐层剥离~CT~数据观察内部结构，同时保持实时帧率（156 FPS~渲染和~36.635 PSNR）。ClipGS-VR（2026.01）~\cite{2601_19310} 面向临床场景中外科医生佩戴~VR~头显自由探索~CT~解剖结构的需求，将~ClipGS~从桌面工作站移植到移动~VR~平台（Meta Quest 3）。

核心挑战在于：ClipGS~的可学习截断与自适应变形依赖运行时神经网络推理，其~VRAM~占用超过~6\,GB，无法在独立~VR~头显上实时运行。ClipGS-VR~的解决方案分两步。首先，将神经推理\textbf{预计算烘焙}为显式高斯属性——对每个病例均匀采样~200~个切片平面状态，离线计算~LT~和~AAM~的输出并直接嵌入高斯基元，消除运行时推理开销。在压缩策略上，分离球谐系数为全局共享数组（因为所有切片状态共享相同的高阶~SH~系数），并仅保留切片平面处的差异基元而非复制整个体积，将每个病例从约~6\,GB~压缩至平均~40\,MB。其次，将~ClipGS~原有的单轴裁剪平面扩展为\textbf{6-DoF~任意角度切片}：对每个高斯基元，根据其中心到切片平面的符号距离与沿法线方向的投影半径之比计算透明度系数~$\sigma = \text{clamp}(1/2 + (\mathbf{\mu} \cdot \mathbf{n} - c)/(2 \cdot s_{\mathbf{n}}), 0, 1)$，再调制不透明度~$\alpha' = \alpha \cdot \sigma$，使与切片平面相交的高斯按几何穿透深度平滑淡出，消除硬截断的锯齿伪影。

定量评价表明，该方法在单轴切片模式下达到~33.40\,dB~PSNR~和~0.9698~SSIM（对比基线~30.55\,dB / 0.9542），用户研究的系统可用性量表（SUS，System Usability Scale）评分也显著优于固定单轴切片模式。与~ClipGS~对整个视口进行均匀分辨率渲染不同，Hybrid Foveated Path Tracing（2026.01）~\cite{2601_22026} 利用人眼注视点渲染策略进一步优化性能：解剖可视化对交互性要求极高（医生需要实时裁剪、旋转、调整传递函数），全视口光线追踪在~VR~所需的分辨率和帧率下计算代价过高，因此该系统将渲染分为两个区域——注视点中心区域（约~$20°$~视场角）执行流式体路径追踪（path tracing），外围区域使用基于~Mini-Splatting2~的轻量~3DGS~近似。

系统由三个独立模块组成：轻量~VR~实时查看器、GPU~密集型体路径追踪器和~3DGS~外围模型训练模块，架构上支持拆分渲染——查看器可运行在任何支持~DirectX 12~或~Vulkan~的消费级设备上（包括独立~VR~头显如~Meta Quest 3），而路径追踪和模型训练在远程高性能~GPU~上执行。外围模型的初始构建极为快速：从体数据直接生成~20,000~个表面对齐初始点（约~50\,ms），然后以~12~张~8~spp~图像经~95~步简化训练即可在~300\,ms~内生成约~10,000~个高斯的初始模型（完整重建约~1~秒），远快于标准~3DGS~所需的~1-2~分钟。路径追踪器的注视点图像同时用于持续优化外围模型（以超过~500~次迭代/秒的速度训练），使模型在用户交互过程中不断改善。为掩盖路径追踪的延迟，系统采用\textbf{深度引导重投影}策略：将注视点图像按第一命中点深度数据投影到动态变形网格上，重投影公式为~$p'_{\text{clip}} = P_{\text{new}}V_{\text{new}}\,p_{\text{world}}$，其中~$p_{\text{world}}$~由旧相机参数和线性深度重建，从而将渲染帧率与显示刷新率解耦。

实验表明，该方法在注视点区域的感知质量~(LPIPS)~优于同等帧时间下的纯路径追踪，同时将准备时间从~3DGS~的~1-2~分钟缩短至不到~1~秒，在~Meta Quest 3~上以~72 FPS~以上运行。

NLI4VolVis（2025.07）~\cite{2507_12621} 将体积可视化的交互范式从空间裁剪推进到语义级自然语言交互。\textit{锚点回顾。}ClipGS~的交互模式是几何操作——用户通过移动裁剪平面逐层"解剖"体积数据，系统响应的是空间位置关系。然而，对于非专业用户（如医学生、安全检查人员），他们更希望用自然语言表达意图（如"显示消化系统相关器官"或"把背包里的方形物体变红"），而非手动调整裁剪面或传递函数参数。NLI4VolVis~正是针对这一需求，构建了一个完整的自然语言驱动的体积可视化系统。

\textit{系统架构与核心设计。}NLI4VolVis~的管线包含四个阶段。（1）\textbf{多视角语义分割}：从体积数据渲染多视角图像后，采用灵活的分割策略——支持预标注数据集、传递函数分割、体素分割、跨帧对象跟踪（SAM 2）以及直接在~3DGS~上分割（LangSplat、SAGD）等五种方式，每个分割出的语义组件独立训练一个可编辑~3DGS~模型（基于~iVR-GS）。（2）\textbf{3D~开放词汇场景理解}：对每个分割组件，用基于信息熵的视角选择机制从均匀采样的多视角图像中挑选信息量最大的~$k$~帧（按~alpha~通道熵值~$H(I)=-\sum_{i=1}^{N} p_i \log p_i$~排序），再通过~CLIP（Contrastive Language-Image Pre-training，对比语言-图像预训练）~模型提取视觉嵌入并取均值~$E_{\text{vision}} = \frac{1}{k}\sum_{j=1}^{k} E_{\text{frame}_j}$，若存在文本标签则融合~$E_{\text{object}} = (E_{\text{vision}} + E_{\text{text}})/2$。用户查询时计算~CLIP~余弦相似度匹配目标对象。（3）\textbf{可编辑高斯训练}：每个语义组件训练独立的~iVR-GS~模型，该模型在标准~3DGS~基础上为每个高斯基元增加法线~$n$、偏移颜色~$\Delta c$、环境光/漫反射/镜面反射系数~$(k_a, k_d, k_s)$~及光泽度~$\gamma$~等可编辑属性，用~Blinn-Phong~反射模型计算视角相关颜色，并支持多个独立模型的参数直接拼接（composability）以组合完整场景。（4）\textbf{多智能体~LLM~自然语言交互}：采用层次化协作的多智能体架构，各~Agent~遵循~VPA~循环（Visualize-Perceive-Act）——核心~Agent~将用户的自然语言意图解析为声明式渲染命令（调整颜色、不透明度、光照方向、视角等），通过函数调用工具执行操作后，再感知当前可视化结果判断是否达成目标，迭代至用户满意或达到最大迭代次数。

\textit{与~ClipGS~的具体差异。}ClipGS~的核心交互是空间裁剪（用裁剪平面切割体积），而~NLI4VolVis~将交互维度从几何层面提升到语义层面——用户不再需要理解裁剪面的空间含义，而是直接描述想看到的结构和期望的视觉呈现。具体而言有三点关键不同：（1）分割粒度不同：ClipGS~按空间位置一刀切，NLI4VolVis~按语义含义分割为独立可编辑的组件（如将鲤鱼~CT~数据分为骨骼、鱼鳍、内脏等七个组件）；（2）编辑自由度不同：ClipGS~只能控制裁剪面的位置和角度，NLI4VolVis~可以对每个语义组件独立调整颜色、不透明度、光照和风格迁移；（3）智能程度不同：ClipGS~需要用户手动操作，NLI4VolVis~通过~LLM~多智能体自动解析意图、选择最佳视角、执行~2D~风格化（基于~InstructPix2Pix），并提供语音输入/输出支持。用户研究表明，8~名参与者（含非~CS~背景）平均~9.9~分钟熟悉系统后，~35.1~分钟内完成~23~项任务，自然语言使用率达~88.1\%，领域知识问答准确率~8.75/9，系统可用性评分~4.38/5。上述方法均需对每个数据集单独训练~3DGS~模型，这在批量临床应用中效率受限。Render-FM（2025.05）~\cite{2505_17338} 提出基础模型从~CT~体积直接回归~6DGS（六维高斯~Splatting）参数，消除逐扫描优化。\textit{与~ClipGS~的差异。}ClipGS~仍需对每个~CT~数据执行逐场景优化训练（~2~阶段策略），Render-FM~的核心洞察是医学~CT~的渲染不需要逐场景优化——不同患者的解剖结构虽有差异，但~CT~数据的密度分布遵循可学习的通用模式。

\textit{模型架构。}Render-FM~采用~3D U-Net~骨干网络（受~nnU-Net~启发），输入~6~通道~CT~体积：归一化~HU~值、分割掩码和~4~通道~RGBA~传递函数值。解码器输出~37~通道参数体积~$\Theta \in \mathbb{R}^{37 \times D/2 \times H/2 \times W/2}$（半分辨率），每个体素位置对应一个潜在~6DGS~基元。37~个通道编码：均值方向~$\mu_d$（3~通道）、球谐系数~$L=0$~偏移和~$L=1$（12~通道）、不透明度偏移（1~通道）和~$6\times6$~协方差矩阵~$\Sigma$（21~通道，对角线取对数、非对角线取~$\tanh$~激活保证正定性）。

\textit{解剖引导启动（AGP）。}标准~3DGS~用~SfM~点云或随机初始化，对体积医学数据不适用。Render-FM~提出~AGP~先验：每个体素位置直接作为高斯中心，从传递函数推导基础颜色和不透明度，分割掩码决定仅在前景区域实例化基元——将解剖先验编码到~6DGS~参数的初始化中。

\textit{6DGS~视角相关渲染。}实例化后的~6DGS~基元通过协方差矩阵的视角条件化切片（slicing）实现视角相关渲染：将~$6\times6$~协方差分块为空间~$\Sigma_{pp}$、方向~$\Sigma_{dd}$~和交叉~$\Sigma_{pd}$~项，给定视角方向~$v_i$~后计算有效~3D~空间属性~$\mu'_{p,i} = \mu_{p,i} + \Sigma_{pd}\Sigma_{dd}^{-1}(v_i - \mu_{d,i})$~和~$\Sigma'_{pp,i} = \Sigma_{pp} - \Sigma_{pd}\Sigma_{dd}^{-1}\Sigma_{dp}$，以及方向分量导出的不透明度调制因子~$w_i = \mathcal{N}(v_i | \mu_{d,i}, \Sigma_{dd})$。切片后的~3D~高斯经标准~3DGS~分块光栅化渲染。

\textit{训练与推理。}端到端训练使用~$L_1$ + D-SSIM~渲染损失。推理时，新~CT~体积经单次前向推理（~$<3$~秒，GPU）即可得到~6DGS~表示，支持实时交互渲染。无分割掩码时自动调用~TotalSegmentator~生成。整个流程从原始~CT~到交互渲染仅需数秒（对比逐场景优化的~$\sim$1~小时），可选~2~分钟微调进一步提升质量。Niedermayr et al.（2024.04）~\cite{2404_11285} 展示了压缩~3DGS~在消费级设备上实现电影解剖学的可行性。\textit{与~ClipGS~的差异。}ClipGS~聚焦交互式裁剪渲染的算法创新，但仍依赖桌面级~GPU~运行。Niedermayr et al.~从工程部署角度出发：如何将高质量~3DGS~解剖可视化从工作站移植到消费级设备。

\textit{电影解剖学管线。}该工作提出完整的电影解剖学（Cinematic Anatomy）管线：首先对~CT~数据选择预设（传递函数、材质分类、裁剪平面），用物理路径追踪器为每个预设渲染多视角高保真图像，然后用~3DGS~拟合这些图像生成高斯表示，最后通过压缩实现消费级设备上的实时渲染。管线嵌入~Mip-Splatting~处理不同细节层级的平滑过渡。

\textit{压缩策略。}核心创新是灵敏度感知向量量化（Sensitivity-Aware Vector Quantization）加熵编码：对高斯属性按其对渲染质量的影响排序，高灵敏度参数保留更高精度，低灵敏度参数激进量化。量化后各属性流经增量编码和~LZMA~熵压缩。最终将数~GB~数据集压缩至~70\,MB~以下，压缩率超过~$50\times$，同时在~WebGPU~渲染器中保持视觉质量。渲染器在浏览器中运行，支持~GPU~排序和~2D splat~光栅化，使~3DGS~解剖可视化首次可在平板电脑和手机等消费级设备上实时运行，为临床教育和患者沟通提供了低成本、零安装的部署方案。

\subsubsection{牙科与口腔重建}

本路线介绍从极稀疏口内照片重建牙齿和口腔 3D 结构的方法。

牙科场景的独特挑战在于：标准临床采集中仅有~3~张稀疏口内照片（正面和双侧颊侧），且口内环境光照不一致、牙齿表面高度反光。这使得牙科成为~3DGS~在极稀疏视角条件下的极限测试场景——从仅~3~个视角恢复完整的~3D~牙齿和口腔结构，远低于自然场景中常用的~50--200~个视角。

DentalSplat（2025.11）~\cite{2511_03099} 对~3DGS~的初始化和训练管线进行了针对性的修改——用~DUSt3R~先验引导替代~SfM~点云初始化，用尺度自适应剪枝替代标准致密化，并新增光流几何约束与梯度碰撞修正——3DGS~的可微分光栅化渲染核心则保持不变。

\textit{Baseline 做法与问题。}标准~3DGS~依赖~COLMAP~通过特征匹配恢复~SfM~点云，从仅~3~个无位姿口内照片初始化时，特征匹配几乎完全失败——正面与双侧颊侧照片之间视角差异巨大（约~$60^\circ$--$90^\circ$），重叠区域有限，COLMAP~无法建立足够对应关系，点云极度稀疏甚至完全为空。即使强行用随机初始化训练，稀疏视角下高斯缺乏充分的多角度约束，优化过程中高斯基元倾向于通过增大尺度"补丁"式覆盖未观测区域，产生大量退化的扁平高斯。此外，稀疏视角还引发\textbf{梯度碰撞}（Gradient Collision）问题：每个~3D~高斯影响多个像素，而每个像素受多个高斯影响，当不同像素的梯度方向相互冲突时，累积梯度相互抵消，导致~3DGS~的分裂（split）操作无法确定正确的致密化方向，结果在过重建区域产生模糊伪影。实验表明，标准~3DGS~在~3~视角设定下完全无法收敛。

\textit{核心改进：先验引导初始化。}DentalSplat~的第一步是用~DUSt3R（密集立体先验模型）替代~COLMAP~获取初始点云。DUSt3R~是一种预训练的几何基础模型：给定两张无标定的图像，它直接回归逐像素的~3D~点图（pointmap）——即每个像素对应的~3D~坐标——同时输出相对相机位姿。DUSt3R~通过大规模多场景训练获得了强大的跨视角稠密对应能力，能够在无需特征匹配的情况下为~3~张口内照片生成全局对齐的稠密点云。这解决了~COLMAP~在极端稀疏视角下的初始化失败问题，为~3DGS~提供了"粗但合理"的几何起点。

\textit{核心改进：尺度自适应剪枝（SAP）。}DUSt3R~生成的点云极为稠密（可达数百万点），直接用于~3DGS~初始化会严重拖慢优化并降低收敛质量。DentalSplat~提出尺度自适应剪枝策略（Scale-Adaptive Pruning, SAP），在高斯初始化后根据空间影响过滤异常点。直觉上，尺度过大的高斯通常是不合理的"补丁"，而非对精细结构的真实刻画。SAP~对全体高斯的缩放参数~$\mathbf{S} \in \mathbb{R}^{N \times 3}$~构建双掩码过滤：第一个掩码筛除最大缩放分量超过全局均值的高斯，$\mathcal{M}_1 = \{\mathbf{S}_i \mid \max(\mathbf{S}_i) > \mu_S\}$，其中~$\mu_S$~为所有缩放分量的均值幅度；第二个掩码根据点云规模自适应调阈值——当点云较小（$N < 5 \times 10^6$）时使用分位数阈值~$\mathcal{M}_2 = \{\mathbf{S}_i \mid \max(\mathbf{S}_i) > Q_i(\mathbf{S})\}$，否则使用更严格的~$4\mu_S$~阈值。最终掩码取交集~$\mathcal{M}_{\text{final}} = \mathcal{M}_1 \cap \mathcal{M}_2$，仅保留通过双重检验的高斯。这种双掩码设计确保在大场景中不会过度剪枝（丢失关键几何），在小场景中依然能有效剔除离群点。

\textit{核心改进：光流几何约束。}极稀疏视角下，仅靠光度损失不足以约束~3D~几何。DentalSplat~引入光流作为额外的几何约束。直觉是：如果~3DGS~正确重建了场景，那么将训练好的高斯从视角~$t$~投影到视角~$t+1$~时，像素位移应与真实的光流一致。具体而言，对于像素处覆盖的~$K$~个高斯，DentalSplat~通过~alpha~加权混合计算"高斯流"：$\text{flow}^G(\mathcal{T}_{t+1}, \mathcal{G}_t) = \sum_{i=1}^{K} w_i (\boldsymbol{\mu}_{i,t+1} - \boldsymbol{\mu}_{i,t})$，其中~$w_i$~为归一化混合权重，$\boldsymbol{\mu}_{i,t}$~为第~$i$~个高斯在视角~$t$~下的~2D~投影中心。由于~DentalSplat~使用各向同性高斯（协方差矩阵经~Cholesky~分解简化为单位矩阵），高斯流等价于投影中心的位移。光流损失定义为：$\mathcal{L}_{\text{flow}}(\mathbf{x}_{t_1}) = \|M(\mathbf{x}_{t_1}) \odot \text{flow}^G(\mathbf{x}_{t_1}) - M(\mathbf{x}_{t_1}) \odot \text{flow}^{Gt}(\mathbf{x}_{t_1})\|_2$，其中~$\text{flow}^{Gt}$~为由双向光流模型预测的伪真值光流，$M(\mathbf{x}_{t_1})$~为双向光流得到的逐像素置信度掩码——稀疏场景中光流预测噪声较大，置信度掩码仅保留可靠像素的光流约束，避免噪声梯度污染优化。

\textit{核心改进：梯度碰撞修正。}针对稀疏视角下的梯度碰撞问题，DentalSplat~对累积梯度取绝对值操作：$\hat{g}_{i,x} = \sum_{j=1}^{n} |\partial L_j / \partial \mu_{i,x}|$，$\hat{g}_{i,y} = \sum_{j=1}^{n} |\partial L_j / \partial \mu_{i,y}|$。这一操作的本质是将不同像素的梯度方向沿坐标轴对齐——不再让正负梯度相互抵消，而是让它们在各自坐标轴上累积——从而确保~3DGS~的分裂操作能够沿正确的方向致密化。消融实验表明，移除梯度约束后~PSNR~从~28.41~降至~27.58，是影响最大的单项组件。

\textit{实现细节与训练目标。}整个训练过程联合优化相机位姿~$\hat{\mathcal{T}}_{t+1}$~和~3DGS~参数~$\hat{\mathcal{G}}$~，目标函数为~$\hat{\mathcal{T}}_{t+1}, \hat{\mathcal{G}} = \arg\min_{\mathcal{T}, \mathcal{G}} (\lambda_1 \mathcal{L}_{\text{rgb}} + \lambda_2 \mathcal{L}_{\text{flow}})$，其中~$\mathcal{L}_{\text{rgb}}$~为渲染图像与真实图像的~$\mathcal{L}_1$~残差。训练仅需~2000~次迭代即可收敛。

\textit{设计动机。}牙科~3~视角重建是~3DGS~在极稀疏条件下的极限测试场景。关键洞察是：从~3~张照片直接训练~3DGS~几乎不可能收敛到合理解——初始化质量决定了上限。DUSt3R~等密集立体重建模型提供了跨视角的稠密对应关系，将~3DGS~的起点从"零"提升到"粗但合理"。然而~DUSt3R~的稠密输出又是~3DGS~的负担——SAP~策略将稠密先验修剪为~3DGS~可高效优化的稀疏表示。光流约束和梯度修正则分别从外部几何一致性和内部优化动力学两个维度弥补稀疏视角的信息不足。

\textit{实验验证。}在~950~个临床病例（3~视角极稀疏图像集）和~195~个视频测试集（3/6/9/12~视角递减测试）上验证。在~3~视角设定下，标准~3DGS~完全无法收敛；InstantSplat~达到~32.78 PSNR / 0.945 SSIM；DentalSplat~达到~34.50 PSNR / 0.954 SSIM / 0.135 LPIPS，训练仅需~69~秒（单卡~RTX 4090）。

Dental3R（2025.11）~\cite{2511_14315} 进一步将牙科极稀疏重建推进到完全无位姿设定——既无相机位姿也无内参，仅需~3~张未标定的口内照片即可完成重建。\textit{与~DentalSplat~的锚点关系。}DentalSplat~用~DUSt3R~替代~COLMAP~初始化，但仍假设相机内参已知，且将所有~$\binom{n}{2}$~图像对直接送入~DUSt3R~进行全局点图回归。在仅~3~张照片的极稀疏设定下，这种"全配对"策略暴露出三个失败模式：（i）稠密配对（dense pairing）在传播~DUSt3R~特征时内存和计算量爆炸，而单参考配对（one-reference star pairing）则对大基线约束不足；（ii）正面与双侧颊侧照片之间大基线、不一致光照和釉质高反光导致相机恢复和辐射场优化不稳定，产生漂移、浮空伪影和过重建；（iii）稀疏视角下纯光度监督产生频率偏差，将釉质边缘和邻接面细节过度平滑。Dental3R~针对这三个问题分别设计了图引导配对策略和小波正则化目标。

\textit{核心改进一：几何感知配对策略（GAPS）。}Dental3R~的关键洞察是：DUSt3R~的配对质量直接决定重建上限，而非所有~$\binom{n}{2}$~配对都等价——选择重叠度不足或基线过大的配对不仅浪费计算，还会引入错误的几何约束。GAPS~将图像排列在环图~$C_n$~上，定义顶点集~$V = \{0, 1, \ldots, n-1\}$，并通过多尺度邻域生成候选边集：$E_{\text{cand}} = \bigcup_k \{\{i, (i+k) \bmod n\} : i \in V\}$，其中~$k$~为固定的正整数偏移集合。边的几何重要性由环图上的环绕索引距离建模：$d(i,j) = \min(|i-j|, n - |i-j|)$，该距离经过单调衰减函数~$\gamma$~映射为重叠度代理值，再转换为范围自适应重要性权重：$w(i,j) = \alpha_{r(i,j)} \cdot \gamma(d(i,j)) + \beta_{r(i,j)}$，其中~$r(i,j) \in \{\text{local, global}\}$~为由距离~$d(i,j)$~推导的粗粒度范围类别。最后，GAPS~在候选图中求解度约束的最大权重子图：$\max_{E' \subseteq E_{\text{cand}}} \sum_{\{i,j\} \in E'} w(i,j) \quad \text{s.t.} \quad \deg_{E'}(v) \leq b \;\; \forall v \in V$，其中~$b$~为统一的度上限，防止出现枢纽节点。这一设计使得~DUSt3R~的对应关系计算聚焦于高价值配对，在~3~视角设定下仅需~3~个有效配对（而非~3~个全配对），内存占用大幅降低。

\textit{核心改进二：小波正则化目标。}DentalSplat~仅使用~$\mathcal{L}_1$~光度损失加光流约束训练~3DGS，在极稀疏视角下产生频率偏差——低频分量被充分约束但高频细节（釉质边缘、邻接面轮廓）被过度平滑。Dental3R~引入基于二级离散小波变换（DWT）的频域正则化。具体而言，对~GT~图像~$I_{\text{gt}}$~和渲染图像~$\hat{I}$~分别执行二级~DWT，得到四个子带：近似子带~LL、水平细节~LH、垂直细节~HL、对角细节~HH。每个子带的残差定义为~$\Delta_x = W_x(I_{\text{gt}}) - W_x(\hat{I})$，其中~$W_x$~为提取子带~$x \in \{\text{LL, LH, HL, HH}\}$~的算子。总小波损失为各子带加权~$\mathcal{L}_2$~范数之和：$\mathcal{L}_{\text{wav}} = \sum_{x \in \{\text{LL,LH,HL,HH}\}} \lambda_x \|\Delta_x\|_2^2$。关键在于~LH/HL/HH~三个高频子带的显式约束——它们分别捕获水平边缘、垂直边缘和对角纹理，正是釉质边界和邻接面细节在频域中的对应分量。通过在频域强制这些高频分量与~GT~一致，小波正则化有效抵消了稀疏视角下光度损失的频率偏差。

\textit{与~DentalSplat~的具体差异总结。}（1）在初始化策略方面：DentalSplat~直接使用~DUSt3R~的全配对输出进行初始化，假设内参已知；Dental3R~则通过~GAPS~选择最优配对子集后再用~DUSt3R~回归点图，同时求解外参和内参。（2）配对策略：DentalSplat~未考虑配对选择问题；Dental3R~的~GAPS~通过环图多尺度邻域和度约束子图选择，将内存占用从~$O(n^2)$~降至~$O(nb)$。（3）训练目标：DentalSplat~使用~$\mathcal{L}_{\text{rgb}} + \mathcal{L}_{\text{flow}}$~联合优化；Dental3R~使用~$\mathcal{L}_{\text{rgb}} + \mathcal{L}_{\text{wav}}$~，用频域约束替代光流约束，直接在频域保持临床关键的高频细节。

\textit{设计动机。}Dental3R~的出发点是远程正畸场景——患者在家用手机拍摄~3~张口内照片上传，系统自动完成~3D~重建，无需诊所级硬件（CBCT/口内扫描仪）和专业操作人员。这一场景的约束比~DentalSplat~更严格：不仅视角极稀疏，而且相机完全未标定（无内参无外参）。GAPS~解决的是"在有限计算资源下选择最有价值的几何约束"；小波正则化解决的是"在缺乏充分多角度约束时如何保持临床关键的高频细节"。

\textit{实验验证。}在~195~个视频测试集（3/6/9/12~视角递减测试）上，Dental3R~在所有视角设定下均优于~CF-3DGS~和~InstantSplat。在~6~视角设定下，CF-3DGS~输出存在明显模糊和浮空伪影，InstantSplat~在下颌牙齿区域产生严重几何扭曲；Dental3R~的重建保持几何保真且无伪影。在~3~视角极稀疏设定下（950~个临床病例），标准~3DGS~无法收敛，Dental3R~仍能产生临床可用的重建结果。消融实验表明：移除~GAPS~（改用全配对或单参考配对）导致~PSNR~下降且~GPU~内存显著增加；移除小波约束导致釉质边缘细节模糊。DentalSCR（2025.11）~\cite{2511_14343} 研究口内扫描（IOS，Intraoral Scanner）3D~模型与头影测量~X~光片（CR，Cephalometric Radiograph）的~3D-2D~跨模态对齐问题。\textit{与~DentalSplat~的差异。}DentalSplat~从口内照片重建~3D~牙齿模型，DentalSCR~解决不同模态间的空间对齐——将口内扫描的~3D~表面模型与头颅~X~光片配准，为正畸诊断和手术规划提供统一的~3D~坐标系。

\textit{U-中线牙弓轴（UMDA）坐标系。}DentalSCR~首先为上下颌~IOS~模型建立统一的解剖参考系~$(X, Y, Z)$：$Z$~轴对齐咬合面法向（通过顶点法向与咬合面参考对齐后的~PCA（Principal Component Analysis，主成分分析）~最小特征向量确定），$X$~轴沿~U~形牙弓前向（通过臼齿基线与切牙参考点拟合），$Y$~轴取双颌平均以确保跨颌一致性。UMDA~提供三个临床可解释的解剖平面：咬合面（$X$-$Y$）、正中矢状面（$X$-$Z$）和冠状面（$Y$-$Z$）。

\textit{表面~DRR~投影。}传统~DRR~需要体积密度数据（CT/CBCT），而~IOS~模型仅为三角网格表面，缺乏内部密度信息。DentalSCR~提出\textbf{表面~DRR~公式}：将网格视为均匀密度~$\rho_0$~的空壳，通过对每个表面元素在探测器平面上累积投影贡献生成等效厚度图~$L(u,v) = \sum_{i=1}^{N} w_i \exp(-((u - y'_i)^2 + (v - z'_i)^2)/(2\sigma_i^2))$，再代入~Beer-Lambert~定律~$I(u,v) = I_0 \exp(-\rho_0 L(u,v))$~生成~X~射线模拟图像。进一步，采用高斯~Splatting~策略近似连续厚度分布，解决稀疏表面点直接投影产生的锯齿和断裂伪影。

\textit{轮廓配准。}投影后的~IOS~轮廓与~CR~轮廓通过对称~Chamfer~距离最小化进行配准：$\mathcal{L}_{\text{Chamfer}} = \frac{1}{|C_{\text{IOS}}|}\sum_{p \in C_{\text{IOS}}}\min_q \|p-q\|_2 + \frac{1}{|C_{\text{CR}}|}\sum_{q \in C_{\text{CR}}}\min_p \|q-p\|_2$，估计~4~参数相似变换（缩放~$s$、旋转~$\theta$、平移~$t_x, t_y$）。采用三层粗到精优化策略：粗阶段探索大参数范围拒绝伪假设，精阶段缩窄搜索窗口提升对齐度，超精阶段进行亚像素精修。

\subsection{成像与配准方法}
\label{ss:19_8}

本节讨论两条技术路线：（1）超声成像；（2）医学图像配准。

\subsubsection{超声成像}

本路线介绍将 3DGS 适配到超声成像物理以实现体积重建的方法。

超声成像具有安全、廉价和实时的优势，但其~2D~视野受限且成像物理与光学相机截然不同。核心差异在于：光学相机通过反射光测量物体表面，受遮挡影响大；超声探头则通过发射超声波测量不同深度的反射回波，能够穿透不透明体积。这意味着标准~3DGS~的累积透射率~$T_j$~概念不再适用，因为不透明体积不应被视为遮挡。

UltraGauss（2025.05）~\cite{2505_05643} 首次提出超声专用的高斯~Splatting~框架，系统性地重新设计了成像模型。其核心创新在于将渲染机制从"投影"改为"相交检测"：不将~3D~高斯投影到~2D~图像平面，而是检测~3D~高斯与探头平面在~3D~空间中的交集。具体而言，像素~$x$~处的不透明度通过~3D~马氏距离计算：$\alpha_i(x) = \sigma_i \exp(-\frac{1}{2} d_{3D}^2)^{1/2}$，其中~$d_{3D}^2$~是将~2D~图像点提升到~3D~探头坐标系后的马氏距离。这实现了概念上的反转——RGB~渲染评估~2D~投影后的~2D~高斯，而超声渲染在~3D~空间中评估原始~3D~高斯与探头平面的关系。

在协方差参数化方面，UltraGauss~发现标准~3DGS~的四元数-缩放参数化在超声场景中出现数值不稳定（四元数归一化约束导致梯度异常），因此提出三角协方差参数化：用下三角矩阵~$L$~表示逆协方差~$\Sigma^{-1} = LL^T$，对角线元素强制为正（加小常数~$\epsilon$），既保证正定性又无需归一化。该参数化还带来高效的优势——通过前向替换即可求逆，并通过~$y = \mu + L^{-T}z$（$z \sim \mathcal{N}(0, I)$）直接采样。此外，UltraGauss~设计了基于~$\chi^2$~分布的两阶段光栅化边界剔除策略：第一阶段剔除与探头平面不相交的高斯（$b_3^- > 0$~或~$b_3^+ < 0$），第二阶段仅对通过的高斯在~2D~包围盒内迭代像素，实现线程间负载均衡。在~5~分钟内完成重建，SSIM~达~0.99。

UltraGS（2025.11）~\cite{2511_07743} 引入物理解耦渲染（PD Rendering），将超声信号分解为低阶球谐函数捕获组织的宏观反射模式，同时用一阶波动效应建模声波在组织中的散射和衰减，两者结合实现高效的强度合成，达到~64.69 FPS~的实时渲染。\textit{与~UltraGauss~的差异。}UltraGauss~重新设计了超声渲染的几何管线（从投影改为相交检测），但强度合成仍采用简单的不透明度加权。UltraGS~进一步适配超声的声学物理：PD~Rendering~将超声信号分解为两个物理分量——低阶~SH~捕获组织界面的宏观镜面反射（超声图像中的亮线结构），一阶波动效应建模声波在组织中的散射和衰减（超声图像中的散斑噪声和深度衰减）。此外，UltraGS~引入深度感知高斯基元（depth-aware Gaussian primitives）配合可学习视场角（learnable FOV，Field of View），改善无约束探头运动下的几何一致性。在三个数据集上达到~SOTA：PSNR~达~29.55，SSIM~达~0.89，渲染速度~64.69 FPS（ICME 2026）。

TVGS（2025.08）~\cite{2508_14552} 针对经阴道超声的自由扫掠提出无监督~3D~重建框架。核心挑战在于自由扫掠既无外部跟踪设备也无法预知扫掠轨迹，系统必须同时估计每帧切片的~6-DoF~位姿和~3D~解剖结构。TVGS~引入切片感知的可微光栅化器适配超声物理（非光学投影），联合优化切片位姿和解剖结构。\textit{与~UltraGauss~的差异。}UltraGauss~假设已知探头位姿（或通过外部跟踪获取），专注于从已知位姿的~2D~切片重建~3D~体积。TVGS~解决更根本的问题：自由扫掠中连探头位姿都未知。其核心贡献是联合优化方案——将切片的~6-DoF~位姿作为可学习参数与高斯基元的几何属性一起优化，通过图像级监督信号同时精修位姿和解剖结构。切片感知光栅化器确保超声物理（非光学投影）在可微管线中得到正确建模。结果是一个紧凑、灵活、内存高效的体积表示，以高空间保真度捕获解剖细节。

Ultra-Wide-NeRF（2026.04）~\cite{2604_24187} 利用多元~3D~高斯（MVG）实现宽视野超声重建，通过距离依赖的凸体积采样和各向异性~3D~高斯建模复杂波束几何。\textit{与~UltraGauss~的差异。}UltraGauss~处理标准线阵或凸阵探头的~3D~重建。Ultra-Wide-NeRF~面向宽视野（Wide Field-of-View）重建——将多个重叠扫掠拼接为覆盖更大解剖区域的~3D~体积。核心挑战在于凸阵探头产生的发散声束几何：随着深度增加，声束宽度增大、分辨率降低，导致传统拼接方法产生严重的混叠伪影。MVG~表示通过距离依赖的凸体积采样（distance-dependent convex volumetric sampling）显式建模这一波束几何，各向异性~3D~高斯沿声束方向拉伸以匹配深度相关的分辨率变化。这种方法天然提供抗混叠效果，并生成连续神经表示支持从任意虚拟轨迹合成高保真新视角。在心内超声（ICE）的体模和猪心数据上验证，扩展了术中导航所需的空间上下文。

\subsubsection{医学图像配准}

本路线介绍利用 3DGS 的显式高斯控制点表示实现医学图像可变形配准的方法。

3DGS~的显式稀疏表示在医学图像配准中展现出独特优势：通过少量高斯控制点即可建模连续变形场，兼具几何可解释性和计算效率。本节介绍基于高斯控制点的可变形图像配准方法。

GaussianDIR（2024.06）~\cite{2406_03394} 将~3DGS~的显式表示思想引入医学图像可变形配准（DIR），核心目标是在~3D~体积图像间估计连续位移向量场（DVF），将运动图像对齐到固定图像。

\paragraph{基线做法与局限}
传统~DIR（Deformable Image Registration，可变形图像配准）~方法（如~B-spline FFD（Free-Form Deformation）、LDDMM）将配准建模为迭代优化问题，虽然精度高但耗时长且需大量超参数调优。深度学习方法（如~VoxelMorph）通过~3D U-Net~直接预测~DVF，推理速度快但依赖大量训练数据且泛化能力差。隐式神经表示（INR）方法（如~IDIR）用~MLP~直接优化坐标映射函数，精度好但训练仍远慢于实时，且隐式表示缺乏可解释性。

\paragraph{高斯控制点表示}
GaussianDIR~的核心改进是将~DVF~从密集体素/隐式~MLP~表示替换为稀疏的高斯控制点集合。每个控制点~$G_i$~由高斯参数（中心位置~$\mu_i \in \mathbb{R}^3$、缩放~$s_i$、旋转~$r_i$）和局部变换参数（平移~$T_i \in \mathbb{R}^3$、旋转~$r_i'$）联合定义。直觉上，这些控制点相当于``空间锚''——少数稀疏锚点各自管辖一片局部区域，通过插值即可覆盖整个~3D~体积的变形场。初始控制点布置在~$10 \times 10 \times 10$~的均匀~3D~网格上（仅~1000~个），参数量仅占现有方法的不到~5\%。

\paragraph{变形场插值：线性混合蒙皮}
对于每个体素~$x_j$，系统通过~KNN~搜索找到~$K$~个最近高斯控制点，用高斯概率密度计算归一化权重：

\begin{equation}
w_{jk} = \frac{\exp\bigl(-\tfrac{1}{2}(x_j - \mu_k)^\top \Sigma_k^{-1}(x_j - \mu_k)\bigr)}{\sum_{k' \in \mathcal{N}_j} \exp\bigl(-\tfrac{1}{2}(x_j - \mu_{k'})^\top \Sigma_{k'}^{-1}(x_j - \mu_{k'})\bigr)},
\end{equation}

体素的位移向量通过加权求和得到：$\Phi_j = \sum_{k \in \mathcal{N}_j} w_{jk}\,\phi_k$，其中~$\phi_k = R_k'(x_j - \mu_k) + \mu_k + T_k$~是第~$k$~个高斯的局部刚体变换。这一策略不仅简化了体素运动的计算，还天然提供了平滑正则化——相邻体素共享相似的最近高斯，因此位移场自动保持空间连续性。实验表明仅使用~3 DoF~平移（不含旋转）即可在精度-速度权衡中取得最佳效果。

\paragraph{自适应密度控制}
不同解剖区域的变形复杂度差异很大：肺部边界区域变形剧烈需密集控制点，而均匀软组织区域稀疏控制点即可。GaussianDIR~借鉴~3DGS~的自适应密度策略，根据高斯位置梯度的范数动态调整控制点数量——梯度范数超过阈值~$\tau_{\max}$~的高斯被克隆（沿随机方向偏移复制），低于~$\tau_{\min}$~的被剪枝。此外，系统采用多尺度训练（迭代步数从~125~逐步增加到~4000），先捕获粗粒度全局变形再细化局部细节，减少陷入局部最优的风险。

\paragraph{优化与结果}
整个管线通过反向传播端到端联合优化高斯参数和变换参数，损失函数仅使用归一化互相关（NCC）而无需额外正则化项。在~DIR-Lab~肺部~4D-CT~数据集上，GaussianDIR~以平均~TRE~1.01\,mm、处理时间~2.43\,秒超越~IDIR（1.07\,mm）和~ccIDIR（1.04\,mm），参数量仅~0.019M（不足~IDIR~的~5\%）。在~ACDC~心脏~MRI~数据集上，平均~Dice~系数~0.879~超越~IDIR（0.862）和~ccIDIR（0.852）。作为无需训练的逐案例优化方法，GaussianDIR~打破了``迭代方法本质上很慢''的刻板印象。

\subsection{多模态融合与体积重建}
\label{ss:19_9}

本节介绍多模态医学成像联合重建与体积数据的 3DGS 表示方法。

\subsubsection{多模态医学成像与体积重建}

本路线介绍多模态联合重建、切片式体积重建、扩散光学断层成像和动态表面跟踪等方法。

3DGS~的显式高斯表示在多模态医学成像和体积重建中展现出灵活的适配能力：从切片式数据的时空重建到内部结构可视化，从光学漫射层析到多模态融合。本节介绍~7~种方法，涵盖折叠高斯建模、聚焦感知渲染、条件~Splatting、扩散光学断层成像和动态表面跟踪等方向。

MedGS（2025.09）~\cite{2509_16806} 提出多模态~3D~医学成像的高斯~Splatting~框架，支持~CT~和~MRI~等多种模态的联合~3D~重建——核心挑战在于不同模态的成像物理差异（CT~测量~X~射线衰减，MRI~测量磁共振信号），需要统一的~3DGS~表示来桥接。

\paragraph{基线做法与局限}
医学切片数据的~3D~重建传统上依赖~Marching Cubes（移动立方体算法）~等方法从体素网格提取网格，但容易产生连通性伪影且受限于分割精度。隐式神经表示（INR，Implicit Neural Representation）虽能实现高保真重建，但计算量大、难以解释和手动编辑。纯~3DGS~则缺乏对切片间连续性的建模——当输入帧稀疏时（如每隔~2--5~帧才有一张切片），帧间几何容易断裂。

\paragraph{折叠高斯与双任务架构}
MedGS~基于~VeGaS~所提出的折叠高斯（Folded-Gaussian）分布构建。与标准~3DGS~中每个高斯仅表示静态~3D~位置不同，折叠高斯将时间维度作为高斯的一维，通过可学习的非线性平移函数~$f(t)$~和缩放函数~$a(t)$~建模切片间的时空变形：

\begin{equation}
\mathcal{FN}(\mathbf{m}, \mathbf{\Sigma}, a, f)(\mathbf{x}) = \mathcal{N}(m_s + f(m_t - t),\, a(t)\mathbf{\Sigma}_s)(s|t) \cdot \mathcal{N}(m_t, \sigma_t^2)(t),
\end{equation}

其中~$s|t$~是条件空间变量，$m_s$~和~$\mathbf{\Sigma}_s$~控制空间位置和形状，$f(\cdot)$~用~7~次多项式参数化以捕获复杂的非线性切片间变形。MedGS~的关键设计是每个高斯基元同时携带两组颜色属性——图像强度系数~$\mathbf{c}_i^{\text{img}}$（灰度重建）和分割掩码系数~$\mathbf{c}_i^{\text{seg}}$（二值占用场）——共享几何参数但外观分离，迫使图像合成的密集梯度信号正则化几何学习。

\paragraph{帧间正则化（IBFR）}
医学切片数据常有噪声和采样稀疏问题。MedGS~引入帧间正则化（In-Between Frame Regularization）机制：对相邻训练帧~$I_t$~和~$I_{t+1}$，按~$\alpha \sim \mathcal{U}(0.2, 0.8)$~线性混合生成伪真值帧~$I_{t_\alpha} = \alpha \cdot I_t + (1-\alpha) \cdot I_{t+1}$，强制折叠高斯在切片间平滑过渡，防止过拟合噪声或退化为离散层。时间扩展参数~$\sigma_t$~被约束在~$[2/T, 1]$~范围内（$T$~为训练帧数），防止高斯过于局部或过于弥散。总损失函数为~$\mathcal{L}_{\text{total}} = \mathcal{L}_{\text{img}} + \lambda_{\text{seg}}\mathcal{L}_{\text{seg}} + \lambda_{\sigma}\mathcal{L}_{\sigma}$，其中~$\lambda_{\text{seg}} = \lambda_{\sigma} = 0.001$。

\paragraph{结果}
MedGS~在前列腺超声、MRI~脑部（TotalSegmentator）和心脏数据集上验证。帧插值任务中，隔帧采样时~PSNR~达~33.52\,dB（线性插值~32.70，光流~29.07）。网格重建中，Chamfer~距离在前列腺数据集上平均~0.187（Poisson~0.208，FUNSR~0.226），双任务联合模型优于仅用分割分支的单任务变体（MedGS$_s$），证实图像信号对几何正则化的贡献。训练使用~100k~个折叠高斯组件，在~A100~上约~20~分钟完成。GaussianPile（2026.03）~\cite{2603_20611} 提出统一的稀疏高斯~Splatting~框架用于切片式体积重建。\textit{与~MedGS~的差异。}MedGS~用折叠高斯建模切片间的时空连续性，通过时间扩展参数~$\sigma_t$~控制切片间的平滑过渡。GaussianPile~采用根本不同的物理路线：将成像系统的点扩散函数（PSF）直接嵌入渲染管线，从成像物理层面建模切片厚度效应。

\textit{聚焦感知建模。}GaussianPile~的核心创新是将成像系统的有限聚焦深度显式纳入~3DGS~的渲染管线。定义系统灵敏度图~$h(-z_c) = \exp(-\sigma_c^2/(2\sigma_z^2))$，其中~$\sigma_z$~反映成像系统在轴向（层析方向）的聚焦能力。对每个~3D~高斯基元，渲染强度为高斯与灵敏度图的卷积——这在数学上等价为三步操作：(1)~\textbf{轴向重参数化}——将~$\sigma_z$~注入逆协方差矩阵~$\Sigma_e^{-1} = \Sigma_c^{-1} + \sigma_z^{-2}e_3e_3^\top$（$e_3 = [0,0,1]^\top$），生成聚焦高斯（Focus Gaussian），使轴向支撑范围收缩以匹配有限切片厚度；(2)~\textbf{不透明度调制}——通过~Mahalanobis~距离变化量~$\text{opacity}_r = \exp(-\frac{1}{2}(\mu_c^\top\Sigma_c^{-1}\mu_c - \mu_e^\top\Sigma_e^{-1}\mu_e))$~抑制离焦基元的贡献，自然消除跨切片的"鬼影"伪影；(3)~\textbf{加法光栅化}——切片成像中像素强度来自沿投影线的积分叠加而非遮挡合成，因此采用加法累积~$I(p) = \sum_{i} \gamma_i \exp(-\frac{1}{2}d_i^\top\Sigma_{2d,i}^{-1}d_i)$~替代标准~3DGS~的~$\alpha$-混合。

\textit{3D-2D~一致性保证。}聚焦感知模型的一个关键优势是天然保证~3D~体积一致性：同一高斯基元被多个切片图像观测时共享相同的协方差参数，梯度信息直接依赖于聚焦高斯表示——高斯无法在切片位置拟合图像的同时在焦平面外维持不同内部结构。这种内置正则化消除了标准~3DGS~中~2D~投影看似合理但~3D~结构塌缩的问题。

\textit{压缩与量化。}训练收敛后，GaussianPile~采用空间~Morton~排序暴露参数局部相关性，各属性流经自适应精度量化（位置~14~位、不透明度~12~位、缩放对数~12~位、四元数~12~位）和~LZMA~熵编码，实现~$16\times$~压缩。在~ABUS~（自动乳腺超声）和~LSM~（光片显微镜）两种模态上验证，3~分钟完成重建，比~NeRF~方法快~11$\times$。InnerGS（2025.08）~\cite{2508_13287} 针对内部场景重建提出因式分解~3DGS。\textit{与~MedGS~的差异。}MedGS~的折叠高斯从时间维度建模切片间关系，InnerGS~从空间维度分解~3DGS——标准~3DGS~通过投影将~3D~高斯映射到~2D~图像平面，本质上面向外表面合成，高斯基元集中在物体表面而内部体积留空；且投影渲染依赖相机内外参，切片数据（如~MRI~体积扫描）本身不含相机参数，无法直接适用。

\paragraph{条件~Splatting（Conditional Splatting）}
InnerGS~的核心创新是因式分解方案：将~3D~高斯密度沿深度轴~$t$~分解为条件~2D~高斯与边缘~1D~高斯的乘积，$p_i(u,v,t) = p_i(u,v|t) \cdot p_i(t)$。其中~$p_i(t) = \mathcal{N}(\mu_z, \sigma_{zz})$~是沿深度方向的边缘高斯，$p_i(u,v|t)$~是给定切片深度后的条件~2D~高斯，其均值随深度~$t$~线性偏移（$\mu_{u,v|t} = [\mu_x, \mu_y]^\top + \frac{t - \mu_z}{\sigma_{zz}}[\sigma_{xz}, \sigma_{yz}]^\top$），反映~3D~高斯在不同切片上的投影中心移动。条件协方差为~$\Sigma_{uv|t} = \begin{bmatrix}\sigma_{xx} & \sigma_{xy} \\ \sigma_{xy} & \sigma_{yy}\end{bmatrix} - \frac{1}{\sigma_{zz}}\begin{bmatrix}\sigma_{xz}^2 & \sigma_{xz}\sigma_{yz} \\ \sigma_{xz}\sigma_{yz} & \sigma_{yz}^2\end{bmatrix}$，固定深度后平面内不确定性降低。与~MedGS~直接用折叠高斯建模连续变形不同，InnerGS~通过因式分解避免了逐切片重复计算完整~3D~协方差，每个切片获得自适应的紧致包围盒，相比~3D~椭球投影方法将平均包围盒面积从~209.80~降至~63.09、误检数从~15.78~降至~1.82~（像素级），候选高斯数从~16.70~降至~2.62。

\paragraph{重建流程与动态扩展}
高斯在~3D~规则网格（$42^3$）上初始化，对每个~2D~切片通过条件~Splatting~采样候选高斯并按深度排序，沿用~3DGS~的~$\alpha$-混合~渲染，通过标准可微光栅化反向传播优化。重建目标直接是~3D~体素密度而非~2D~投影，因此能合成任意深度的切片。对动态数据（如腕部运动~MRI、脑~fMRI），InnerGS~引入变形场网络：静态~3D~高斯表示解剖基准态，轻量~MLP~预测逐帧的位置、协方差、颜色和不透明度偏移，冻结几何参数仅优化颜色和不透明度即可捕获~fMRI（functional MRI，功能磁共振成像）~的~BOLD（Blood-Oxygen-Level Dependent，血氧水平依赖）~信号时变。

\paragraph{文本引导分割}
下游临床任务中，InnerGS~在每个~3D~高斯上附加语言特征。通过知识增强的医学视觉语言模型（VLM，Vision-Language Model）提取高维语义嵌入，经场景专用自编码器蒸馏为紧凑隐向量~$f_i$，再用可微~Splatting~方程聚合为~2D~语义图：$F(v) = \sum_i^N f_i \alpha_i \prod_{j=1}^{i-1}(1 - \alpha_j)$，最后由轻量解码器投影回语义空间实现开放词汇对齐——避免了显式高维建模的内存爆炸，同时保持体积一致性。

\paragraph{实验结果}
在~BrainWeb~脑部~MRI~上，InnerGS~准确重建皮层褶皱和深部脑边界，均匀区域（白质、脑脊液）保持平滑过渡。在~HVSMR-2.0~心脏~MRI~上，心腔和血管细节清晰。分割任务在~SAT-DS~数据集上取得~mDSC~88.17\%、mIoU~83.72\%、像素精度~98.97\%，肝脏区域~DSC~稳定在~84.62\%~以上。腕部动态~MRI（38~帧）上~4D~重建在~RTX~3090~上约~45~分钟收敛，在大幅运动的时间戳上显著优于静态基线。

GS-DOT（2026.04）~\cite{2604_23675} 首次将~3DGS~适配到漫射光学断层成像（DOT，Diffuse Optical Tomography）。\textit{与~MedGS~的差异。}MedGS~处理标准医学成像模态（CT、MRI、超声），其成像物理基于自由传播的~X~射线或声波。DOT~面临根本不同的物理体制：近红外光在生物组织中经历大量散射（$\mu_s' \gg \mu_a$），光传输遵循扩散方程而非自由传播的射线方程。

\textit{高斯基元参数化。}标准~3DGS~用不透明度和颜色渲染场景外观。GS-DOT~用高斯基元参数化吸收系数分布（逆问题的未知量），每个基元仅携带单一标量参数——吸收对比度振幅~$\alpha_k$：$\mu_a(r) = \sum_{k=1}^{K} \alpha_k \exp(-\|r - c_k\|^2/(2s_k^2))$，其中~$c_k$~为中心、$s_k$~为各向同性尺度。二维情形下扩展为各向异性~2D~高斯，支持半轴~$s_{x,k}$、$s_{y,k}$~和旋转角~$\theta_k$~的独立调整，以适应细长或不规则形状的吸收体。所有参数采用指数重参数化（$\alpha_k = e^{\tilde{\alpha}_k}$，$s_{x,k} = e^{\tilde{s}_{x,k}}$）保证严格正性。

\textit{正向模型与逆问题。}光-组织交互的物理模型为时域扩散方程：$[-\nabla \cdot D(r)\nabla + \mu_a(r) + \frac{1}{v_m}\frac{\partial}{\partial t}]\Phi(r,t) = q(r,t)$，其中~$\Phi$~为光子 fluence~率，$D = 1/(3\mu_s')$~为扩散系数。解析~Green~函数在~Born~近似下给出灵敏度矩阵（Jacobian）~$J_{sd}(r,t)$，将测量扰动与吸收扰动线性关联。逆问题通过优化~$6K$~个高斯参数拟合~100~个源-探测器对的时域测量数据，损失函数包含~$L_2$~数据拟合、正则化（惩罚大振幅和过度各向异性）和排斥项（防止空间重叠的基元冗余）。

\textit{实验验证。}在~6\,cm~直径圆形~2D~域（模拟新生儿头部截面）上，10~源-10~探测器均匀分布，时间分辨率~20\,ps，总窗口~6\,ns。四种测试案例（单包含体、三包含体、月牙形、空心环形），GS-DOT~在干净和~2\%~泊松噪声条件下均展现鲁棒性：RMSE~变化~$< 0.002$\,cm$^{-1}$，SSIM~$> 0.81$，质心定位误差~$< 0.16$\,cm。关键优势在于参数空间从~2184~（体素化未知量）压缩至~$6K$（$K = 1$--$16$），压缩率~22.75--364$\times$。

Nguyen et al.（2026.04）~\cite{2604_06671} 将~4DGS~从静态体积重建扩展到动态血管表面跟踪，用于体外血栓切除（thrombectomy）实验的运动学测量。\textit{与~MedGS~的差异。}MedGS~从切片数据重建静态~3D~体积，聚焦切片间的时空连续性建模；Nguyen et al.~面向完全不同的应用场景——机械血栓切除术中血管壁在器械回撤过程中的时变变形，需要同时重建~4D~几何并跟踪表面运动，而非静态重建。

\paragraph{设计动机}
机械血栓切除是急性大血管闭塞性卒中的有效治疗手段，但器械回撤过程可能造成血管过度拉伸和蛛网膜下腔出血等并发症。现有体外实验多关注再通率、远端栓塞或器械-血栓摩擦力等全局指标，缺乏空间分辨的、时间连续的全场~3D~血管表面运动学测量。Nguyen et al.~的目标是建立从多视角采集到运动学指标提取的标准化端到端管线，量化器械操作期间血管表面的局部变形集中区域。

\paragraph{9~相机同步采集系统}
硬件方面，Nguyen et al.~设计了低成本多视角采集平台：9~台~Arducam IMX586~相机（2160p、20\,fps）安装于十二面体状~PVC~框架上（5~台下环、4~台上环），相机-目标距离约~230--271\,mm，可视角向公共中心区域。整个硬件（PVC~结构、3D~打印件、相机和照明）成本低于~\$1500。硅胶血管模型表面涂覆红色荧光微球（约~300~个），配合~400\,nm UV~照明增强表面纹理对比度。每实验录制~20--30\,s，时间降采样后约~1000~帧用于~4DGS~处理。

\paragraph{4DGS~重建管线}
采集管线包含六个阶段：(i)~多相机标定（$6\times8$~棋盘格，平均重投影误差~0.134~像素）；(ii)~SAM2~视频分割逐视角掩码血管区域、抑制背景；(iii)~4DGS~动态重建（Wu et al.~实现，每序列训练~20\,000~迭代，RTX 4070 Ti~上约~20~分钟）；(iv)~高斯基元聚类为固定连通边图（KMeans~颜色分组~+~DBSCAN~空间聚类，3D Delaunay~邻域图经边长裁剪保留局部拓扑）；(v)~空间相干性滤波——基于固定图邻域的~Huber~型~IRLS~加权平滑，抑制孤立撕裂伪影，保留分段光滑的运动场；(vi)~ROI~位移与应力代理指标提取。

\paragraph{表面运动学测量与应力代理}
运动学分析的核心是在固定连通图上定义两个~ROI~级指标。

\textbf{位移指标}：顶点位移~$\mathbf{u}_k(t) = \mathbf{x}_k(t) - \mathbf{x}_k(0)$，每帧~ROI~中位数为~$d^{\text{med}}_{\text{ROI}}(t) = \text{median}_k \|\mathbf{u}_k(t)\|_2$，条件级汇总取时间最大值~$d^{*}_{\text{ROI}} = \max_t d^{\text{med}}_{\text{ROI}}(t)$。

\textbf{应力代理}：对边~$e=(i,j)$，定义拉伸~$\lambda_e(t) = \ell_e(t)/\ell_e(0)$，映射为不可压缩~Neo-Hookean~单轴应力代理~$\sigma_e = \mu(\lambda_e^2 - \lambda_e^{-1})$，其中~$\mu \approx E/3 \approx 0.383$\,MPa（硅橡胶~$E=1.15$\,MPa）。应力代理是相对比较指标而非绝对壁应力——刚体平移控制实验中，中位应力代理~$\approx 0$\,MPa（90~百分位~0.028\,MPa），验证了其抑制刚体运动伪信号的能力。

\paragraph{实验验证}
验证分为合成控制和体外比较两部分。合成实验用~Blender~渲染已知变形的血管网格，施加~1--5\,mm~局部拉伸。重建精度：对称~Chamfer~距离~1.714--1.815\,mm，$\tau=1$\,mm~时精度~0.964--0.972。运动学一致性：ROI~位移回归斜率~0.987、$R^2=0.992$，Bland-Altman~偏差~-0.142\,mm；应力代理回归斜率~1.068、$R^2=0.969$。体外比较中，颈段~ICA~放置的抽吸导管相比~ICA~末端放置，在远端~M1~段的峰值中位位移为~3.776\,mm vs.\ 1.151\,mm，MCA~分叉处为~3.144\,mm vs.\ 0.846\,mm，对应力代理值亦一致更高（0.109\,MPa vs.\ 0.058\,MPa），表明近端放置导致更广泛的血管表面变形。

\begin{takeaway}
临床应用与专用场景展示了 3DGS 的广泛适用性：从牙科稀疏视角重建到超声弹性成像，从血管运动学测量到医学图像配准，3DGS 的实时渲染和显式表示优势在多样化的临床需求中得到验证。
\end{takeaway}

%% ============================================================

\subsection{近期方法脉络：医学3D补全综述}

这一组近期工作延续 Gaussian splatting，主要把改动集中在增量改进等环节。下面只保留相对既有路线的变化，便于和前文方法直接对照。

\paragraph{Deep Learning Approaches for 3D Medical Scene Completion}
Deep Learning Approaches for 3D Medical Scene Completion（2026.06）~\cite{2606_24180} 延续 Gaussian splatting 在“医学3D补全综述”上的基本设定，主要改变是：医学3D scene completion综述，覆盖voxel、point、implicit和Gaussian splatting相关生成/渲染范式。3DGS 的显式高斯表示和可微渲染能力保持为基础。
其设计动机是：梳理2016-2026年医学3D场景补全表示和生成范式演进，连接扩散先验与实时渲染技术。
具体而言，论文按voxel、point cloud、implicit representation和生成式/GS方向组织文献。
适用边界是：不是3DGS专门医学应用论文；对具体高斯方法技术细节覆盖有限。

\section{近期研究（2026 年 8--11 月）}
\paragraph{XClipGS：半空间精确积分与切面监督}
2026.08~\cite{2608_07760}。交互式切割需要处理横跨平面的高斯。按中心点硬裁会把整个原语保留或删除，容易在切面附近留下错误贡献。XClipGS 在仿射 EWA 模型下直接积分高斯落在保留半空间的部分，并用同一个可微算子训练。切割视图覆盖不同平面位置和相机距离，为原本被外层遮住的体内结构提供监督。

作者在六个 CT 和两个 MRI 体数据上比较切面边缘和错误侧漏。实验包括在同一个体内表示上替换裁剪算子，以区分算子与训练监督的作用。算子可计算新的半空间，无需再次训练，但切面视觉质量仍依赖体内表示。留出偏移位于训练区间内，这组结果主要支持对该区间内新切面的泛化。

\paragraph{用多色 X 射线投影降低锥束 CT 金属伪影}
2026.08~\cite{2608_13159}。锥束 CT 的射线跨越金属时，低能成分往往比高能成分衰减更多，单能模型难以解释随路径变化的衰减。该方法用高斯基元表示局部体结构，以每个高斯的两个系数分别表达近似与能量无关的康普顿项和随能量变化的光电项。可微多色投影将不同能量下的透射信号按系统响应合成，并从实际投影联合估计体场和全局系统响应。

原文将肺、牙和西兰花三个模拟场景分别嵌入铁、钛、铝物体，按体 PSNR 与 SSIM 对照重建，并对五组实际扫描作比较。实际样本的对照条件各异，例如金属杆样本未取得单独的无金属参考。结果支持所测扫描条件下的伪影抑制，不能直接外推到其他扫描设备、材料组合或临床应用。

\paragraph{TR-GS：用重尾原语和射线置信度重建稀疏视角 CT}
2026.08~\cite{2608_16042}。TR-GS 用可投影的 Student t 分布取代放射高斯，以缓解稀疏观测和异常值对体密度估计的影响。局部射线的几何支撑决定原语自由度的正则目标。置信度同时引导三维小波约束，在压制稀疏采样伪影时保留可辨识的高频边界。

作者使用 R²-Gaussian 的合成和真实 CT 数据，在整圈扫描轨迹上均匀选取 12、18 和 25 个视角比较，另外用 6 和 50 个视角做消融。论文报告大部分设置下改善最显著的是少视角情形。整圈均匀稀疏与只覆盖窄角度属于不同的逆问题，后者留待验证。

\paragraph{统计形状约束下的少视角 X 射线股骨重建}
2026.09~\cite{2609_10376}。少视角 X 射线中，不同的三维密度可能产生近似投影。该方法先用群体平均模板初始化高斯的位置、取向和密度，再把原语的位置绑定到统计形状模型，让它们按群体上可见的形变模式共同移动。密度先验另负责约束体强度。

论文从六例留出股骨 CT 合成投影，在五视角下相对 R²-Gaussian 的平均 PSNR 从 34.55 升至 37.38 dB，十视角下提升较小。先验在投影约束最少时作用更明显，结果目前限于该股骨模板和模拟投影设置。

\paragraph{LINGO：改善稀疏 X 射线重建的初始化和低密度梯度}
2026.09~\cite{2609_22849}。LINGO 在 R2-Gaussian 上从投影视图的掩码构造三维体素筛选，抑制背景噪声并初始化高斯。训练用动态损失缩放加强低密度区域的梯度，并以 IPSD 专门评价初始点云结构。

论文在 X3D 比较 5、10 和 25 张 X 射线输入。结果对应数据集的器官和投影条件；IPSD 衡量初始化形态，不能代替最终 CT 体素重建或临床评价。

\paragraph{FactorSplat：保持体渲染代理几何并编辑传递函数}
2026.10~\cite{2610_02382}。体渲染的区域传递函数把组织强度映射为颜色和不透明度。FactorSplat 保留一套高斯位置、协方差和方向相关响应，把传递函数的变化作用于视角无关颜色和不透明度。局部组织与强度查表给出物理外观改动，低秩残差弥合代理表示和原体渲染器间的差异。修改传递函数后，交互转动相机可复用更新后的高斯属性。

作者在七个 CT 与 MRI 扫描中比较留出的传递函数，报告相对于区域感知 VEG 的平均重建和编辑精度提升。逐扫描结果仍有例外，例如心脏及膝关节的编辑误差由 VEG 取得更低值。模型对所支持的传递函数族训练，更广泛的编辑需要单独验证。

\paragraph{CellSplat4D：将显微成像轴向模糊与稀疏时间补全分开}
2026.10~\cite{2610_04199}。CellSplat4D 将遗漏的活细胞显微体数据视为时空补全。它用物理测量固定轴向点扩散函数，避免优化误把光学模糊归咎于原语伸展；时间门以帧为单位，在固定预算下按门控残差增密。连续时间表示可用来渲染未采样时刻。

CE02 和 CHO 的比较让各方法共享同一前向光学模型，并分别检查细胞核、留出轴向层等指标。作者发现离最近采集帧越远时本法的保真度衰减较慢，但部分中心定位指标仍落后。结果说明特定采集协议下的补全表现，未证明任意动态都能忠实恢复。

\paragraph{AD-GS：把稀疏牙科 CBCT 的原语分配与结构约束分开}
2026.10~\cite{2610_14963}。AD-GS 从 FDK 初值建立放射高斯，用衰减密度引导马尔可夫链蒙特卡罗采样，将更多原语分配给牙齿及颌面硬组织。基于 MedSAM 的语义约束负责维持结构一致性，分阶段训练把重点从整体形状转向局部细节。投影一致性仍是重建的主要观测约束。

论文在 ToothFairy3 上用 TIGRE 模拟含噪的 20、30、50 个圆轨迹投影。实验同时报告体图像质量、分割重叠和表面距离，因而能观察细边界是否随重建改善。该协议支持评价模拟少视角条件，真实低剂量扫描仍需相应验证。

\paragraph{面向固定高斯预算的跨个体解剖分配}
2026.10~\cite{2610_20397}。对已知 CT 体进行可交互渲染时，固定预算里每个高斯放在哪里会影响解剖边界的表现。作者从训练个体逐例优化的高斯模型提取每个体素格的典型分配，训练标签感知的神经分配器，并给出一个考虑可见性的解析方案。部署时先按分配密度抽取恰好 B 个中心，再对当前体数据拟合颜色、尺度等属性。

作者比较相同预算下的均匀分配、剪枝和所提方法，在 20 例留出和 19 例外部个体上评估。评价分配器时固定高斯中心，以隔离位置的作用。两种分配器各有外观适用范围：解析法依赖解剖标签仍能解释渲染外观，学习法也受训练时传递函数对比度的影响。

\paragraph{KFS⁴LAM：让内镜关键帧筛选与形变建图协同}
2026.10~\cite{2610_22505}。内镜视频既含冗余相邻帧，也含局部组织运动。KFS⁴LAM 按累计运动提出关键帧候选，再根据图像冗余筛选。建图端按保留帧调节优化预算，用解剖褶皱边界监督与受限时间形变维持结构一致性。

作者在 C3VD 和 C3VD-v2 上以轨迹误差、图像相似度和深度误差评估，并在 SplaTAM 与 EndoGSLAM 主干下做匹配比较。自适应优化预算会增加部分条件下的计算量。作者将结果定位为质量与计算的取舍，并指出外部褶皱检测器的可靠性会影响推广。

\paragraph{X-Splat：用同一组高斯重建 RGB 外观和 X 射线体密度}
2026.10~\cite{2610_23237}。X-Splat 面向同一物体的标定 RGB 图像和 X 射线投影，先用粗 X 射线重建建立跨模态配准和共享几何初值。优化时两种模态共享一组高斯几何，颜色等外观属性集中在近表面原语，X 射线重建给出用于 RGB 的表面深度约束。

实验包括 InsideOut 的四个场景和 Flower \& Insect 的五个场景，并从参考 CT 合成统一协议下的锥束投影。冻结 CT 并单独优化 RGB 的消融不及联合模型。作者报告联合训练主要改善 RGB，X 射线质量约等于单模态基线。透明物体或强次表面散射会破坏近表面外观假设。

\paragraph{H-2DGS：用高度场高斯表达皮层表面}
2026.10~\cite{2610_30910}。H-2DGS 先建立体高斯，再从视图深度融合出截断符号距离场，最后在零水平集上布置曲面高度场原语。高度场拟合局部皮层曲率，颜色和不透明度基函数保留纹理变化。术后配准以 MRI 网格渲染的深度约束目标形状，同时使用颜色和重叠约束维持表面外观。

论文在五例中报告更低的多视图深度不一致，但 PSNR 与 SSIM 逊于体高斯基线。另一个视频中的生理脉动影响 SfM 位姿，导致后续重建失败。结果体现了显式薄表面带来的几何优势及其对外观表达与输入位姿的要求。

\paragraph{CTBakeGS：从已知 CT 体直接烘焙渲染代理}
2026.10~\cite{2610_37406}。CTBakeGS 处理的是已有 CT 密度体及选定传递函数的显示。它直接将体数据和渲染外观转换为标准三维高斯，使用毫米单位的自适应八叉树划分，并在转换中补偿透明度和抗锯齿。输出可用普通高斯查看器显示，转换时无需先渲染训练相机或逐例拟合。

作者在六个 TotalSegmentator v2 扫描、三种传递函数预设下与参考体渲染比较，报告宽视图均方根误差约为满量程的 2.3\% 至 3.1\%。这衡量的是对既有渲染外观的复现，适用范围与传递函数和参考渲染配置绑定。

\paragraph{CryoGaussian：用显式密度高斯表示冷冻电镜异质性}
2026.11~\cite{2611_26128}。异质冷冻电镜重建需要同时估计粒子状态和每种状态的三维密度。CryoGaussian 从随机初始化的原语出发，用粒子与高斯层级的潜变量调节中心、幅值和尺度，在可微加和投影中拟合粒子图像。这里高斯用于表达密度，优化利用的是冷冻电镜观测模型。

作者按 CryoBench 的逐构象和逐图像协议比较傅里叶壳相关，并在真实粒子数据上检验重建。方法输入包含固定的粒子朝向、平移和对比传递函数。真实实验数据没有真值，论文采用已有参考重建等证据评价，证据条件不同于模拟基准。


\section*{本章小结}
\addcontentsline{toc}{section}{本章小结}

\begin{futurebox}{本章总结与展望}{ch19_summary}
本章从医学影像角度综述了~3D Gaussian Splatting~的应用方法，涵盖以下核心内容：
\begin{itemize}
\item \textbf{内窥镜与手术场景重建}（\ref{sec:19_1}节）：4DGS~是动态内窥镜重建的主流范式，针对光照变化、稀疏视角和~SLAM~等核心挑战，发展出光照嵌入、几何正则化和变形感知~SLAM~等系列方法；手术场景重建需要处理不可逆变形和组织剪切，器械建模需要精确的几何先验和可控关节参数。
\item \textbf{CT/MRI~断层成像重建}（\ref{sec:19_2}节）：3DGS~从~$\alpha$-混合~到~X~射线积分的物理适配是核心技术突破，辐射高斯~Splatting~和光线追踪两条路线各有优势；MRI~重建的核心创新在于复值~k~空间数据的显式高斯表示和运动分辨重建。
\item \textbf{显微成像与生物结构重建}（\ref{sec:19_3}节）：3DGS~在~cryo-EM~中展现出显著的效率优势（训练快~48\%、内存低~12\%），层次高斯混合模型有效处理了生物分子的构象异质性，电子断层扫描和光学显微镜也获得了专用适配。
\item \textbf{临床应用与专用场景}（\ref{sec:19_4}节）：3DGS~的实时渲染和紧凑表示（多~GB~压缩到~70 MB~）使其成为消费级设备上交互式解剖可视化的理想工具；3~张稀疏照片的牙科重建是极稀疏视角的极限挑战；超声成像需要重新设计投影模型实现专用重建。
\end{itemize}

当前医学影像~3DGS~面临的主要挑战包括：
\begin{itemize}
\item \textbf{临床部署的可靠性}：手术导航等安全关键应用对~3D~重建的鲁棒性要求极高，现有方法在光照变化、组织遮挡和运动模糊等临床干扰下的性能仍需系统评估。
\item \textbf{跨模态统一表示}：内窥镜~RGB~、CT~投影、MRI k~空间和超声回波等模态的成像物理差异巨大，缺乏统一的~3DGS~表示框架来桥接不同模态。
\item \textbf{实时~4D~重建}：手术场景的实时~4D~重建（变形组织+动态器械）仍面临计算瓶颈，当前最先进方法在实时性和精度之间需要仔细权衡。
\item \textbf{数据隐私与合规}：医学数据的隐私保护要求对~3DGS~模型的存储、传输和共享提出额外约束，需要与安全水印技术（第~21~章）相结合。
\end{itemize}
\end{futurebox}
